\documentclass[acmsmall, screen]{acmart}
\let\Bbbk\relax

\usepackage{quiver}
\usepackage{mathpartir}
% \usepackage{tikz-cd}
\usepackage{enumitem}
\usepackage{wrapfig}
\usepackage{fancyvrb}
\usepackage{comment}
\usepackage{array}

% \theoremstyle{remark} 
% \newtheorem{theorem}{Theorem}
% \newtheorem{corollary}{Corollary}
\theoremstyle{remark}
\newtheorem{remark}{Remark}
\newenvironment{DIFnomarkup}{}{}
% \pagestyle{empty}
% \settopmatter{printfolios=false}

%\newcommand{\bibliographypath}{../../paper-new/references}
%\newcommand{\defspath}{../../paper-new/defs}
\newcommand{\bibliographypath}{references}
\newcommand{\defspath}{defs}

\input{\defspath}


\begin{document}

\title{Denotational Semantics for Effectful Gradually Typed Languages using Guarded Type Theory}
\author{Eric Giovannini}
\affiliation{
  \department{Electrical Engineering and Computer Science}
  \institution{University of Michigan}
  \country{USA}
}
\email{ericgio@umich.edu}
\orcid{0009-0003-6871-1714}


\author{Max S. New}
\affiliation{
  \department{Electrical Engineering and Computer Science}
  \institution{University of Michigan}
  \country{USA}
}
\email{maxsnew@umich.edu}
\orcid{0000-0001-8141-195X}

\begin{abstract}
    Gradually-typed langugaes allow for soundly mixing static and dynamically
    typed programming styles and support migration from a dynamic to static
    typing style. Establishing the desired metatheoretic properties for
    real-world gradually-typed languages has proved challenging.
    %
    In recent work, Giovannini et al. gave a denotational semantics for a simple
    gradually-typed lambda calculus using the tools of guarded type theory, and
    furthermore showed that their semantics was sufficient for establishing the
    crucial metatheoretic property known as \emph{graduality}, which states that
    migrating a program by making the types more precise does not fundamentally
    chance the meaning of the program except to possibly catch more type errors.
    %
    
    In this work, we address two limitations of the work of Giovannini et al.
    First, we generalize their model to allow for languages with algebraic
    effects, such as I/O, global state, nondeterminism, etc. Second, we relate
    the denotational semantics to a more familiar operational semantics.
\end{abstract}

\maketitle

\newif\ifdraft
\drafttrue
\renewcommand{\max}[1]{\ifdraft{\color{blue}[{\bf Max}: #1]}\fi}
\newcommand{\eric}[1]{\ifdraft{\color{orange}[{\bf Eric}: #1]}\fi}
\newcommand{\tingting}[1]{\ifdraft{\color{red}[{\bf Tingting}: #1]}\fi}

% Question: should we present algebraic effects as having a request and response
% type, or should we fold the request type into the type of operations?

\section{Introduction}

\emph{Gradually typed languages} \cite{siek-taha06, tobin-hochstadt06} allow for
safe interoperability between static and dynamic typing disciplines in the same
codebase, and they support the smooth, \emph{gradual} migration from dynamic
typing to static typing via the addition of explicit typing information.
%
% As a result, the programmer can begin with fully dynamically-typed code and
% \emph{gradually} add explicit typing information to migrate portions of the
% codebase to a statically typed style without needing to rewrite the entire
% project in a completely different language.
%
One of the fundamental theorems for gradually typed languages is
\emph{graduality}, also known as the \emph{dynamic gradual guarantee},
originally defined by Siek, Vitousek, Cimini, and Boyland
\cite{siek_et_al:LIPIcs:2015:5031, new-ahmed2018}.
%
Informally, graduality says that migrating code from dynamic to static typing
should only allow for the introduction of static or dynamic type errors, and not
otherwise change the behavior of the program.
%

Establishing graduality for a given cast calculus has traditionally been a
tedious endeavor. One common approach involves defining a step-indexed logical
relation over an operational semantics \cite{TODO}; these methods are inherently
dependent on the particular calculus being modeled and offer little opportunity
for reuse of mathematical abstractions across different developments. Another
method, developed by New and Licata \cite{new-licata-18}, constructs a
denotational semantics using \emph{classical} domain theory. While
compositional, this approach suffers from a lack of expressiveness: classical
domain-theoretic techniques appear incapable of modeling viciously
self-recursive structures such as higher-order store
\cite{Birkedal-Stovring-Thamsborg-2009}, which are common in real-world
gradually-typed languages.
%
Thus, it remained an open question whether one could give a semantics for
gradual typing that combined the expressiveness of the step-indexed approach
with the compositional nature of denotational semantics.

% Graduality has remained a difficult property to establish, with existing
% methods being largely tied to syntactic or operational specifics of the
% language being modeled and offering little opportunity for reuse of common
% mathematical abstractions.

% Recently, Giovannini et al. \cite{gtt-sgdt} combined the expressiveness of the
% step-indexed approach with the compositional nature of denotational semantics
% by...

Recently, Giovannini et al. \cite{gtt-sgdt} answered this question affirmatively
by constructing a denotational semantics for a simple gradually-typed
programming language using the tools of guarded type theory
\cite{birkedal-mogelberg-schwinghammer-stovring2011}, and proved that their
model validates the graduality property. Their model combines the expressive
power of step-indexing techniques with the compositional nature of denotational
semantics.
%
However, the semantics developed by Giovannini et al. has the limitation that it
is a single model for a quite bare-bones gradually-typed language. Real-world
languages can include effects, polymorphism, dependent types, etc.
%
Furthermore, it is unclear how to relate their model to a more typical
operational semantics for a cast calculus. It is desirable to establish an
adequacy result relating their denotational semantics to an operational model.
More specifically, writing $\sem{M}$ for the denotation of a term $M$, adequacy
states that an equality of denotations $\sem{M} = \sem{N}$ implies that $M$ and
$N$ behave the same operationally. This allows one to use facts proved in the
denotational semantics to reason about the operational behavior of terms.
%
More specifically, an adequacy result provides a means of translating the
graduality theorem --- a statement about the \emph{denotations} of terms ---
into a statement about the \emph{operational behavior} of terms.

% In this work, we seek to address these limitations by
% In this work, we seek to address the limitations of the prior model by 
In this work, we take a step towards modeling more realistic gradually-typed
languages using guarded type theory: we extend the denotational semantics of
Giovannini et al. to account for arbitrary algebraic effects, and we establish
an adequacy theorem relating the resulting denotational model to an effectful
operational semantics.

% The level of generality needed to model arbitrary algebraic effects means that
% the semantics we produce is correspondingly generic. This is true both
% denotationally and operationally. On the denotational side, we generalize the
% semantics of Giovannini et al. to allow for arbitrary algebaic effects. The
% (guarded) graduality theorem we obtain is essentially a statement relating term
% precision to the behavior of the interaction trees that the terms denote.

% TODO: update
% Likewise, we construct a generic effectful operational semantics for the cast
% calculus in guarded type theory.

% While this approach is applicable to any
% algebraic theory representing the programming language's effects, what is
% missing is a way to relate the small-step operational semantics it produces to a
% more traditional effectful small-step semantics. For example, a small-step
% semantics for a language with global state would normally involve a transition
% relation between pairs of a term and state.

% TODO: update
% The missing step in both the operational and denotational semantics is supplied
% by a suitable notion of algebra that interprets the guarded interaction trees as
% a concrete element of an ``observation'' type.




\subsection{Overview of Approach}\label{sec:overview-of-approach}

The main contribution of this paper is a framework for giving an adequate
operational semantics to an effectful gradually-typed cast calculus. The input
to the framework consists of a specification of the algebraic effects present in
the cast calculus, given in the form of an \emph{algebraic theory}, which
contains the operations along with equations that they satisfy. More details on
algebraic theories are given in Section \ref{sec:modeling-effects}. The cast
calculus contains a term corresponding to each of the operations present in the
provided algebraic theory $\calA$, allowing the programmer to perform the
effects contained in the theory.

Every algebraic theory $\calA$ induces a monad, called the \emph{free-model
monad} for $\calA$. In addition to an algebraic theory $\calA$ specifying the
effects, our framework also takes as input a presentation of the free-model
monad for the theory $\calA$, which we denote $T_\calA$.

Given this data as input, we define an operational semantics for the cast
calculus with effects specified by $\calA$. The operational semantics has the
form
%
\[ \run : (\term{A} \to T'(\val{A})). \]
%
where $T'$ is the free-model monad $T_\calA$ augmented with an additional delay
effect we call ``administrative delays''. The additional delay effect accounts
for the fact that a priori, we do not know whether the given term will
terminate.

The graduality property says that if a term $M$ is syntactically more precise
than a term $N$, then the behavior of $M$ when run is a refinement of the
behavior of $N$. More concretely, graduality says that for any closed terms $M,
N$ of type $\nat$, if $M \ltdyn N$, then $\run{M} \lesssim \run{N}$. Here,
$\lesssim$ is an error ordering on monadic values capturing when two
computations behave the same up-to the left side terminating with an error.

To prove graduality, we proceed in several steps:
% We leverage the existing denotational model of gradual typing of Giovannini et
% al. \cite{gtt-sgdt}.
%

%\begin{enumerate}
% First, we develop a guarded denotational model of gradual typing with
% effects, taking as inspiration the existing model of Giovannini et al.
First, we develop a guarded denotational model of gradual typing with effects
based on the existing model of Giovannini et al. The result is a \textbf{guarded
denotational graduality} theorem:
%
\begin{theorem}[Guarded Denotational Graduality]
  If $\cdot \vdash M \ltdyn N: \nat$, then $\sem{M} \ltbisim \sem{N}$.
\end{theorem}
%
Here, $x \ltbisim y$ means that computation $x$ is the same as $y$ up to $x$
being an error, and $x$ and $y$ may take a different number of steps.

Next, we relate the operational semantics to the denotational semantics
via an \textbf{adequacy theorem}:
%
\begin{theorem}[Adequacy]
  Let $\cdot \vdash M : \nat$. Then we have $\run{M} \equiv_b \sem{M}$.
\end{theorem}
%
Here, $x \equiv_b y$ relates elements belonging to two different monads, where
the monad in which $\run{M}$ lives supports a special ``bureaucratic''
stepping operation that the monad where $\sem{M}$ lives does not. The reason
for this discrepancy will be explained in the section on operational
semantics, but for now it is sufficient to note that the relation $x \equiv_b
y$ means that $x$ and $y$ are the same up-to a \emph{finite} number of
bureaucratic steps in $x$. The adequacy theorem then says that running a term
$M$ does not result in an infinite sequence of bureaucratic steps.

We prove the adequacy theorem using a logical relation constructed in the
ambient guarded type theory. More details are given in Section
\ref{sec:relating-the-semantics}.

Putting the above two results together, we obtain \textbf{guarded
operational graduality}:
%
\begin{corollary}[Guarded Operational Adequacy]
  If $\cdot \vdash M \ltdyn N: \nat$, then $\run{M} \lesssim \run{N}$.
\end{corollary}
%
In the above, $x \lesssim y$ (where $x$ and $y$ are monadic values) means that,
disregarding the number of steps taken by each side, $x$ behaves the same as $y$
except that $x$ may be an error where $y$ is not.

Lastly, we translate the guarded operational graduality result to the ordinary
set-theoretic world via a process known as \emph{globalization} \cite{TODO}. The
motivation is that the statements of the results up to this point involve
features of guarded type theory, e.g., the later modality. The technique of
globalization produces a semantics that does not rely on any features of guarded
type theoty, and the result we obtain is a statement that can be made in
ordinary set theory.
%
Unlike the prior steps, the details in this step are sensitive to the
presentation of the free-model monad $T_\calA$, and there is not a general
result that can be stated at the level of an arbitrary algebraic theory
$\calA$.

However, there are two special cases where we can describe the globalization.
First, if we know that $T_\calA$ can be presented as application of a
suitably-generalized concept of monad transformers applied to the guarded lift
monad $T(\li(X))$, then we can directly generalize the globalization technique
of Giovannini et al.
%
Second, if the algebraic theory $\calA$ has \emph{no equations}, then there is a
presentation of the global version of the monad in terms of \emph{finite
computation traces}. We describe both of these special cases in Section
\ref{sec:globalization}. 


The organization of the remainder of the paper is as follows:
\begin{enumerate}
    \item In Section \ref{sec:modeling-effects}, we provide background on
    algebraic theories which will be used in the remainder of the paper.
    \item In Section \ref{sec:syntax}, we specify the syntax for gradual typing
    used in the remainder of the paper. Given a collection of algebraic effects
    specified as an algebraic theory $\calA$, we define a cast calculus
    supporting those effects. The cast calculus includes an inequational theory
    in the form of axioms for type and term precision.
    \item In Section \ref{sec:operational-semantics}, we define an operational
    semantics for the cast calculus defined in the previous section.
    \item In Section \ref{sec:denotational-semantics}, we show how to construct
    a denotational semantics for the effectful cast calculus by adapting the
    construction of Giovannini et al. \cite{gtt-sgdt} to account for arbitrary
    algebraic effects.
    \item In Section \ref{sec:relating-the-semantics}, we prove an adequacy
    theorem relating the operational and denotational models in the expected
    manner.
    \item In Section \ref{sec:globalization}, we discuss how to globalize the
    results of the previous two sections.
    \item Lastly, in Section \ref{sec:discussion}, we summarize our results,
    compare to prior work, and conclude with directions for future work.
\end{enumerate}


\subsection{Metatheoretic Preliminaries}

We begin with a brief discussion of the ambient guarded type theory in which we
work. \emph{Guarded type theory} refers to a class of type theories with
primitives that allow one to solve non-well-founded recursive equations and
reason about such solutions. It can be viewed a \emph{synthetic} approach to
step-indexed techniques \cite{birkedal-mogelberg-schwinghammer-stovring2011};
the benefit of this approach is that we avoid the tedious natural number
manipulations present in explicit step-indexing.

There are many variants of guarded type theory, but they each allow a form of
recursion to be encoded at the level of types through a modality $\later \colon
\type \to \type$ (pronounced ``later''). The use of a modality to express
guarded recursion originates in the work of Nakano \cite{Nakano2000}. More
specifically, a typical guarded type theory has at least the following features:

\begin{itemize}
  \item Given a type $A$, the type $\later A$ represents an element of type $A$
  that is available one ``abstract time step'' later. 

  \item There is an operator $\nxt : A \to\, \later A$ that ``delays'' an element
  available now to make it available later.

  \item There is a principle of \emph{guarded recursion}, implemented by a
  guarded fixpoint operator of the form
  %
  \[ \fix : \forall A, (\later A \to A) \to A. \] 
  %
  That is, to construct a term of type $A$, it suffices to assume that we have
  access to such a term ``later'' and use that to help us build a term ``now''.
  This operator satisfies the axiom that $\fix f = f (\nxt (\fix f))$.

  \item There is a method for solving guarded-recursive type equations $T \cong
  F(\later T)$. In guarded type theory with a universe of types $\type$, we can
  solve such equations by applying the guarded fixpoint operator with $A =
  \type$ \cite{birkedal-møgelberg-2013}.
  %
  In particular, $\fix$ applies when type $T$ is instantiated to a proposition
  $P : \texttt{Prop}$; in that case, it corresponds to $\lob$-induction.
  
\end{itemize}
 
We will use a tilde to denote a term of type $\later A$, e.g., $\tilde{M}$.

% TODO later is an applicative functor, but not a monad

% Clocked Cubical Type Theory
\subsubsection{Ticked Cubical Type Theory}

Following Giovannini et al. \cite{gtt-sgdt}, in this paper we will work in a
specific form of guarded type theory with access to additional constructs called
\emph{ticks} and \emph{clocks}.

% TODO review and consider trimming
Ticked Cubical Type Theory \cite{mogelberg-veltri2019} is an extension of
Cubical Type Theory \cite{CohenCoquandHuberMortberg2017} that has an additional
sort called \emph{ticks}. Ticks were originally introduced by Bahr et al.
\cite{bahr-grathwohl-bugge-mogelberg2017}. A tick $t : \tick$ intuitively acts
as evidence that one unit of time has passed. In Ticked Cubical Type Theory, the
type $\later A$ is the type of dependent functions from ticks to $A$. The type
$A$ is allowed to depend on $t$, in which case we write $\later_t A$ to
emphasize the dependence.
%
For the sake of brevity, we will also write tick application as $M_t$. The
constructions we carry out in this paper will take place in a guarded type
theory with ticks.



\section{Background on Algebraic Theories}\label{sec:modeling-effects}

In this section, we discuss our approach to modeling algebraic effects in the
context of gradual typing. These concepts will be employed in both our
denotational and operational semantics.


\subsection{Guarded Algebraic Theories}\label{sec:guarded-algebraic-theories}

We introduce the notion of a \emph{guarded algebraic theory}, a suitable
generalization of the notion of algebriac theory where we allow for the
algebraic operations to have their arguments guarded by the later modality. We
begin with the notion of a guarded algebraic signature.
% TODO cite paper on algebraic effects

A \textbf{guarded algebraic signature} $\Sigma$ consists of a collection of
operation symbols $\Op$. Each operation symbol $\oper : \Op$ has an
\textbf{ordinary arity} given by a type $\tau_\oper$ and a \textbf{guarded
arity} given by a type $\gamma_\oper$. We write $\oper : (\tau, \gamma)$ to mean
$\oper$ has ordinary arity $\tau$ and guarded arity $\gamma$.
% Each operation \oper : (\tau, \gamma) has ordinary arity \tau and guarded arity \gamma.

% We do not require the arity types to be finite.

\begin{remark}
  In defining algebraic effects, one typically defines an algebraic signature as
  a collection of operations that each have a \emph{request} type $\sigma_\oper$
  and response type $\tau_\oper$. The request type of operation $op$ is the type
  of the value that the program sends to the environment. Similarly, the
  \emph{response type} of the operation is the type of the value that the
  environment sends back to the program once the operation is complete. We can
  encode this definition into ours by having a separate algebraic operation
  $\oper_x$ for each $x : \sigma$.
\end{remark} 

Given a guarded signature $\Sigma$ as described above, we define the
\textbf{polynomial endofunctor} corresponding to $\Sigma$: 
%
\[ \Sigma(X) = \Sigma_{\oper: \Op}, (\tau_\oper \to X) \times (\gamma_\oper \to \later X). \]
%

A \textbf{$\Sigma$-algebra} is simply an algebra of the endofunctor $\Sigma$
defined above, i.e, a type $B$ equipped with a function $\Sigma(B) \to B$. More
concretely, a $\Sigma$-algebra consists of a set $B$ equipped with
interpretations for each of the operations in $\Sigma$. That is, for each $\oper
: (\tau, \gamma)$, there is a function
%
\[ \sem{\oper} : ((\tau \to B) \times (\gamma \to (\later B))) \to B. \]
%
% This view makes it clear that $\sigma$ is the request type, i.e., the type of
% value sent to the environment, while $\tau$ is the response type, i.e., the type
% of value returned from the environment.

% We will often write B to denote either a $\Sigma$-algebra or its underlying
% set.

There is a forgetful functor $U_\Sigma$ from the category of $\Sigma-$algebras
and morphisms to the category of sets. The left adjoint to this functor takes a
set $A$ to the \textbf{free $\Sigma$-algebra} $F_\Sigma(A)$ on $A$. Explicitly,
the construction is by induction and guarded recursion. We seek to solve the
equation
%
\[ F_\Sigma(A) \cong A + \Sigma(F_\Sigma(A)), \]
%
i.e.,
%
\[ F_\Sigma(A) \cong A + \Sigma_{\oper: \Op}, 
    (\tau_\oper \to F_\Sigma(A)) \times G_\oper(\later F_\Sigma(A)). \]
%
First, consider the parameterized inductive type 
%
\[ P(X) := \mu Y. (A + \Sigma_{\oper : \Op}, (\tau_\oper \to Y) \times G_\oper(X)). \]
%
Now, using guarded recursion, define $F_\Sigma(A)$ to be the unique solution to
the guarded-recursive equation
%
\[ F_\Sigma(A) \cong P(\later F_\Sigma(A)). \]
%
Then by the definitions of least fixpoint and guarded fixpoint, we have that
$F_\Sigma(A)$ is a solution to the original equation.

Elements of $F_\Sigma(A)$ can intuitively be thought of as $\Sigma$-terms
generated by the elements of $A$ and the operations in $\Sigma$. We can view
such terms as a generalization of inductively-defined trees: a term can either
be a leaf, consisting of a variable in $A$, or a node labeled by an operation
$\oper : (\tau, \gamma)$, with $\tau$-many children available $\emph{now}$ and
$\gamma$-many children available $\emph{later}$.

% A similar notion of guarded-recursive trees has been studied under the name
% \emph{guarded interaction trees} \cite{guarded-interaction-trees}.

The type $F_\Sigma(A)$ has the universal property that given any
$\Sigma$-algebra $B$ and any map of variables $f : A \to B$, there is a unique
map $F_\Sigma(A) \to B$ extending $f$. In other words, if we specify how to
interpret the variables $A$ as elements of $B$, this extends uniquely to an
interpretation of any term. We write $\sem{t}$ to denote such an interpretation.
%
% Given any equation $e = (\Gamma_e, l_e, r_e)$, we can interpret $e$ in any
% $\Sigma$-algebra $B$ by asking that each side maps to the same term in $B$. This
% will be useful in defining the notion of guarded algebraic theory.

We now consider the addition of equations to a signature.
% A \textbf{guarded algebraic theory} $\calA$ consists of a guarded signature
% $\Sigma_\calA$ and a collection of equations $E_\calA$. 
A \textbf{guarded algebraic theory} $\calA$ is given by a signature
$\Sigma_\calA$ and a collection of equations $E_\calA$. An equation
$e$ consists of a context $\Gamma_e$ of free variables, and a pair of terms
$l$ and $r$ in $F_{\Sigma_\calA}(\Gamma_e)$.

Models of a guarded algebraic theory are simply algebras of the associated
signature (i.e., a set $B$ with an interpretation of the operations in
$\Sigma_A$), in which all of the equations of the theory hold. More concretely,
a \textbf{model of a theory} $\calA = (\Sigma, E)$ is a model $B$ of the
signature $\Sigma$ such that for each equation $e = (\Gamma, l , r)$ in $E$, we have
$\sem{l} = \sem{r}$ for all instantiations $f : \Gamma_e \to B$ of the free
variables of $e$.
% i.e., for all $f : \Gamma_e \to B$.


% A \emph{model} of an algebraic theory $\calA$ is a set $B$ with an interpretation
% of the operations in $\Sigma_A$, such that all equations in $E_A$ hold. Formally
% speaking, this means that when we interpret the LHS and RHS of each equation as
% elements of $B$, those elements are equal.

For a guarded algebraic theory $\calA$, the category of $\calA$-models has a
forgetful functor $U_\calA$ to the category of sets, which simply takes the
underlying set of the model. This functor has a left adjoint $F_\calA$, which
given a set $X$ constructs the \textbf{free model} of $\calA$ on $X$. Concretely,
the free model of $\calA = (\Sigma_\calA, E_\calA)$ on a set $X$ is given by
quotienting the set of $\Sigma_\calA$-terms by the equations in $E_\calA$.
%
We can construct the free model in an analogous manner to our construction of
the free $\Sigma$-algebra for a signature $\Sigma$, using quotient-inductive
types in place of ordinary induction.
%
Under this construction, elements of $F_\calA(X)$ are mixed
inductive/guarded-recursive trees quotiented by the equations of the theory.
%
This adjunction induces a monad $T_\calA := U_\calA \circ F_\calA$ on the
category of sets, which is known as the \textbf{free model monad}.
% TODO cite
% For a given theory $\calA$, we denote the free model monad for $\calA$ by
% $T_\calA$.


% Because we are modelling gradual typing, at a minimum we want to ensure that our
% theory contains an operation corresponding to error, and since our semantics is
% intensional, i.e., step-aware, we also include an operation corresponding to
% taking an observable step. Such a step will occur whenever a non-well-founded
% recursive definition is unfolded; in the case of gradual typing this would be
% when inspecting a value of the dynamic type.
% % TODO cite GTT-SGDT paper?

\subsection{Operations on Algebraic Theories}

Given two guarded algebraic theories $\calA_1$ and $\calA_2$, we define their
\emph{coproduct} $\calA_1 + \calA_2$ by combining the operations and the
equations in $\calA_1$ and $\calA_2$. We do not impose any additional equations.
% TODO cite paper on combining algebraic theories.
% TODO example

Given two theories $\calA_1$ and $\calA_2$, we define their \emph{tensor
product} $\calA_1 \otimes \calA_2$ by combining the operations and equations as
in the coproduct, but also imposing additional equations that state the
commutativity of each of the operations in $\calA_1$ with all of those in
$\calA_2$.
% TODO example


\subsection{Examples of Algebraic Theories}

We introduce two specific guarded algebraic theories that will be used thoughout
the paper.

\begin{example}[Errors]\label{ex:err}
    The theory of computations that can error $\calA_{err}$ is defined as
    follows. Let $\Sigma_{err}$ contain a single operation $\errop$ with
    $\tau_{\errop} = 0$ and $G_{\errop}(X) = 1$. There are no equations
    associated with this theory.

    Then a model of this theory is a set $B$ equipped with a function $((0 \to
    B) \times 1) \to B$, i.e., a function $1 \to B$, which is just an element of
    $B$. In other words, models of this theory are simply pointed sets $(B,
    \mho)$ where $\mho \in B$.

    The free-model monad for $\calA_{err}$ is given by freely adjoining a point
    to its argument:
    %
    \[ F_{\calA_{err}}(X) \cong X + 1. \]
    
\end{example}

In our next example, the guarded arity is nontrivial:

\begin{example}[Computational Steps]\label{ex:steps}
    %We define the theory $\calA_{step}$ of computational steps as follows. 
    The theory of computational steps $\calA_{step}$ models computations that
    can take an observable computational step. The theory is defined as follows.
    Let $\Sigma_{step}$ contain a single operation $step$ with $\tau_{step} = 0$
    and $G_{step}(X) = X$. There are no equations associated with this theory.

    Then a model of this theory is a set $B$ equipped with a function
    $\sem{step} : ((0 \to B) \times (\later B)) \to B$, i.e., a function $\later
    B \to B$. We call this function $\theta_B$ in accordance with prior work
    using guarded type theory.

    The free-model monad for $\calA_{step}$ is the so-called \emph{guarded-lift
    monad} defined as the unique solution to the guarded-recursive equation
    %
    \[ L(X) \cong X + \later L(X). \]
    % TODO cite paper on guarded lift monad
\end{example}


\subsection{Relational Aspects}

An important aspect that we have elided so far is the fact that, since we are
modeling gradual typing, we need our mathematical structures to be equipped with
a notion of relation to model the graduality ordering. This theory was developed
extensively by Giovannini et al. in their denotational model for gradual typing
\cite{gtt-sgdt}. The implication of this is that the categories involved in the
adjunction described above are not merely the category of sets and the category
of $\calA$-models. We require additional relational structure: a relation to
model the graduality ordering as well as a notion of weak bisimilarity that
relates terms that differ only in the number of steps they take. These relations
are crucial ingredient in the graduality proof. Since these relational aspects
are not relevant for the operational semantics, we defer a detailed discussion
to the section on denotational semantics.

% for now, it is sufficient to keep in mind that the categories involved, and the
% free-model monad $T_\calA$ for a given algebraic theory $\calA$, enjoy this
% additional structure.

\section{A Generic Syntax for Gradual Typing with Effects}\label{sec:syntax}

To begin, we fix a framework for specifying gradual cast calculi with algebraic
effects. The framework is parameterized by an ordinary algebraic theory $\calA$
specifying the effects available in the syntax.

The cast calculus is call-by-value and is presented in Figure
\ref{fig:gtlc-syntax}.

% TODO should we use CBPV syntax instead?
% TODO what about coproducts and lazy products?
% TODO if we are doing fine-grain CBV, do we need return and bind?

Types include natural numbers, functions $A \ra A'$, products $A \times A'$, and
a dynamic type, which is the type assigned to a dynamically-typed value, i.e., a
value whose type is unknown at compile-time.

As is typical in cast calculi, there is an ordering on types known as \emph{type
precision} specifying when one type is more precise (i.e., less dynamic) than
another. Following prior work, our calculus uses an explicit syntax of type
precision \emph{derivations} of the form $c : A \ltdyn A'$. Such a derivation
should be thought of as a proof that $A \ltdyn A'$. This presentation is
particularly convenient for carrying out proofs by induction.

The cast calculus has special coercion terms called casts, which come in two
forms: upcasts and downcasts. Given $c : A \ltdyn A'$ and $M$ of type $A$, the
term $\upc c M$ has type $A'$, and given $N$ of type $A'$, the term $\dnc c N$
has type $A$. Upcasts are pure, i.e., they always succeed (at least to a first
order), whereas downcasts are effectful.

% Lastly, we have terms $\synop{op}{M}$ corresponding to each operation in the
% algebraic theory $\calA$. More concretely, if $\eff{op}{\sigma}{\tau}$ is an
% operation in $\calA$ and $M$ is of type $\sigma$, then $\synop{op}{M}$ is of
% type $\tau$.

Lastly, we have terms $\synop{op}{M}$ corresponding to each operation in the
algebraic theory $\calA$. More concretely, if $\oper$ is an operation in $\calA$
of arity $\tau_\oper$, then there is a term $\oper$ of type $\tau_\oper$.

% We mention two notes about the relationship between the given algebraic theory
% and the syntax. First, 

We assume that the meta-level arity types of each operation exist as
object-level types within the cast calculus. While these types are not
\emph{equal}, there will be an isomorphism allowing us to translate between
them. For example, if we have an operation $\oper$ of arity $\mathbb{N}$, then
in the syntax the corresponding term will have type $\nat$. We leave implicit
the conversion between the syntactic and semantic versions of the arities.


% Terms:
% upcasts and downcasts
% effect operations

% explain type precision

% New, Licata, and Ahmed gave a syntax and logic axiomatizing the graduality
% property \cite{new-licata-ahmed2019}.

\begin{figure}
  \begin{mathpar}
    \begin{array}{rcl}
    \text{Types } A &::=
      &\nat 
      \alt \,\dyn 
      \alt A \ra A' 
      \alt A \times A'\\
    \text{Type Precision } c &::=
      & r(A) 
      \alt c c'     \alt \iarr 
      \alt \inat    \alt \itimes 
      \alt c \ra c' \alt c \times c'\\
    \text{Values } V &::=
      & x 
      \alt \upc c V \alt \zro 
      \alt \suc\, V \alt \lda{x}{M} 
      \alt (V,V') \\ 
    \text{Terms } M,N &::=
      & \err 
      \alt \upc c M  \alt \dnc c M 
      \alt \zro      \alt \suc\, M 
      \alt \lda{x}{M} \\ 
      &&\alt M\, N   \alt (M,N) 
      \alt \textrm{split } (x,y) = M \textrm{ in } N 
      \alt \textrm{let } x = M \textrm{ in } N 
      \alt \synop{op}{V}{x}{M} \\
    \text{Contexts } \Gamma &::= &\cdot \alt \Gamma, x : A \\
    \text{Ctx Precision } \Delta &::=& \cdot \alt \Delta,x:c
  \end{array}

  \inferrule*[right=Up]
  {\Gamma \vdash M : A \and c : A \ltdyn A'}
  {\Gamma \vdash \upc c M : A'}

  \inferrule*[right=Dn]
  {\Gamma \vdash N : A' \and c : A \ltdyn A'}
  {\Gamma \vdash \dnc c N : A}

  \inferrule*[right=Eff]
  {\eff{op}{\sigma}{\tau} 
    \and \Gamma \vdash V : \sigma 
    \and \Gamma,x : \tau \vdash M : A}
  {\Gamma \vdash \synop{op}{V}{x}{M} : A}

  \inferrule{}{\Gamma \vdash \mho : A}
  \end{mathpar}
  \begin{mathpar}
    \inferrule{}{r(A) : A \ltdyn A}\and
    \inferrule{c : A \ltdyn A' \and c' : A' \ltdyn A''}{cc' : A \ltdyn A''}\and
    \inferrule{}{\iarr \colon \dyn \ra \dyn \ltdyn \dyn}\and
    \inferrule{}{\inat \colon \nat \ltdyn \dyn}\and
    \inferrule{}{\itimes \colon \dyn \times \dyn \ltdyn \dyn}\and
    \inferrule{c_i : A_i \ltdyn A_i' \and c_o : A_o \ltdyn A_o'}{c_i \ra c_o : (A_i \ra A_o) \ltdyn (A_i' \ra A_o')}\and
    \inferrule{c_1 : A_1 \ltdyn A_1' \and c_2 : A_2 \ltdyn A_2'}{c_1 \times c_2 : (A_1 \times A_2) \ltdyn (A_1' \times A_2')}\and
     r(A)c \equiv c\and
     c \equiv cr(A')\and
     c(c'c'') \equiv (cc')c''\and
     r(A_i \ra A_o) \equiv r(A_i) \ra r(A_o)\and
     r(A_1\times A_2) \equiv r(A_1) \times r(A_2)\and
     (c_i \ra c_o)(c_i' \ra c_o')\equiv (c_ic_i' \ra c_oc_o') \and
     (c_1\times c_2)(c_1'\times c_2')\equiv (c_1c_1' \times c_2c_2')
  \end{mathpar}
  \caption{Syntax for Gradual Typing with Effects}
  \label{fig:gtlc-syntax}
\end{figure}



% rules for term precision
Next, we have the axioms for reasoning about terms. First is the equational
theory, including the usual CBV $\beta\eta$ rules for function and product
types.
%
We also have term precision, the extension of type precision to terms. Figure
\ref{fig:term-prec} presents selected rules for the term precision relation.
%
As in prior work, term precision rules have the form $\Delta \vdash M \ltdyn M'
: c$ where $\Delta$ is a context mapping variables to type precision
derivations.
%
% The judgment is only well formed when every use of $x : c$ for $c : A
% \ltdyn A'$ is used with type $A$ in $M$ and $A'$ in $M'$ and similarly
% the output types match $c$.

The term precision rules are as follows. First, the rule \textsc{ErrBot} says
that the error term $\mho$ is the most precise term.
%
Next, we have a congruence rule for each term constructor, e.g., the rule for
application says that if $M \ltdyn M'$ and $N \ltdyn N'$ then $M\,N \ltdyn
M'\,N'$. For the sake of brevity, the congruence rules are elided from the
figure.
%
Using these congruence rules, reflexivity $M \ltdyn M$ is admissible. As in the
syntax of Giovannini et al. \cite{gtt-sgdt}, and in accordance with the original
formulation of the dynamic gradual guarantee \cite{siek_et_al:LIPIcs:2015:5031},
we do not take transitivity of term precision as a primitive rule.
%
% We include a rule that says that equivalent type precision derivations
% $c \equiv c'$ are equivalent for the purposes of term precision.

Lastly, we have the four fundamental rules for reasoning about casts (\upl,
\upr, \dnl, \dnr), which axiomatize the upcasts and downcasts as behaving like
least upper bounds and greatest lower bounds, respectively. 

% these say that the upcast is a kind of \emph{least upper bound} and dually that
% the downcast is a \emph{greatest lower bound}.


\begin{figure}
  \begin{mathpar}

  % TODO do we need this rule?
  % \inferrule*[right=EquivTyPrec]
  % {\Delta\vdash M \ltdyn M' : c \and c \equiv c'}
  % {\Delta\vdash M \ltdyn M' : c'}

  \inferrule*[right=ErrBot]
  {}
  {\Delta \vdash \mho \ltdyn M : c}

  \inferrule*[right=UpL]
  {M \ltdyn M' : cc_r}
  {\upc {c} M \ltdyn M' : c_r}

  \inferrule*[right=UpR]
  {M \ltdyn M' : c_l}
  {M \ltdyn \upc {c} M' : c_lc}

  \inferrule*[right=DnL]
  {M \ltdyn M' : c_r}
  {\dnc {c} M \ltdyn M' : cc_r}

  \inferrule*[right=DnR]
  {M \ltdyn M' : c_lc}
  {M \ltdyn \dnc {c} M' : c_l}
  \end{mathpar}
  \caption{Term Precision Rules (Selected)}
  \label{fig:term-prec}
\end{figure}

Our goal in the remainder of this work is to define an operational semantics for
this cast calculus and prove that it satisfies graduality. For an outline of our
approach to doing so, refer to Section \ref{sec:overview-of-approach}.

% compositional denotational semantics of types, terms, type and term precision
% from which we can easily extract a big step semantics that satisfies graduality
% and respects the equational theory of the calculus.


\section{Operational Semantics}\label{sec:operational-semantics}

In this section, we define a small-step operational semantics for the generic
cast calculus introduced in Section \ref{sec:syntax}. Our approach is based
loosely on the recent work of Kavvos \cite{kavvos-2025}, who gave a generic
operational semantics for a language with algebraic effects.

Recall that our framework is parameterized by an \emph{ordinary} algebraic
theory $\calA$; recall that this means there are no guarded operations in
$\calA$.
%
Because we are modelling gradual typing, at a minimum we want to ensure that our
theory contains an operation corresponding to error, and since our semantics is
intensional, i.e., step-aware, we also include an operation corresponding to
taking an observable step. Such a step will occur whenever a non-well-founded
recursive definition is unfolded; in the case of gradual typing this would be
when inspecting a value of the dynamic type.
%
In fact, for reasons that will become clear later in the present section, we
distinguish between two different kinds of semantic stepping operations:
\textbf{true steps}, which occur when inspecting a non-well-founded definition
such as the dynamic type, and \textbf{bureaucratic steps}, which represent all
other steps in the operational semantics. A key difference between the two is
that bureaucratic steps are invisible in the denotational semantics; they are
merely a bookkeeping tool necessary in defining the operational semantics.
%
If we view steps as having an associated ``cost'', we can think of true steps as
having a nonzero cost, while bureaucratic have zero cost.
%
% TODO: is this accurate?
The logical relation defined in Section \ref{sec:relating-the-semantics} will
imply that when running any term, there can be at most finitely many consecutive
bureaucratic steps between any two operations.

From the given ordinary algebraic theory $\calA$, we define two related guarded
algebraic theories:
%
\begin{enumerate}
  \item We define the theory combining $\calA$ with errors:
  %
  $\calA_\mho := \calA \otimes \th{\mho}$,
  %
  where $\th{\mho}$ is the theory of errors.

  \item We define the theory combining $\calA$ with errors and true steps:
  %
  $\calA_{\mho s} := \calA \otimes (\th{\mho} + \th{\step})$,
  %
  where $\th{\step}$ and $\th{\mho}$ are the theories of computational steps and
  error, respectively, and $\otimes$ and $+$ are the tensor and coproduct of
  algebraic theories (see Section \ref{TODO}). Notice that we take the coproduct
  with $\th{\mho}$ rather than the tensor. This is because otherwise, by Remark
  \ref{TODO}, we would end up with the equation $\mho = \fix\, \theta$.

  \item We define the theory combining $\calA$ with errors and both true and
  bureaucratic steps:
  % 
  $\calA_{\mho sb} := 
    (\calA \otimes (\th{\mho} + \th{\step} + \th{\step}))$. 
%
  We denote the first step operation by $\theta^{s}$ and the second step
  operation by $\theta^{b}$.
\end{enumerate}

We write $T_\mho$, $T_{\mho s}$ and $T_{\mho sb}$ to denote the free-model
monads for the theories $\calA_{\mho}$, $\calA_{\mho s}$, and $\calA_{\mho +
sb}$, respectively. (Recall that the free-model monad for a theory is unique up
to isomorphism of monads.)
%
The monad $T_{\mho}$ supports each of the operations of the input theory
$\calA$, in addition to a nullary error operation $\mho$. This means that for
each operation $\oper$, of arity $\tau_\oper$, there is a function
$\sem{\oper}^{T X} : (\tau_\oper \to T X) \to T X$.
%
The monad $T_{\mho s}$ supports each of the above, in addition to a guarded step
operation $\theta^s : \later T_{\mho s}(X) \to T_{\mho s}(X)$.
%
The same is true of $T_{\mho sb}$, but it has two guarded step operations
$\theta^s, \theta^b : \later T_{\mho sb}(X) \to T_{\mho sb}(X)$.

Since $\calA_{\mho}$ is a subtheory of $\calA_{\mho + s}$, there is a morphism
of monads $\iota : T_\mho \to T_{\mho s}$ and likewise there is a morphism of
monads $\iota' : T_{\mho s} \to T_{\mho sb}$.

% TODO: are specific presentations of the monads taken as input?

\subsection{Defining the Operational Semantics}

% TODO: define the theory of \calA + stepping

In what follows, we use $\term{A}$ to denote the (meta-theoretic) type of closed
terms of type $A$ in the cast calculus, and likewise $\val{A}$ denotes the type
of closed values of type $A$ in the cast calculus.

Our goal is to define a guarded big-step operational semantics, i.e., a function
%
\[ \run : \term{A} \to T_{\mho sb} (\val{A}). \]
%
Observe that the output of running a term is a monadic value. The monad accounts
for all possible effects that may occur as the term is being run, including both
true and bureaucratic steps.
%
We write $M \Dwn \hat{V}$ to mean that $\run(M) = \hat{V}$ and we write $M
\Uparrow$ to mean that $\run(M) = \fix \thetatrue$.


To define the $\run$ function, we first model small-step reduction as a function
%
\[ \step : (M : \term{A}) \to (\Sigma_{V \in \val{A}} (M = V)) 
    + (\term{A} + \Sigma_\mho(\term{A}) \times \steptype). \]
%
Here, $\Sigma_\mho$ is the signature of the algebraic theory $\calA_\mho$, and
$\steptype$ is simply a type with two inhabitants: $\steptype := \tstep \alt
\bstep$, where $\tstep$ represents a true step and $\bstep$ represents a
bureaucratic step.
%
Then the $\step$ function states that for every closed term $M$ of type $A$,
either $M$ is a value or it takes a step. When a term steps, it can be a
\textbf{pure reduction step} to some other term $M'$, or an \textbf{effectful
reduction step}, involving an operation in the input theory $\calA$ or error.
For both kinds of reduction, we also track the nature of the step (true or
bureaucratic).

% Notation for stepping
% For the sake of brevity, we use the following notation for stepping. If
% $\step{M} = \inr(\hat{N}, ty)$, we say that $M \stepsin{0} \hat{N}$ if $ty =
% \bstep$ and $M \stepsin{1} \hat{N}$ if $ty = \tstep$. That is, the natural
% number (either $0$ or $1$) indicates whether the step is a true step or an
% bureaucratic step. In the cost analogy discussed above, the natural number
% represents the ``cost'' of performing the step, whereby true steps have cost $1$
% and bureaucratic steps have cost $0$.
% %
% For a pure reduction step $\hat{N} = \inl N$, we omit the $\inl$ and simply
% write $M \stepsin{k} N$.

We define an endofunctor $\Sigmahat(X) = X + \Sigma_\mho(X)$. Its action on a
function $f : X \to Y$ is defined in the obvious way, using $f$ in the first
case, and the action of $\Sigma_\mho$ on $f$ in the second case.
%
We write $\cpure{M}$ for the left injection into $\Sigmahat(X)$, and
$\ceff{\bar{M}}$ for the right injection.

For the sake of brevity, we use the following notation for stepping. We write $M
\stepsin{k} \hat{N}$, where $\hat{N}$ is a metatheoretic value of type
$\Sigmahat(\term{A})$, to mean that $\step(M) = \inr (\hat{N}, t)$ where $t =
\tstep$ if $k = 1$ and $t = \bstep$ if $k = 0$. That is, the number $k$ (either
$0$ or $1$) indicates whether the step is a true step or a bureaucratic step. In
the cost analogy discussed above, $k$ represents the ``cost'' of performing the
step, where true steps have cost $1$ and bureaucratic steps have cost $0$.
%
For a pure reduction step $\hat{N} = \cpure N$, we sometimes omit the
$\cpure{}$ and simply write $M \stepsin{k} N$.



% Old version:
% To define the $\run$ function, we first model small-step reduction as a function
% %
% \[ \step : (M : \term{A}) \to (\Sigma_{V \in \val{A}} (M = V)) + (T(\term{A}) \times \steptype). \]
% %
% That is, given a term of type $A$, either that term is a value, or it takes an
% effectful step. Here, $\steptype$ is simply a type with two inhabitants:
% $\steptype := \tstep \alt \bstep$, where $\tstep$ represents a true step and
% $\bstep$ represents a bureaucratic step.
% %
% Notice that the small-step function, by virtue of its type signature, includes
% the usual progress, preservation, and determinacy theorems.



% Goal: operational graduality
% TODO: should we move this to before the operational semantics section?
% Our main goal in the rest of the paper will be to prove that this big-step
% semantics validates the graduality property:

% \begin{theorem}[Graduality]
%   Suppose $M \ltdyn N : \nat$. Then either:
%   \begin{enumerate}
%     \item $M \Dwn \mho$, where $\mho \in T_{db} (\val{\nat})$ is the error element of the monad.
%     \item $M \Dwn \eta\, n$ and $N \Dwn \eta\, n$ where $n \in \nat$.
%     \item $M \Uparrow$ and $N \Uparrow$.
%     \item $M \Dwn \text{op}(V)$ and $N \Dwn \text{op}(V)$ % M and N both terminate with a 0-ary effect
%   \end{enumerate}
% \end{theorem}

\subsection{Defining the Small-Step Semantics}

Figure \ref{fig:operational-semantics} defines the small-step semantics. The
presentation is as a series of rules, but these can be converted to the form of
the $\step$ function whose signature we have discussed above:
%
\[ \step : (M : \term{A}) \to (\Sigma_{V \in \val{A}} (M = V)) 
    + (\term{A} + \Sigma_\mho(\term{A}) \times \steptype). \]
%   
In particular, our notational conventions are as follows:
%
\begin{itemize}
  \item A reduction rule $M \stepsin{k} \hat{N}$, where $\hat{N}$ is of type
  $(\term{A} + \Sigma_\mho(\term{A})) \times \steptype)$, should be translated as
  $\step(M) = \inr (\hat{N}, t)$ where $t = \tstep$ if $k = 1$ and $t = \bstep$
  if $k = 0$. That is, the number $k$ (either $0$ or $1$) determines whether the
  step is considered a true step or a bureaucratic step.

  \item A \emph{pure reduction rule} is a special case of the above where
  $\hat{N} = \cpure{N}$ for some term $N$. In this case, we write $M
  \stepsin{k} N$, where $N$ is a term, leaving the $\cpure{}$ constructor
  implicit.

  \item An \emph{impure reduction rule} is one where $\hat{N} = \ceff{(\oper,
  u)}$ for some $\oper : \calA_\mho.\Op$ and $u : \tau_\oper \to \term{A}$. Note
  that every operation in $\calA_\mho$ has trivial guarded arity, so there are
  no guarded values.
  
  \item Any term $M$ that does not match one of these rules is equal to a value
  $V$, and hence $\step(M)$ returns $\inl$ applied to the value $V$ and the
  proof that $M$ is equal to $V$.

  \item We write $\mho$ to denote the element $(\err, *)$ of
  $\Sigma_\mho(\term{A})$.
\end{itemize}

% TODO cite wrapping semantics for casts
% TODO define/explain tags
% TODO make sure to distinguish between meta-level lambda and object-level lambda
% TODO fix the congruence rules for evaluation contexts

\begin{figure}

  \textbf{Beta rules:}

  \begin{mathpar}

    % Usual beta rules:
    \inferrule*[right=Beta]
      { }
      {(\lambda x : A.M)(V) \stepsin{0} M[V/x]}

    \inferrule*[right=Split]
      { }
      {\synsplit{(V_1, V_2)}{x_1}{x_2}{M} \stepsin{0} M[V_1, V_2 / x_1, x_2]}

    % \inferrule*[right=Bind]
    %   { }
    %   {\synbind{\ret\, V}{x}{M} \stepsin{0} M[V/x]}
  \end{mathpar}

  \textbf{Casts (functorial actions and composition):}
  \begin{mathpar}

    % Upcast of a function steps when applied to a value:
    \inferrule*[right=Up-Arr]
      { c_i : A_i \ltdyn A_i' \and c_o : A_o \ltdyn A_o' 
        \and \Gamma \vdash V_f : A_i \ra A_o 
        \and \Gamma \vdash V : A_i'}
      { (\upc (c_i \ra c_o) V_f) \stepsin{0} (\lambda x. (\upc {c_o} (V_f\, (\dnc {c_i} x)))) }

    % Downcast of a function steps when applied to a value:
    \inferrule*[right=Dn-Arr]
      { c_i : A_i \ltdyn A_i' \and c_o : A_o \ltdyn A_o' 
        \and \Gamma \vdash V_f : A_i \ra A_o 
        \and \Gamma \vdash V : A_i'}
      { (\dnc (c_i \ra c_o) V_f) \stepsin{0} (\lambda x. (\dnc {c_o} (V_f\, (\upc {c_i} x)))) }


    % Up/down cast of a paired value steps to the pair of the up/down casted
    % values:
    \inferrule*[right=Up-Pair]
      { c_1 : A_1 \ltdyn A_1' \and c_2 : A_2 \ltdyn A_2'}
      { \upc (c_1 \times c_2) (V_1, V_2) \stepsin{0} (\upc {c_1} V_1, \upc {c_2} V_2) }

    \inferrule*[right=Dn-Pair]
      { c_1 : A_1 \ltdyn A_1' \and c_2 : A_2 \ltdyn A_2'}
      { \dnc (c_1 \times c_2) (V_1, V_2) \stepsin{0} (\dnc {c_1} V_1, \dnc {c_2} V_2) }
  

    % "Composition" of casts with inj-arr and inj-times:
    \inferrule*[right=$\iarr$CompUp]
      { c : A \ltdyn (\dyntodyn) \and \Gamma \vdash M : A }
      { (\upc (\injarr c) M) \stepsin{0} (\upc \iarr (\upc c M))}

    \inferrule*[right=$\iarr$CompDn]
      { c : A \ltdyn (\dyntodyn) \and \Gamma \vdash M : \dyn }
      { (\dnc (\injarr c) M) \stepsin{0} (\dnc c (\dnc \iarr M))}

    \inferrule*[right=$\itimes$CompUp]
      { c : A \ltdyn (\dyn \times \dyn) \and \Gamma \vdash M : A }
      { (\upc (\injtimes c) M) \stepsin{0} (\upc \itimes (\upc c M))}

    \inferrule*[right=$\itimes$CompDn]
      { c : A \ltdyn (\dyn \times \dyn) \and \Gamma \vdash M : \dyn }
      { (\dnc (\injtimes c) M) \stepsin{0} (\dnc c (\dnc \itimes M))}

  \end{mathpar}

  \textbf{Casts involving dyn:}
  \begin{mathpar}
    % Downcasts of upcasted-to-dyn values:
    \inferrule*
      { }
      { \dnc \inat (\upc \inat V) \stepsin{0} V}

    \inferrule*
      { c : A \ltdyn (\dyn \times \dyn)}
      { \dnc {\injtimes c} (\upc {\injtimes c} V) \stepsin{0} V}

    \inferrule*
      { c : A \ltdyn (\dyntodyn)}
      { \dnc {\injarr c} (\upc {\injarr c} V) \stepsin{1} V}

    \inferrule*
     { \text{tag}_1, \text{tag}_2 \in \{ \mathbb{N}, \times, \ra \} }
     { \dnc {\text{Inj}_{\text{tag}_1}} (\upc {\text{Inj}_{\text{tag}_2}} V) \stepsin{0} \ceff{(\mho, *)} }


  \end{mathpar}

  \textbf{Effects and congruence rules:}
  \begin{mathpar}

    % Effects
    % \inferrule*[right=Eff]
    %   { \eff{op} \sigma \tau }
    %   { \synop{op} V x M \stepslabel{\eff {op} V W} M[W/x] }

    \inferrule*[right=Eff, leftskip=2em]
      { \oper : \tau }
      { \lett {x} {\oper} {N} \stepsin{0} \ceff{(\oper, (\lambda (z : \tau).M[z/x]))} }  

    % Congruence
    % Accounts for the fact that reduction may be impure
    \inferrule*[right=Cong]
      { i \in \{0,1\} \and M \stepsin{i} \hat{N}}
      { E[M] \stepsin{i} \Sigmahat(E[\cdot])(\hat{N}) }

    % \inferrule*[right=EffCong]
    %   {M \stepslabel {\eff {op} V W} N}
    %   {E[M] \stepslabel {\eff {op} V W} E[N]}


  \end{mathpar}
  \caption{Operational Semantics (Selected Rules)}
  \label{fig:operational-semantics}
\end{figure}

There are $\beta$ rules for functions, pairs, and bind, which are standard. Next
are the rules involving casts. We first have rules involving upcasts and
downcasts of functions and products.
%
For function casts, we use the standard \emph{wrapping semantics} of contracts
\cite{findler-felleisen02}. In particular, when a function $V_f$ is upcasted, it
steps to a lambda expression. When the upcasted function is applied to a value,
the value is downcasted according to the domain and applied to the original
function $V_f$, and then the result is upcasted according to the codomain. The
dual holds for a function that is downcasted.
%
When a pair of values $(V_1, V_2)$ is upcasted, we simply upcast each component
of the pair, and likewise for downcasting.
%
Next we have rules for up/down casts involving composed precision derivations,
i.e., $\injarr c$ and $\injtimes c$. These rules decompose such casts into the
cast for $c$ and the primitive injection casts $\iarr$ or $\itimes$. 
% These rules express the behavior of such casts in
% terms of the behavior of the cast for $c$ and the behavior of the primitive
% injection casts $\iarr$ or $\itimes$. 
For example, to upcast a term using the precision derivation $\injarr c$ where
$c : A \ltdyn (\dyntodyn)$, we first upcast the term by $c$, and then upcast the
result by $\iarr : \dyntodyn \ltdyn \dyn$.
%
Finally, we have the rules involving downcasting a value that was cast to $\dyn$
using the injection casts $\inat$, $\injarr c$, and $\injtimes c$. In each case,
we check that the tag of the upcast matches the tag of the downcast, e.g., the
upcast is $\iarr$ and the downcast is $\iarr$. If the tags match, then the cast
steps to the original value. Otherwise, the term takes an effectful step to
$\ceff{(\mho, *)}$, representing that an error was raised. Note that when
downcasting a dynamically-typed function, we record one recursive unfolding. In
the denotational semantics, this corresponds to inserting a $\theta$.
%
% Note also that these rules only mention the primitive injection forms $\inat$,
% $\iarr$, and $\itimes$ rather than the composed forms $\injarr c$ and $\injtimes
% c$. This is sufficient because casts involving $\injarr c$ or $\injtimes c$ will
% first step to terms involving $\iarr$ and $\itimes$, which can then be handled
% using the rules just discussed.

Next, we have our second impure reduction rule. This rule states that for an
operation $\oper$ with arity $\tau$, when we bind $x$ to the result of $\oper$
in a term $N$, the let expression takes an effectful step to the metatheoretic
value $\ceff{(\oper, (\lambda (z : \tau).M[z/x]))}$. Note that since $\tau$ is a
metatheoretic type, strictly speaking a conversion is necessary to substitute
$z$ in the object-level term $M$. As mentioned in Section \ref{sec:syntax}, we
assume that the arity of each operation is isomorphic to a corresponding
syntactic type. For the sake of simplicitly, we elide these conversions.

% for any possible request value $V$ and
% response value $W$, the term $\synop{op} V x M$ steps to $M[W/x]$. This
% represents the fact that the environment has responded with value $W$, after
% which the program continues as $M$ with $W$ in place of $x$.

The final rule says that we can perform pure and impure reductions under an
evaluation context. Recall that $\Sigmahat$ is the endofunctor defined by
$\Sigmahat(X) = X + \Sigma_\mho(X)$. The function $E[\cdot]$ is shorthand for
the function taking a term $M$ to $E[M]$.
% The final rules are simply congruence rules allowing us to perform pure and
% impure reductions under an arbitrary evaluation context.

\subsection{Defining a Runnable Semantics}

Now that we have defined a small-step semantics, we discuss how to extract a
runnable ``big-step'' semantics. More concretely, recall that our goal is a function
%
\[ \run : \term{A} \to T_{\mho sb} (\val{A}). \]
%
Notice that the codomain allows for two kinds of steps: true steps, and
bureaucratic steps. It is necessary that the monad be able to accommodate both
kinds of steps. 

% On the one hand, we want to preserve the true steps taken by terms so that we
% can relate them to the recursive unfoldings on the denotational side. On the
% other hand, because it is not necessarily the case that there are an 

We define $\run$ by guarded recursion to be the unique function satisfying the
properties
%
\begin{mathpar}
  \inferrule*
    { \step{M} = \inl (V, eq) }
    { \run{M} = T_{\mho sb}.\eta(V) }

  \inferrule*
    { \step{M} = \inr (\hat{N}, \mathit{ty}) }
    { \run{M} = \delta_{\mathit{ty}} (\run^{\dag} {(\iota(f(\hat{N})))}) }
\end{mathpar}
%
In the above:
%
\begin{itemize}
  \item $f : (\term{A} + \Sigma_\mho(\term{A})) \to T_{\mho}(X)$ is defined by
  cases, mapping $\cpure{M}$ to $\eta^{T_{\mho}}(M)$, and mapping $\ceff{(\oper, u, g)}$ 
  to $T_{\mho}(\term{A}).\oper(\eta \circ u, (\later \eta) \circ g)$.
  (Note that $u : \tau_\oper \to \term{A}$ and $g : \gamma_\oper \to \later (\term A)$.)
  \item $\iota : T_{\mho}(X) \to T_{\mho sb}(X)$ is the inclusion of the free
  model monad $T_{\mho}$ into $T_{\mho sb}$. 
  \item The variable $\mathit{ty}$ has type $\steptype$.
  \item $\delta_{\mathit{ty}} = \theta_{\mathit{ty}} \circ \nxt$ is the
  ``delay'' function on $T_{\mho sb}$. The type of delay is determined by
  $\mathit{ty}$.
\end{itemize}
%
Notice that the recursive call to $\run$ is guarded by a $\theta$, so the
guarded fixpoint operator provides a unique function satisfying the above
properties.

We next define a function that ``runs'' an evaluation context $E$. Specifically,
if $E : A \multimap B$, then we define $\run{E} : T_{\mho sb}(\term A) \to
T_{\mho sb}(\val B)$ as
%
\[ \run{E} := (\lambda M. \run{(E[M])})^\dag, \]
%
where $\dag$ is the Kleisli extension operator for the monad $T_{\mho sb}$.

Then we have the following lemma, which says that running commutes with filling
an evaluation context:
%
\begin{lemma}[Run is a Homomorphism]\label{lem:run-fill}
  $\run{(E[M])} = (\run{E})(\run{M})$.
\end{lemma}
% TODO cite appendix.
The proof uses \lob-induction; see the Appendix for details.



% The $\run$ function satisfies the following additional properties that will be
% useful in relating the operational and denotational semantics in Section
% \ref{sec:relating-the-semantics}:


\begin{comment}
  Following the work of Kavvos, our operational semantics has two kinds of
  reductions. There are \emph{pure reductions}, which comprise the usual $\beta$
  reductions in an ordinary lambda calculus as well as reductions involving casts.
  There are also \emph{impure reductions}, which result from the program making a
  request to its environment and receiving a value in response.

  The pure reductions are labelled by a natural number $i \in \{0, 1\}$, which
  intuitively tracks the number of recursive unfoldings involved in each step. The
  source of such steps is the downcast from the dynamic type to a function, since
  in the denotational semantics the function is only available later and hence the
  downcast involves an observable computational step. We track the recursive
  unfolding of terms in the operational semantics to help us relate the
  operational and denotational semantics, which we address in a subsequent
  section.
  %

  The impure reductions do not involve any recursive unfolding, so they are not
  labeled by a natural number. A general impure reduction has the form
  %
  \[ M \stepseff{op}{V}{W} N, \]
  %
  which means that $M$ sends value $V$ to the environment and continues as $N$
  upon receiving response $W$.

  % TODO if this is for fine-grain call-by-value, then do we need rules involving
  % function application?
  % TODO define evaluation contexts
  % TODO need to define a subset of "first-order values" to be used in the effect
  % operations.


  \subsection{A Guarded Interaction Tree Semantics}
\end{comment}

\section{Denotational Semantics}\label{sec:denotational-semantics}

In this section, we discuss how to adapt the denotational model of Giovannini et
al. \cite{gtt-sgdt} to accommodate algebraic effects. After reviewing the key
aspects of their model, we discuss the generalizations necessary to extend their
semantics to accommodate an arbitrary algebraic theory. We additionally correct
a subtle limitation with their semantics related to the denotation of certain
casts. While this limitation was not problematic for their result, this change
will prove crucial for us in relating the operational and denotational semantics
in Section \ref{sec:relating-the-semantics}.

% Recall that in their model, pure values were modeled as elements of a
% \emph{predomain}, and effectful computations were modeled as elements of an
% \emph{error domain}, which is a predomain equipped with algebraic operations
% supporting error and computational steps. Our approach will be to generalize the
% notion of error domain to be an algebra of a suitable endofunctor corresponding
% to the signature of effects being modeled. Then the free error domain
% construction will be generalized to construct the free such algebraic structure,
% as an \emph{inductive-guarded interaction tree}.


\subsection{Review of the Prior Denotational Model}

% In this section, we discuss how to build a denotational model of effectful
% gradual typing from scratch in guarded type theory.
%
% Recall that our framework is parameterized by an algerbaic theory $\calA$.
%
% Our goal will be to generalize the notion of \emph{error domain} from prior work
% to accommodate these effects.
%
We begin by reviewing the key ingredients in the denotational semantics of
Giovannini et al. \cite{gtt-sgdt}. For further discussion, see their paper. For
the sake of brevity, we will refer to their model as the \emph{prior guarded
model} of gradual typing.

The prior model is based on call-by-push-value \cite{levy99}, a categorical
model that refines Moggi's monadic semantics \cite{MOGGI199155}. In
call-by-push-value, there is a category of value types $\calV$ and pure
morphisms $f : A \to A'$, and a category $\calE$ of computation types $\phi : B
\multimap B'$. A pair of adjoint functors $F \vdash U$ mediates between the two
categories, where $F : \calV \to \calE$ takes $A$ to the type $F A$ of
computations that return values in $A$, and $U : \calE \to \calV$ takes $B$ to
the type $U B$ of suspended computations of type $B$.

In the prior model, the value types are modeled by sets equipped with a partial
ordering relation that models type precision on closed values. However, since
step-insensitive reasoning does not play well with transitivity, each type also
has an additional relation called \textbf{weak bisimilarity}. They call such a
set a \textbf{predomain}. Morphisms of predomains are functions on the
underlying sets that respect both of the relations.
%
The precise definitions are given below:
%
\begin{definition}
A \textbf{predomain} $A$ consists of a set $A$ along with two relations:
\begin{itemize}
    \item A partial order $\ltdyn_A$.
    \item A reflexive, symmetric ``bisimilarity'' relation $\bisim_A$.
\end{itemize}

A \textbf{morphism of predomains} $f : A \to A'$ is a function between the
underlying sets that is both \emph{monotone} (if $x \ltdyn_A x'$, then $f(x)
\ltdyn_{A'} f(x')$), and \emph{bisimilarity-preserving} (if $x \bisim_A x'$,
then $f(x) \bisim_{A'} f(x')$).
\end{definition}

Computation types are equipped with the same two relations, in addition to a
special distinguished error element that is the least element in the ordering,
as well as a distinguished morphism of predomains $\theta_B : \later B \to B$.
%
The precise definitions are given below:
%
\begin{definition}
An \textbf{error domain} $B$ consists of a predomain $B$ along with the
following data:
\begin{itemize}
    \item A distinguished ``error'' element $\mho_B \in B$ such that $\mho
    \ltdyn_B x$ for all $x$
    \item A morphism of predomains $\theta_B \colon \later B \to B$
    \item For all $x, y : B$ with $x \bisim_B y$, we have $\theta_B(\nxt\, x)
    \bisim_B y$
\end{itemize}
\end{definition}

A \textbf{morphism of error domains} is a morphism of the underlying predomains
that preserves the error element and the $\theta$ map.
%
For an error domain $B$, we define the predomain morphism $\delta_B := \theta_B
\circ \nxt$.

In the prior model, type precision is modeled as a kind of heterogeneous
relation on the denoted predomains. More concretely, A \textbf{predomain
relation} $c : A \rel A'$ is a relation between $A$ and $A'$ that is upward
closed in the ordering on $A'$ and downward closed in the ordering on $A$
(no condition is required involving the bisimilarity relations).
%
\begin{definition}
  A \textbf{predomain relation} $c : A \rel A'$ is a relation $c$ between the
  underlying sets of $A$ and $A'$ satisfying:
  \begin{enumerate}
    \item (upward closure): For all $x \in A$ and $y, y' \in A'$, if $x
    \binrel{c} y$ and $y \ltdyn_{A'} y'$, then $x \binrel{c} y'$.
    \item (downward closure): For all $x, x' \in A$ and $y \in A'$, if $x'
    \ltdyn_{A} x$ and $x \binrel{c} y$, then $x' \binrel{c} y$.
  \end{enumerate}
\end{definition}
%
The ``reflexive'' predomain relation is written $r(A)$ and is given by the
ordering $\ltdyn_A$. The relational composition of two predomain relations is
also a predomain relation, and $r(A)$ is the identity for this composition.

Error domain relations are predomain relations that are additionally a
congruence with respect to the operations $\theta$ and $\mho$:
%
\begin{definition}
  An \textbf{error domain relation} $d : B \rel B'$ is a relation $d$ on the
  underlying predomains satisfying
  \begin{enumerate}
     \item ($\mho$ congruence): $\mho_B \binrel{d} \mho_{B'}$.
     \item ($\theta$ congruence): For all $\tilde{x}$ in $\laterhs B$ and
     $\tilde{y} \in \laterhs B'$, if $\laterhs_t (\tilde{x}_t \binrel{d}
     \tilde{y}_t)$, then $(\theta_B (\tilde{x}) \binrel{d} \theta_{B'}
     (\tilde{y}))$. 
  \end{enumerate}
\end{definition}
%
There is again a reflexive error domain relation $r(B)$ given by the ordering.
The relational composition of two error domain relations in general is not an
error domain relation ($\theta$ congruence in particular is not preserved), but
we instead use an alternative ``free composition'' $dd' : B \rel B''$ that is
inductively generated from the relational composition.

Lastly, term precision is modeled via the notion of \emph{square}, which
captures the logical-relations style relationhip between the functions denoted
by two terms. More concretely:
%
\begin{definition}
Let $f : A_i \to A_o$ and $g : A_i' \to A_o'$ be predomain morphisms, and let
$c_i : A_i \rel A_i'$ and $c_o : A_o \rel A_o'$ be predomain relations. There is
a \textbf{predomain square} $f \ltsq{c_i}{c_o} g$ if for all $x \binrel{c_i} y$,
we have $f(x) \binrel{c_o} g(y)$. 

Likewise, an \textbf{error domain square} $\phi \ltsq{d_i}{d_o} \psi$ is defined
similarly, with no extra interaction with the algebraic structure.
\end{definition}
%

\subsection{Generalizing to Algebraic Effects}

In the remainder of this section, we discuss the changes we must make to the
semantics of Giovannini et al. to accommodate arbitrary algebraic effects. The
key question we answer is how to generalize the notion of error domain from
their model to account for the presence of effects besides error and
computational steps.

\subsubsection{Models of Guarded Algebraic Theories in Predomains}

As a first step, we work over an arbitrary guarded algebriac theory that has
suitable operations for error and steping. Even though our framework is
parameterized by an \emph{ordinary} algebraic theory, everything in this section
applies to the more general setting of guarded algebraic theories. This section
can be seen as analogous to Section \ref{sec:guarded-algebraic-theories}, except
now, instead of defining models of guarded algebraic theories in the category of
sets, we define models in the category of predomains and their morphisms.

Let $\calA = (\Sigma, E)$ be a guarded algebraic theory (see Section
\ref{sec:guarded-algebraic-theories} for the relevant definitions). Recall that
$\Sigma$ is a signature consisting of a type of operation names $\Op$, where
each $\oper : \Op$ has ordinary arity $\tau_\oper$ and guarded arity
$\gamma_\oper$. (For brevity, we write $\oper : (\tau_\oper, \gamma_\oper)$.)

Furthermore, suppose that $\calA$ has a distinguished error operation $\err :
(0, 0)$ and a distinguished step operation $\theta : (0, 1)$.

In the following, we will need the fact that the induced polynomial endofunctor
$\Sigma : \type \to \type$ extends to an endofunctor on the category of
predomains and their morphisms. In particular, for a predomain $X$, we define
%
\[ \Sigma(X) := \Sigma_{\oper: \Op}, (\tau_\oper \to X) \times (\gamma_\oper \to \later X). \]
%
In the above, the ordering on $\later X$ is defined such that $\tilde{x}
\ltdyn_{\later X} \tilde{y}$ if one time step later, the values are related in
$X$. Likewise for the bisimilarity relation.
%
In addition, the type of functions $\tau_\oper \to X$ inherits a predomain
structure from that of $X$: the ordering and bisimilartity relations are defined
pointwise using the corresponding relations for $X$.

We define the action of $\Sigma$ on morphisms as follows. Given a function $f :
X \to Y$, we define $\Sigma(f) : \Sigma(X) \to \Sigma(Y)$ by
%
\[ \Sigma(f)(\oper, (u, g)) = (\oper, (u', g')), \]
%
where $u' : \tau_\oper \to Y$ is defined by $u' = f \circ u$ and $g' :
\gamma_\oper \to \later Y$ is defined using the functorial action of $\later$ by
$g' = (\later f) \circ g$.

The analogue of error domains in our semantics will then be an algebra $B$ of
the endofunctor $\Sigma$. Using the action of $\Sigma$ on predomains defined
above, an algebra of $\Sigma$ is a predomain $B$ equipped with a morphism of
predomains $\Sigma(B) \to B$. 


% a distinguished error element $\mho$ and a distinguished $\theta_B \colon \later
% B \to B$. 

% As in the definition of an error domain from the denotational model of
% Giovannini et al., we require that $\mho$ is the smallest element in the
% ordering, that $\theta$ respects the relations (i.e., it is a morphism of
% predomains), and that $\delta_B := \theta_B \circ \nxt$ is bisimilar to the
% identity.

% As in the definition of an error domain, we also require that the algebra comes
% with ordering and a bisimilarity relations, that $\mho$ is the smallest element
% in the ordering, that $\theta$ respects the relations (i.e., it is a morphism of
% predomains), and that $\delta_B := \theta_B \circ \nxt$ is bisimilar to the
% identity.


More formally, we make the following definition:
%
\begin{definition}\label{def:sigma-domain} 
  
  Let $\calA = (\Sigma, E)$ be a guarded algebraic theory as defined above. An
  \textbf{$\calA$-domain} is a predomain $B$ with a morphism of predomains
  $\alpha : \Sigma(B) \to B$. In particular, for each $\text{op} \in
  \Sigma.\Op$, we have a morphism of predomains
  %
  \[ \sem{\oper} : (\tau_\oper \to B) \times (\gamma_\oper \to \later B) \to B. \]
  %
  
  Furthermore, each equation $e = (\Gamma_e, l_e, r_e)$ in $E$ holds when
  interpreted in $B$.

  Lastly, the predomain $B$ has a distinguished element $\mho_B$ such that
  $\mho_B \ltdyn_B x$ for all $x \in B$, along with a distinguished morphism of
  predomains $\theta_B : \later B \to B$, such that if $x \bisim_{B} y$, then
  $\delta_B \bisim_{B} y$ where $\delta_B := \theta_B \circ \nxt$.
\end{definition}
%
As with ordinary error domains, there are corresponding notions of morphisms,
relations, and squares of $\calA$-domains. Before introducing these, we first
discuss the action of the endofunctor $\Sigma$ on predomain morphisms and
predomain relations.

Let $c : A \rel A'$ be a predomain relation. The action of the endofunctor
$\Sigma$ on $c$, written $\Sigma(c)$, is a predomain relation between
$\Sigma(A)$ and $\Sigma(A')$, defined as follows:
%
First, note that any predomain relation $c : A \rel A'$ extends to a predomain
relation $\tau \to c$ between $\tau \to A$ and $\tau \to A'$ defined pointwise,
i.e., $f \binrel{(\tau \to c)} f'$ iff for all $s : \tau$ we have $(f s)
\binrel{c} (f' s)$.
%
Similarly, given $c : A \rel A'$ we define $\later c$ such that $\tilde{x}
\binrel{(\later c)} \tilde{y}$ iff after one step passes, the left and right
sides (now in $A$ and $A'$ respectively) are related in $c$.
%
We now define $(\oper, (u, g)) \binrel{(\Sigma(c))} (\oper', (u', g'))$ iff
$\oper = \oper'$ and $u \binrel{(\tau_\oper \to c)} u'$ and 
$g \binrel{(\gamma_\oper \to \later c)} g'$.


Now we can introduce the definitions of morphisms and relations of
$\calA$-domains.
%
\begin{definition}\label{def:sigma-domain-mor-rel-sq} Let $B$ and $B'$ be
  $\calA$-domains. A \textbf{morphism of $\calA$-domains} from $B$ to $B$' is a
  morphism $\phi$ of the underlying predomains that preserves the algebraic
  structure.
  %
  Spelling this out concretely, we have:
  %
  $\phi (\sem{\oper}_B(u, g)) = \sem{\oper}_{B'}((\phi \circ u), ((\later \phi) \circ g))$,

  A \textbf{relation of $\calA$-domains} $d : B \rel B'$ is a relation $|d|$ of
  the underlying predomains that is a congruence with respect to the algebraic
  structure. In particular:
  %
  If $x \binrel{(\Sigma(|d|))} y$, then $\sem{\oper}_B(x) \binrel{|d|} \sem{\oper}_{B'}(y)$.

  A \textbf{square of $\calA$-domains} $\phi \ltsq{d}{d'} \phi'$ is defined as a
  square of the underlying predomains. 
\end{definition}

\begin{remark}
  Spelling out the definition of morphism for the distinguished guarded stepping
  operation $\theta : \later B \to B$, we get that
  %
  $\phi(\theta_{B}(\tilde{x})) = \theta_{B'}((\later \phi)(\tilde{x}))$.

  The requirement that $\mho_B$ be related to $\mho_{B'}$ in the definition of
  relation may seem weaker than the corresponding requirement in the definition
  of a relation of error domains in the previous section, where we required that
  $\mho_B \binrel{|d|} y$ for \emph{all} $y \in B'$. However, the former implies
  the latter using the fact that predomain relations are closed under the
  ordering on both sides and that $\mho_{B'} \ltdyn y$ for all $y$.
\end{remark}

The usual relational composition of $\calA$-domain relations is not in general a
$\calA$-domain relation. This is because congruence with respect to guarded
operations (such as $\theta : \later B \to B$) is not preserved by the
relational composition. Thus, we instead define a free composition $dd' : B \rel
B''$ that is inductively generated from the relational composition. See the
Appendix for the precise definition.
% TODO cite appendix.

\subsection{Perturbations and Quasi-Representable Relations}

We next review how the model of Giovannini et al. validates the graduality
property. In the New-Licata denotatinoal model of graduality on which Giovannini
et al. based their guarded model, the behavior of casts was encoded in the
semantics by \emph{representable relations}. In prior work, these were
sufficient for interpreting the four cast rules crucial for graduality. At a
high level, type precision $c : A \ltdyn A'$ is modeled by representable
relations, i.e., relations that behave like the graph of a function, e.g., $x
\binrel{\sem c} y$ if and only if $\sem{\upc c}(x) \leq_{A'} y$
\cite{shulman2007}.

However, as Giovannini et al. showed, the notion of representability cannot be
applied unchanged to the guarded setting because of a fundamental tension
between transitivity and step-insensitive reasoning \cite{gtt-sgdt} (see Section
4 of their paper for details). As a result, the authors worked with a transitive
but step-sensitive ordering where terms must be in lock-step. They weakened the
representability requirement to a notion that they called
\emph{quasi-representability}, which is essentially representability up to an
explicit delay. They then developed a theory of these delays, which they called
\emph{syntactic perturbations}.

More concretely, in the prior guarded semantics, value types denote predomains
along with a monoid of syntactic perturbations $M_{A}$ and an interpretation
homomorphism $i_A : M_{A} \to \semptb{A}$, where $\semptb{A}$ is the set of
endomorphisms on $A$ that are weakly bisimilar to the identity. Computation
types were defined analogously. The idea is that the perturbations mimic the
steps taken by casts, allowing the perturbed term to be in lock-step with a
cast.
% and because they are specified as syntactic data, it is possible to associate 

Then to model type precision, the authors required that each predomain relation
$c : A \rel A'$ is quasi-left-representable, and $Fc$ is
quasi-right-representable. Likewise, error domain relations $d : B \rel B'$ are
required to be quasi-right-representable and $Ud : UB \rel UB'$ must be
quasi-left-representable.
%
Lastly, in order to make the notions of quasi-representability compositional,
the authors additionally require that the relations must interact with the
perturbations on either side, in that given $c : A \rel A'$, it is possible to
push a perturbation on $A$ to one on $A'$ and pull a perturbation on $A'$ to one
on $A$.
%
On the other hand, morphisms are not required to interact with the
perturbations.
% TODO more detail on the definition of quasi-representability?

It turns out that the construction of perturbations and quasi-representable
relations is independent of the details of the adjunction between predomains and
(our generalized notion of) error domains. This means that the constructions
used in the model of Giovannini et al. will work in our more general setting
without any major modifications.
%
Details as to the definitions of the perturbations associated with each of the
type constructors can be found in the Appendix. % TODO cite appendix
%
The intuitive reason for why no modifications are needed is that the casts do
not cause side effects other than stepping or error. Thus, the perturbations for
$F(A)$ can still be defined as $\mathbb{N} \oplus M_A$.



\subsection{Adjusting the Term Semantics}\label{sec:term-relational-semantics}

% ...we need to make an important technical modification to the model of Giovannini
% et al.
In addition to generalizing error domains, we need to make an important
technical modification to our model concerning the behavior of casts. This change is
motivated by our goal of relating the denotational and operational semantics via
an adequacy result. In the prior model, the denotations of certain casts
involves additional perturbations (i.e., computational steps) in order to place
terms involving those casts in lock-step. These additional steps ultimately did
not matter for the final result, since the graduality property is
step-insensitive. However, intensionally, these casts behave differently from
the corresponding casts without the extra steps.

This difference in intensional behavior presents a difficulty when relating
their denotational model to an operational semantics: in order for the two
semantics to be in lock-step, we must define the operational semantics in such a
way that the casts take additional steps just as they do in the denotational
semantics. But this requirement is not desirable, as it restricts the kinds of
operational semantics to which our framework applies to those that happen to
feature this non-standard cast semantics. Users of our framework would need to
contort their operational semantics into this form if they would like to apply
our adequacy result.

Our remedy for this situation is to note that there is an additional degree of
freedom in the model of Giovannini et al. that up to this point has remained
unexploited. Leveraging this, we define two \emph{different} cast semantics, one
where the casts involve perturbations, and one where they do not. The former is
used in constructing the semantic proofs of the cast axioms, where being in
lock-step is crucial. The latter is what we use for the denotational term
semantics of casts, allowing the denotations of casts to be in lock-step with
the casts in a typical operational model, i.e., where the casts do not take
additional steps. While the intensional behavior of these two kinds of casts are
distinct, they are extensionally related in that they are \emph{weakly
bisimilar}, which is sufficient for the final proof of graduality.

For concreteness, recall that in the prior model, term precision $M \ltdyn N :
A$ is modelled by $\sem{M} \ltbisim \sem{N}$ where
%
\[ f \ltbisim g := \exists f', g'. f \bisim f' \ltdyn g' \bisim g. \]
%
That is, we do not require that $f$ and $g$ are in lock-step; we are allowed to
use any $f'$ bisimilar to $f$ and $g'$ bisimilar to $g$ such that $f'$ is
lock-step related to $g'$ in the error ordering.
%
This is useful because the downcast associated to the composition of type
precision $c : A \ltdyn A'$ and $c' : A' \ltdyn A''$ is a morphism from
$F(\sem{A''})$ to $F(\sem{A})$, but in the prior model it is not simply the
composition of the downcast morphisms for $c$ and $c'$. Instead, an additional
perturbation is inserted in order to make the relational squares line up in the
proof of quasi-representability. More details can be found in the Appendix of
their paper \cite{gtt-sgdt-extended-version}.

On the other hand, observe that we could instead define the downcast morphism
for the composition $cc'$ to be the composition of the corresponding downcasts,
without the additional perturbation, and then note that the resulting morphism
is \emph{weakly bisimilar} to the version with the additional perturbation. This
is sufficient for graduality, since the additional perturbation will be absorbed
by the weak bisimilarity in the definition of $f \ltbisim g$.

% Using the above definition for term precision, we see that the additional
% perturbations in the cast corresponding to a composition will be absorbed by the
% weak bisimilarity.

With this observation in mind, we augment the denotation of type precision in
the prior denotational model by associating to each value relation $c$ in the
semantics two morphisms: a ``relational'' embedding morphism $e$ such that $c$
is quasi-left-representable by $e$, and a ``term'' embedding morphism $e' \bisim
e$. We similarly associate two projection morphisms to each computation relation
$d$. We use the term versions of the embedding and projection morphisms as our
denotations of upcasts and downcasts, and we use the relational versions
internally in the proof of graduality.
%
The result is that our term semantics for casts no longer involves unnecessary
perturbations that were inserted to make the graduality proof work.

% TODO cite appendix
Details on the changes made to the model can be found in the Appendix.


\section{Relating the Operational and Denotational Semantics}\label{sec:relating-the-semantics}

Having specified the operational and denotational semantics, we now seek to
prove a theorem relating them. Our goal is to show that the denotatinoal model
is \emph{adequate} for the operational semantics, that is, if the denotations of
two terms are equal, then their operational behavior is the same.

More concretely, in Section \ref{sec:operational-semantics}, we defined a
function $\run : \term{A} \to T_{\mho sb} (\val{A})$ taking a closed term to its
final behavior. Then in Section \ref{sec:denotational-semantics} we defined the
denotation of a term as a function $\sem{\cdot} : \term{A} \to T_{\mho
s}(\sem{A})$. Our goal is to relate the operational and denotational semantics.
In particular, we seek to show that for closed terms of type $\nat$, we have
%
\[ \run(M) \equiv_b \sem{M}, \]
%
where $\equiv_b$ is a relation between $T_{\mho sb}(\mathbb{N})$ and $T_{\mho
s}(\mathbb{N})$ that ignores the bureaucratic delays in its first argument.

To that end, we will define a logical relation between the denotations of
(closed) terms of type $A$ and syntactic terms of type $A$:
%
\[ \calE_A : \Rel {T_{\mho s}\sem{A}} {(\term A)}, \]
%
and then $\equiv_b$ is defined as $\calE_\nat$ since we have $\sem{\nat} = \mathbb{N}$.

\subsection{The Logical Relation}

% Note: here we *do* care about the fact that the theories are built using
% sum/tensor, not merely that they have error and stepping operations.

Before defining the logical relation, we will need the notion of a
\emph{heterogeneous lifting} which takes a relation $R : X \rel Y$ and lifts it
to $\bbt(R) : T_{\mho s}(X) \rel T_{\mho sb}(Y)$. The lifted relation relates
bureaucratic steps in its right argument to the identity in its left argument.
We define the lifted relation as the least relation that (1) extends $R$, (2) is
a congruence with respect to the operations in $\calA_{\mho s}$ (i.e., the
operations of $\calA$, error, and \emph{true} stepping), and (3) relates the
identity in $T_{\mho s}(X)$ to $\delta_b$ in $T_{\mho sb}(Y)$. The precise
definition is given in Figure \ref{fig:relational-lifting}.

\begin{figure}
  \begin{mathpar}
    \inferrule*[right=RetRet]
    { x \binrel{R} y }
    { \eta^{T_{\mho s}}(x) \binrel{(\bbt(R))} \eta^{T_{\mho sb}}(y) }

    \inferrule*[right=OpOp]
    { \oper \in \calA_{\mho s}.\Op 
      \and \forall s.\, (u\, s) \binrel{(\bbt{R})} (u'\, s) \\
      \and \forall s.\, (g\, s) \binrel{(\later(\bbt{R}))} (g'\, s)  
    }
    { (T_{\mho s}(X).\oper(u, g)) \binrel{(\bbt{R})} (T_{\mho sb}(Y).\oper(u', g')) }

    \inferrule*[right=IdStep]
    { \bar{x} \binrel{(\bbt(R))} \bar{y} }
    { \bar{x} \binrel{(\bbt(R))} \delta_b(\bar{y}) }

  \end{mathpar}

  \caption{Relational Lifting that Ignores Bureaucratic Steps}
  \label{fig:relational-lifting}
\end{figure}


The logical relation is defined in Figure \ref{fig:logical-relation}. We define
a \textbf{value relation}
%
$\calV_A : \Rel{\sem{A}} {\val{A}}$,
%
and an \textbf{expression relation} 
%
$\calE_A : \Rel{T_{\mho s}\,\sem{A}}{\term{A}}$.
%
While logical relations are typically defined by induction on the type $A$, the
case for $\dyn$ is generated by a combination of induction and guarded
recursion. Note that while the product case for $\dyn$ refers to the value
relation at the larger type $\dyn \times \dyn$, this is merely syntactic sugar;
to be completely rigorous, we would simply \emph{define} the relation
$\calV_{\dyn \times \dyn}$ in an analogous manner to the definition for the
product case $\calV_{A \times A'}$.
% TODO explain this further?


\begin{figure}
  \begin{mathpar}
    \begin{array}{rcl}
      \calV_\nat (n, V_n) &:=& (n = V_n) \\
      \calV_{A_i \ra A_o} (f, V_f) &:=& 
        \forall (x \binrel{\calV_{A_i}} V).\, (f\, x) \binrel{\calE_{A_o}} (V_f\, V) \\
      \calV_{A \times A'} ((x, y), V_\times) &:=& \exists V_1, V_2.
        (V_\times = (V_1, V_2)) 
        \wedge (x \binrel{\calV_{A}} V_1) 
        \wedge (y \binrel{\calV_{A'}} V_2) \\
    \end{array}

    \inferrule*[right=DynProd]
    { \vrel{\dyn \times \dyn}{(d_1, d_2)}{V_{d\times}} }
    { \vrel{\dyn}{\itimes(d_1, d_2)}{ \tagtimes{V_{d\times}}} }

    \inferrule*[right=DynFun]
    { \forall x\, V.\, (\later(\vrel{\dyn}{x}{V})) \to \erel{\dyn}{(f\, x)}{(V_f\, V)} }
    { \vrel{\dyn}{\iarr(f)}{\tagarr{V_f}} }


  \begin{array}{rcl}
      \erel {A} {\bar{x}} {N} &:=& \liftvrel{A}{\bar{x}}{(\run N)} % \bar{x} \binrel{(\bbt(\calV_{A}))} (\run{N})
  \end{array}

  \end{mathpar}

  \caption{Logical Relation}
  \label{fig:logical-relation}
\end{figure}

The logical relation satisfies the following properties.
%
First, we have \emph{anti-reduction}:
%
\begin{lemma}[Anti-Reduction]\label{lem:anti-reduction}

  Let $\cdot \vdash M : A$, and suppose $M \stepsin{k} \cpure{M'}$. If
  $\erel{A}{\xbar}{M'}$, then $\erel{A}{\xbar}{M}$.

\end{lemma}

Next, we have a semantic bind lemma for the logical relation, which allows us to
replace arbitrary terms in an evaluation context with values.
%
\begin{lemma}[Monadic Bind]\label{lem:monadic-bind}
  
  % Suppose for all values $V, W$ that $\vrel{A}{V}{W}$ implies
  % $\erel{B}{\phi(V)}{(E[W])}$. Then, for all $\erel{A}{\bar{x}}{N}$, we have
  % $\erel{B}{(\phi(\bar{x}))}{(E[N])}$.
  Let $\phi : T_{\mho s}(\sem{A}) \to T_{\mho s}(\sem{B})$.
  %
  Suppose for all values $x, V$ that $\vrel{A}{x}{V}$ implies
  $\erel{B}{\phi(T_{\mho s}.\eta(x))}{(E[V])}$.
  %
  Then for all $\bar{x} \in T_{\mho s}(\sem{A})$ and terms $N$ of type $A$ such that
  $\erel{A}{\bar{x}}{N}$, we have $\erel{B}{(\phi(\bar{x}))}{(E[N])}$

\end{lemma}


\subsection{Fundamental Lemma}

% TODO define \grel{\Gamma}

% $\gambar$ maps $x : A \in \Gamma$ to $\sem{A}$ and $\gamma'$ maps $x : A \in \Gamma$ to $\val{A}$

We begin by defining the extension of the value relation to contexts.
Concretely, we define $\grel{\Gamma}{\gambar}{\gamma'}$, where $\gambar :
\sem{\Gamma}$ and $\gamma'$ is a mapping of variables $x : A \in \Gamma$ to
$\val{A}$, as the pointwise extension of $\vrel{}{}{}$ to the elements in
$\gambar$ and $\gamma'$.
%
% The definition is by induction, with rules
% %
% $\grel{\bullet}{\_}{\_} := \top$, and
% %
% $\grel{\Gamma, x : A}{\bar{x} :: \bar{xs}}{[x \mapsto V, xs]} :=
%   \vrel{A}{\gambar(x)}{\gamma'(x)} \wedge \grel{\Gamma}{\bar{xs}}{xs}$.
%
Now we define the extension of the logical relation to morphisms $\fbar :
\sem{\Gamma} \to T_{\mho s}(\sem{A})$ and open terms $\Gamma \vdash M : A$, as
follows:
\[ 
  \Gamma \vDash \fbar \sim M : A := 
    \forall (\grel{\Gamma} {\gambar} {\gamma'}).\,
      \erel{A} {(\fbar\, \gambar)} {M[\gamma']}.
\]
%
We now state the fundamental lemma for the logical relation, which implies the
result that $\run(M) \equiv_b \sem{M}$:
%
\begin{lemma}[Fundamental Lemma]
  Let $\Gamma \vdash M : A$. Then $\Gamma \vDash \sem{M} \sim M : A$.
\end{lemma}
%
For the proof of this lemma, see the Appendix. % TODO cite appendix
%
The above implies the adequacy of the denotational semantics with respect to the
operational semantics:
\begin{theorem}[Adequacy]
  Let $\cdot \vdash M : \nat$. Then we have $\run{M} \equiv_b \sem{M}$.
\end{theorem}

% \subsection{Summary of Results So Far}

% Logical relation
% Fundamental Lemma
% Adequacy of the term model
% Adequacy of the relational model


% To this end, we will define a logical relation between (closed) syntactic terms
% of type $A$ and their denotations as elements of $F\sem{A}$, where $F$ is the
% functor constructing the free error-effect algebra as discussed in Section
% \ref{sec:denotational-semantics}.


\section{Globalization and Extensional Collapse}\label{sec:globalization}

The last step of the graduality proof involves using the results obtained in the
prior sections to derive extensional properties of the semantics that do not
mention any features of guarded type theory. The result is a statement about the
operational behavior of terms that can be made in ordinary type theory. The
process we will follow here involves globalization \cite{atkey-mcbride2013},
which is a process that takes a construction in guarded type theory and turns it
into a coinductive object in ordinary type theory. Not all constructions in
guarded type theory can be globalized into coinductive definitions, but the
specific definitions we will be globalizing do satisfy the required technical
condition.
% TODO: mention the notion of a functor commuting with clock quantification.

Unfortunately, carrying out the process of globalization for an arbitrary
algebraic theory $\calA$ appears difficult, and we are not aware of an approach
that works at this level of generality. On the other hand, there are two special
cases where we can explicitly describe the globalization process, which we
describe below.

\subsection{Monad Transformers}

The first special case occurs when we know that the input theory $\calA$ is
constructed from certain ``base'' theories using sum and tensor. The idea is
that in combining one of these base theories with another theory $\calA$, we
obtain for free a corresponding monad transformer taking the free-model monad
for $\calA$ to the free-model monad of the combined theory. 

One example of this phenomenon involves the theory of exceptions:
%
\begin{lemma}[Error Monad Transformer]\label{lem:error-monad-transformer} Let
$\calA$ be a guarded algebraic theory. If $T_\calA$ is the free-model monad for
$\calA$, then the free-model monad for the theory $\calA + \th{\mho}$ is given
by
%
\[ T'(A) := T(A + 1). \] 
\end{lemma}

Another example of this phenomenon involves the theory of global mutable state:
%
\begin{lemma}[State Monad Transformer]\label{lem:state-monad-transformer}
Let $\calA$ be a guarded algebraic theory. Suppose that $\calA$ has an operation
$\theta$ of arity $(0, 1)$. If $T_\calA$ is the free-model monad for $\calA$,
then the free-model monad for the theory $\calA \otimes \th{\text{State}_S}$ is
given by
%
\[ T'(A) := S \to T_\calA(S \times A). \]
%
\end{lemma}

Such results allow us to build up the free-model monad in a compositional
manner. If the innermost monad is the usual guarded lift monad $L(X) \cong X +
\later L(X)$, then we can then globalize the overall monad and use known results
about the globalization of the guarded lift monad to obtain a characterization
of the combined monad.

\subsection{Trace Semantics}

If the provided algebraic theory has no equations, then the free-model monad for
the theory $T(X)$ is isomorphic to the type of unquotiented terms with variables
in $X$.

To motivate the forthcoming definitions, we begin with a simple example where
the only effects are error and computational steps. Consider the guarded
algebraic theory $\calA = \th{\mho} + \th{\step}$ (see Section
\ref{sec:guarded-algebraic-theories} for the definitions). In particular, note
that there are two operations $\mho : (0, 0)$ and $\theta : (0, 1)$. The
free-model monad for this theory is the unique type $T(X)$ solving the
guarded-recursive equation
%
\[ T(X) \cong X + 1 + \later (T(X)). \]

We now review the notion of \emph{clocks}, which are a tool allowing us to
extract coinductive definitions from guarded ones. Clocks were introduced by
Atkey and McBride \cite{atkey-mcbride2013}. The idea is that all guarded
constructions are indexed by a clock variable $k$, e.g., $\later^k A$ is the
type of $A$-values available after one time step on clock $k$.
%
The key aspect of clocks is that universal quantification over clocks gives a
means to turn a guarded definition into a coinductive one
\cite{atkey-mcbride2013, kristensen-mogelberg-vezzosi2022}.

Define $T^{gl}(X) := \forall k. T^k(X)$. Then provided $X$ satisfies a technical
condition known as \emph{clock-irrelevance}, we have
%
\begin{align*}
  T^{gl}(X) &= \forall k. (X + 1 + \later T^k(X)) \\
  &\cong \forall k. X + \forall k. 1 + \forall k. \later T^k(X) \\
  &\cong X + 1 + \forall k. T^k(X).
\end{align*}
%
Moreover, $T^{gl}(X)$ is the final coalgebra for the functor $F_X(Y) = (X + 1) +
Y$. This follows from a theorem of Kristensen et al.
\cite{kristensen-mogelberg-vezzosi2022}. The final coalgebra for $F_X$ can be
expressed using Capretta's \emph{delay monad} $D(X)$, which is defined as the
final coalgebra of the functor $F_X(Y) = X + Y$ \cite{lmcs:2265}. Elements of
$D(X)$ can be thought of as computations that may either return a value in $X$,
or take an observable computational step. Such computations may run forever,
which makes $D(X)$ a constructive method for modelling non-termination in type
theory.

Altenkirch et al. \cite{Altenkirch-Danielsson-Kraus-2017} show that the delay
monad $D(X)$ is equivalent to the type of \emph{monotone sequences} $\mathbb{N}
\to (X + 1)$ (see Section 4.1 of their paper). By a monotone sequence, we mean a
function $f : \mathbb{N} \to (X + 1)$ satisfying that if $f(n) = \inl x$, then
$f(n+1) = \inl x$. The intuition is that while we cannot decide in general
whether an arbitrary computation terminates or runs forever, for any
\emph{finite} amount of ``fuel'' $n$, we can ask whether it has terminated in
$n$ steps or fewer. If it has, then $f$ tells us the value with which the
computation terminated; if not, we are simply told that it is still running.
Monotonicity says that these answers must be consistent: if the computation
terminates in $n$ or fewer steps with result $x$, then it certainly also
terminates in $n+1$ or fewer steps with the same result $x$.

Observe that the delay monad is an intensional representation of
potentially-nonterminating computations. That is, two computations that return
the same value but take a different number of steps are considered distinct.
Since the statement of graduality is oblivious to the number of steps a program
takes, the delay monad is not suitable as a denotational domain for closed
programs. To remedy this situation, one may consider the delay monad quotiented
by weak bisimilarity. However, in order to show that the resulting type is a
monad in type theory, it is believed that the axiom of countable is required
\cite{TODO}.

To generalize the situation to arbitrary effects, we assume that our theory
$\calA$ contains no equations. Then we consider the theory $\calA + \th{\mho} +
\th{\step}$. Then as above, we form the free-model monad for this theory,
and consider the globalized version by performing clock quantification.
%
The analogue of the monotone sequences for the simple setting of error and
stepping will be functions that take a generalized notion of fuel and return a
finite trace of the computation.


\section{Discussion}\label{sec:discussion}

% Comparison with guarded interaction trees
% Extension to arbitrary monads (i.e., not just those arising from an algebraic theory)
% Comparison to Kavvos's approach
% Extension to Kripke models
% Discuss Møgelberg + Paviotti paper on denotational semantics of FPC in SGDT
% Discuss work on runners and comodels
% Discuss What Monads Can and Cannot Do paper

\subsection{Operational Semantics in Guarded Type Theory}

Møgelberg and Paviotti \cite{mogelberg-paviotti2016} use guarded type theory (in
particular, what they call Guarded Dependent Type Theory) to define a
denotational semantics for FPC (STLC extended with recursive types). They then
prove soundness and adequacy for their model using a logical relation
constructed in the type theory using guarded recursion. Particular care is taken
to ensure that the operational and denotational semantics are in lock-step. To
do this, the small-step operational semantics tracks the non-well-founded steps,
and the logical relation ensures that each such step corresponds to a $\theta$
operation in the denotational semantics.

\subsection{Guarded Type Theory and Algebraic Effects}


\subsection{Interaction Trees}

Interaction trees \cite{TODO} are a popular framework for representing effectful
programs in dependent type theory. While the original formulation of interaction
trees was as a coinductive type, this has been extended to \emph{guarded}
interaction trees which are capable of modelling higher-order effects
\cite{guarded-interaction-trees}.
% TODO more to say?

The explicit construction we give of the free model for a guarded algebraic
theory is closely related to guarded interaction trees, but is in one way more
general and in another way less so. Our approach is more general in the sense
that we allow a mix of inductive and guarded operations, whereas with guarded
interaction trees, all operations are guarded. On the other hand, guarded
interaction trees allow for negative occurrences guarded by a later; our
construction does not require this generality since in our notion of guarded
algebraic theories, all occurrences of the type are positive.

\subsection{Adequacy and Algebraic Effects}

Kavvos \cite{kavvos-2025} defines an operational and denotational semantics for
a language parameterized by a signature of effects, and then proves adequacy of
the denotational semantics for operational reasoning. The high-level methodology
is similar to our approach, with a few key differences. Similar to our approach,
Kavvos defines a small-step operational semantics and then extracts a big-step
semantics valued in interaction trees. However, the interaction trees are
defined coinductively, whereas the analogous construction in our approach is
defined using induction and guarded recursion.

Additionally, Kavvos's small-step semantics is inherently nondeterministic, in
that an effect operation can step to its continuation instantiated with all
possible response values. As a result, it is difficult to account for concrete
effectful small-step semantics, such as for state, where a pair of term and
state steps to a new term and a new state. In contrast, in our small-step
semantics, we allow for effectful steps using the endofunctor induced by the
algebraic signature $\Sigma$. As a result, a term is allowed to take an
effectful step

% big-step operational semantics using interaction trees and proves adequacy for
% an operational semantics constructed from

% Additionally, in Kavvos's approach, the small-step semantics is such that an
% effect operation steps nondeterministically...


\subsection{Future Work}

% other operations in the guarded algebraic theory besides just \theta

Our next goal is to extend the effectful semantics to account for more advanced
effects such as higher-order store or dynamic type-tag generation. This was the
original motivation of Giovannini et al., and with the experience gained in
incorporating algebraic effects and operational semantics into their framework,
we believe it will be feasible to extend the semantics to account for such
effects. There are two key challenges that must be overcome to handle
higher-order store. First, dynamic allocation must be modelled. In the
denotational setting this would take the form of a presheaf over a category of
worlds, where the world specifies the currently-allocated names. Second, the
ability to store functions makes the model viciously self-referential. Prior
work has investigated the use of guarded type theory to solve the resulting
domain equation \cite{TODO}. The challenge is to incorporate these approaches
into our denotational model, in particular, the relational aspects (e.g.,
squares) and the perturbations.

Another possible direction would be to adapt the model to allow for gradual
\emph{effect systems}, such as that of the language GrEff
\cite{new-giovannini-licata-2022}. This would necessitate changes to the
denotational semantics, including adding effect tracking. We would also need to
solve a domain equation for the dynamic effect type analogous to the equation
for the dynamic value type.



\bibliographystyle{ACM-Reference-Format}
\begin{DIFnomarkup}
\bibliography{\bibliographypath}
\end{DIFnomarkup}

% EXTENDED VERSION
\newpage
\appendix
% \section{Omitted Details and Proofs for Section \ref{sec:background}}


\subsection{Construction of the Free Algebra for a Signature}\label{sec:app:free-signature-monad}

In this section, we give an explicit construction for the functor that takes a
set $A$ to the free $\Sigma$-algebra $F_\Sigma(A)$ on $A$.
%
We seek to solve the equation
%
\[ F_\Sigma(A) \cong A + \Sigma(F_\Sigma(A)), \]
%
i.e.,
%
\[ F_\Sigma(A) \cong A + \Sigma_{\oper: \Op}, 
    (\tau_\oper \to F_\Sigma(A)) \times G_\oper(\later F_\Sigma(A)). \]
%
First, consider the parameterized inductive type 
%
\[ P(X) := \mu Y. (A + \Sigma_{\oper : \Op}, (\tau_\oper \to Y) \times G_\oper(X)). \]
%
Now, using guarded recursion, define $F_\Sigma(A)$ to be the unique solution to
the guarded-recursive equation
%
\[ F_\Sigma(A) \cong P(\later F_\Sigma(A)). \]
%
Then by the definitions of least fixpoint and guarded fixpoint, we have that
$F_\Sigma(A)$ is a solution to the original equation.

\section{Omitted Details and Proofs for Section \ref{sec:operational-semantics}}

\begin{comment}
\begin{lemma}
    Let $\hat{N} : \Sigmahat(\term{A})$. We have 
    %
    \[ \iota(f (\Sigmahat(E[\cdot])(\hat{N})) ) = (\eta' \circ E[\cdot])^{\ddag} (\iota(f(\hat{N}))). \]
\end{lemma}
\begin{proof}
    We proceed by cases on $\hat{N}$. If $\hat{N} = \cpure{M}$, then we have
    \begin{align*}
        \iota(f (\Sigmahat(E[\cdot])(\cpure{M})) )
        &= \iota(f (\cpure{E[M]}) ) \\
        &= \iota(\eta(E[M])) \\
        &= \eta'(E[M]),
    \end{align*}
    and
    \begin{align*}
        (\eta' \circ E[\cdot])^{\ddag} (\iota(f(\cpure{M})))
        &= (\eta' \circ E[\cdot])^{\ddag} (\iota(\eta(M))) \\
        &= (\eta' \circ E[\cdot])^{\ddag} (\eta'(M)) \\
        &= (\eta' \circ E[\cdot])(M) \\
        &= \eta'(E[M]).
    \end{align*}

    If $\hat{N} = \ceff{(\oper, u)}$, then we have
    \begin{align*}
        \iota(f (\Sigmahat(E[\cdot])(\ceff{(\oper, u)})) )
        &= \iota(f (\ceff{(\oper, E[\cdot] \circ u)}) ) \\
        &= \iota(T_{\mho}(\term{A}).\oper(\eta \circ E[\cdot] \circ u)) \\
        &= T_{\mho sb}(\term{A}).\oper(\iota \circ \eta \circ E[\cdot] \circ u) \\
        &= T_{\mho sb}(\term{A}).\oper(\eta' \circ E[\cdot] \circ u)
    \end{align*}
    and
    \begin{align*}
        (\eta' \circ E[\cdot])^{\ddag} (\iota(f(\ceff{(\oper, u)})))
        &= (\eta' \circ E[\cdot])^{\ddag} (\iota(T_{\mho}(\term{A}).\oper(\eta \circ u))) \\
        &= (\eta' \circ E[\cdot])^{\ddag} (T_{\mho sb}(\term{A}).\oper(\iota \circ \eta \circ u)) \\
        &= (\eta' \circ E[\cdot])^{\ddag} (T_{\mho sb}(\term{A}).\oper(\eta' \circ u)) \\
        &= (T_{\mho sb}(\term{A}).\oper(((\eta' \circ E[\cdot])^{\ddag}) \circ \eta' \circ u)) \\
        &= (T_{\mho sb}(\term{A}).\oper(((\eta' \circ E[\cdot])) \circ u)). \\
    \end{align*}

\end{proof}
\end{comment}


Notation in this section:
%
\begin{itemize}
    \item We let $\eta$ denote $\Tsb.\eta$.
    %
    \item We let $(-)^\dag$ denote Kleisli extension for $\Tsb$.
    %
    \item $f : (\term{A} + \Sigma_\mho(\term{A})) \to \Tsb(\term{A})$ is defined by
    cases, mapping $\cpure{M}$ to $\eta^{\Tsb}(M)$, and mapping $\ceff{(\oper, u, g)}$ 
    to $\Tsb(\term{A}).\oper(\eta \circ u, (\later \eta) \circ g)$.
    (Note that $u : \tau_\oper \to \term{A}$ and $g : \gamma_\oper \to \later (\term A)$.)
    %
    % \item $\iota : T_{\mho}(X) \to T_{\mho sb}(X)$ is the inclusion of the
    % free model monad $T_{\mho}$ into $T_{\mho sb}$. 
\end{itemize}

To begin, we define a function that ``runs'' an evaluation context $E$.
Specifically, if $E : A \multimap B$, then we define $\run{E} : \Tsb(\term A)
\to \Tsb(\val B)$ as
%
\[ \run{E} := (\lambda M. \run{(E[M])})^\dag, \]
%
where $\dag$ is the Kleisli extension operator for the monad $\Tsb$.

Then we have the following lemma, which says that running commutes with filling
an evaluation context:
%
\begin{lemma}[Run is a Homomorphism]\label{lem:run-fill}
  $\run{(E[M])} = (\run{E})(\run{M})$.
\end{lemma}
\begin{proof}
    We want to show that $\run{(E[M])} = (\run{E})(\run{M})$. We proceed by
    guarded recursion and cases on whether $M$ is a value, or $M$ steps.

    If $M$ is a value $V$, then the result is straightforward. 
    % by definition of $\run$ on terms, we have $\run{V} = \eta\, V$.

    Now suppose $M \stepsin{k} \hat{N}$, where $\hat{N} : \Sigmahat(\term A)$.
    Then we have, by definition of $\run$:

    \begin{equation} 
        \run{M} = \delta_{\sty} (\run^{\dag} {(f(\hat{N}))}).
    \end{equation}

    We also have $E[M] \stepsin{k} \Sigmahat(E[\cdot])(\hat{N})$. Thus, by the
    definition of $\run$, we have:
    \begin{equation} 
        \run{(E[M])} = \delta_{\sty} (\run^{\dag} {(f (\Sigmahat(E[\cdot])(\hat{N})) )}).
    \end{equation}

    We continue by cases on $\hat{N}$. If $\hat{N} = \cpure{M}$, then we have
    \begin{align*}
        \run{E[M]}
        &= \delta_{\sty} [\run^{\dag} {(f (\Sigmahat(E[\cdot])(\cpure{M})) )}] \\
        &= \delta_{\sty} [\run^{\dag} {(f (\cpure(E[M])))}] \\
        &= \delta_{\sty} [\run^{\dag} {\eta (E[M])}] \\
        &= \delta_{\sty} [\run E[M]]
    \end{align*}

    On the other hand, we have
    \begin{align*}
        (\run E)(\run M) 
        &= (\lambda V. \run{(E[V])})^{\dag} (\delta_{\sty} (\run^{\dag} {(f(\cpure{M}))}) ) \\
        &= \delta_{\sty} [(\lambda V. \run{(E[V])})^{\dag} (\run^{\dag} {(f(\cpure{M}))})] \\
        &= \delta_{\sty} [(\lambda V. \run{(E[V])})^{\dag} (\run^{\dag} {\eta M})] \\
        &= \delta_{\sty} [(\lambda V. \run{(E[V])})^{\dag} (\run M)] \\
        &= \delta_{\sty} [(\run E) (\run M)]
    \end{align*}
    %
    Since both sides are guarded by a $\delta_{\sty}$, it suffices to show that
    the insides are equal later. This follows from our \lob~induction
    hypothesis.

    Now suppose $\hat{N} = \ceff{(\oper, u)}$. Then for the LHS, we have
    %
    \begin{align*}
        \run{(E[M])} 
        &= \delta_{\sty} [\run^{\dag} {(f (\Sigmahat(E[\cdot])(\ceff{(\oper, u)})) )}] \\
        &= \delta_{\sty} [\run^{\dag} {(f (\ceff{(\oper, E[\cdot] \circ u)}) )}] \\
        &= \delta_{\sty} [\run^{\dag} {(\Tsb(\term{A}).\oper(\eta \circ E[\cdot] \circ u) )}] \\
        &= \delta_{\sty} [\Tsb(\term{A}).\oper(\run^{\dag} \circ \eta \circ E[\cdot] \circ u)] \\
        &= \delta_{\sty} [\Tsb(\term{A}).\oper(\run \circ E[\cdot] \circ u)] \\
        &= \delta_{\sty} [\Tsb(\term{A}).\oper(\lambda w. \run (E[u(w)]))] \\
    \end{align*}

    For the RHS, we have
    %
    \begin{align*}
        (\run E)(\run M) 
        &= (\lambda V. \run{(E[V])})^{\dag} (\delta_{\sty} (\run^{\dag} {(f(\ceff{(\oper, u)}))}) ) \\
        &= \delta_{\sty} [(\lambda V. \run{(E[V])})^{\dag} (\run^{\dag} {(f(\ceff{(\oper, u)}))})] \\
        &= \delta_{\sty} [(\lambda V. \run{(E[V])})^{\dag} (\run^{\dag} { \Tsb(\term{A}).\oper(\eta \circ u) })] \\
        &= \delta_{\sty} [(\lambda V. \run{(E[V])})^{\dag} (\Tsb(\term{A}).\oper(\run^{\dag} \circ \eta \circ u))] \\
        &= \delta_{\sty} [(\lambda V. \run{(E[V])})^{\dag} (\Tsb(\term{A}).\oper(\run \circ u))] \\
        &= \delta_{\sty} [\Tsb(\term{A}).\oper((\lambda V. \run{(E[V])})^{\dag} \circ \run \circ u)] \\
        &= \delta_{\sty} [\Tsb(\term{A}).\oper(\lambda w. (\run E) (\run (u(w))))] \\
    \end{align*}
    %
    As before, since both sides are guarded by a $\delta_{\sty}$, it suffices to
    show that the insides are equal later, which follows from our \lob~induction
    hypothesis.

    % Thus, it will suffice to show that
    % \[ 
    %   \later[ (\run \circ E[\cdot])^{\dag} (\iota(f(\hat{N}))) 
    %     = ((\lambda V. \run{(E[V])})^{\dag} \circ \run)^{\dag} {(\iota(f(\hat{N})))} ]. 
    % \]


\end{proof}

\section{Omitted Details and Proofs for Section \ref{sec:denotational-semantics}}

\subsection{Action of the Signature Endofunctor on Relations}

We describe the action of the endofunctor for the algerbaic signature $\Sigma$
on predomain relations.

Let $c : A \rel A'$ be a predomain relation. The action of the endofunctor
$\Sigma$ on $c$, written $\Sigma(c)$, is a predomain relation between
$\Sigma(A)$ and $\Sigma(A')$, defined as follows:
%
First, note that any predomain relation $c : A \rel A'$ extends to a predomain
relation $\tau \to c$ between $\tau \to A$ and $\tau \to A'$ defined pointwise,
i.e., $f \binrel{(\tau \to c)} f'$ iff for all $s : \tau$ we have $(f s)
\binrel{c} (f' s)$.
%
Similarly, given $c : A \rel A'$ we define $\later c$ such that $\tilde{x}
\binrel{(\later c)} \tilde{y}$ iff after one step passes, the left and right
sides (now in $A$ and $A'$ respectively) are related in $c$.
%
We now define $(\oper, (u, g)) \binrel{(\Sigma(c))} (\oper', (u', g'))$ iff
$\oper = \oper'$ and $u \binrel{(\tau_\oper \to c)} u'$ and 
$g \binrel{(\gamma_\oper \to \later c)} g'$.

\begin{comment}
\subsection{Generalizing Error Domains}

In this section, we discuss the guarded denotational model for gradual typing.

We begin by reviewing the key ingredients in the denotational semantics of
Giovannini et al. \cite{gtt-sgdt}. For further discussion, see their paper. For
the sake of brevity, we will refer to their model as the \emph{prior guarded
model} of gradual typing.

The prior model is based on call-by-push-value \cite{levy99}, a categorical
model that refines Moggi's monadic semantics \cite{MOGGI199155}. In
call-by-push-value, there is a category of value types $\calV$ and pure
morphisms $f : A \to A'$, and a category $\calE$ of computation types $\phi : B
\multimap B'$. A pair of adjoint functors $F \vdash U$ mediates between the two
categories, where $F : \calV \to \calE$ takes $A$ to the type $F A$ of
computations that return values in $A$, and $U : \calE \to \calV$ takes $B$ to
the type $U B$ of suspended computations of type $B$.

In the prior model, the value types are modeled by sets equipped with a partial
ordering relation that models type precision on closed values. However, since
step-insensitive reasoning does not play well with transitivity, each type also
has an additional relation called \textbf{weak bisimilarity}. They call such a
set a \textbf{predomain}. Morphisms of predomains are functions on the
underlying sets that respect both of the relations.
%
\begin{ifextended}
The precise definitions are given below:
%
\begin{definition}
A \textbf{predomain} $A$ consists of a set $A$ along with two relations:
\begin{itemize}
    \item A partial order $\ltdyn_A$.
    \item A reflexive, symmetric ``bisimilarity'' relation $\bisim_A$.
\end{itemize}

A \textbf{morphism of predomains} $f : A \to A'$ is a function between the
underlying sets that is both \emph{monotone} (if $x \ltdyn_A x'$, then $f(x)
\ltdyn_{A'} f(x')$), and \emph{bisimilarity-preserving} (if $x \bisim_A x'$,
then $f(x) \bisim_{A'} f(x')$).
\end{definition}
\end{ifextended}
%
Computation types are equipped with the same two relations, in addition to a
special distinguished error element that is the least element in the ordering,
as well as a distinguished morphism of predomains $\theta_B : \later B \to B$.
%
The precise definitions are given below:
%
\begin{definition}
An \textbf{error domain} $B$ consists of a predomain $B$ along with the
following data:
\begin{itemize}
    \item A distinguished ``error'' element $\mho_B \in B$ such that $\mho
    \ltdyn_B x$ for all $x$
    \item A morphism of predomains $\theta_B \colon \later B \to B$
    \item For all $x, y : B$ with $x \bisim_B y$, we have $\theta_B(\nxt\, x)
    \bisim_B y$
\end{itemize}
\end{definition}

A \textbf{morphism of error domains} is a morphism of the underlying predomains
that preserves the error element and the $\theta$ map.
%
For an error domain $B$, we define the predomain morphism $\delta_B := \theta_B
\circ \nxt$.

In the prior model, type precision is modeled as a kind of heterogeneous
relation on the denoted predomains. More concretely, A \textbf{predomain
relation} $c : A \rel A'$ is a relation between $A$ and $A'$ that is upward
closed in the ordering on $A'$ and downward closed in the ordering on $A$
(no condition is required involving the bisimilarity relations).
%
\begin{definition}
  A \textbf{predomain relation} $c : A \rel A'$ is a relation $c$ between the
  underlying sets of $A$ and $A'$ satisfying:
  \begin{enumerate}
    \item (upward closure): For all $x \in A$ and $y, y' \in A'$, if $x
    \binrel{c} y$ and $y \ltdyn_{A'} y'$, then $x \binrel{c} y'$.
    \item (downward closure): For all $x, x' \in A$ and $y \in A'$, if $x'
    \ltdyn_{A} x$ and $x \binrel{c} y$, then $x' \binrel{c} y$.
  \end{enumerate}
\end{definition}
%
The ``reflexive'' predomain relation is written $r(A)$ and is given by the
ordering $\ltdyn_A$. The relational composition of two predomain relations is
also a predomain relation, and $r(A)$ is the identity for this composition.

Error domain relations are predomain relations that are additionally a
congruence with respect to the operations $\theta$ and $\mho$.
%
\begin{definition}
  An \textbf{error domain relation} $d : B \rel B'$ is a relation $d$ on the
  underlying predomains satisfying
  \begin{enumerate}
     \item ($\mho$ congruence): $\mho_B \binrel{d} \mho_{B'}$.
     \item ($\theta$ congruence): For all $\tilde{x}$ in $\laterhs B$ and
     $\tilde{y} \in \laterhs B'$, if $\laterhs_t (\tilde{x}_t \binrel{d}
     \tilde{y}_t)$, then $(\theta_B (\tilde{x}) \binrel{d} \theta_{B'}
     (\tilde{y}))$. 
  \end{enumerate}
\end{definition}
Note that because $\mho_B$ is a least element this implies the
stronger property $\mho_B \binrel{d} y$ for any term $y: B'$.
%
There is again a reflexive error domain relation $r(B)$ given by the ordering.
The relational composition of two error domain relations in general is not an
error domain relation ($\theta$ congruence in particular is not preserved), but
we instead use an alternative ``free composition'' $dd' : B \rel B''$ that is
inductively generated from the relational composition.

Lastly, term precision is modeled via the notion of \emph{square}, which
captures the logical-relations style relationship between the functions denoted
by two terms.
%
More concretely:
%
\begin{definition}
Let $f : A_i \to A_o$ and $g : A_i' \to A_o'$ be predomain morphisms, and let
$c_i : A_i \rel A_i'$ and $c_o : A_o \rel A_o'$ be predomain relations. There is
a \textbf{predomain square} $f \ltsq{c_i}{c_o} g$ if for all $x \binrel{c_i} y$,
we have $f(x) \binrel{c_o} g(y)$. 

Likewise, an \textbf{error domain square} $\phi \ltsq{d_i}{d_o} \psi$ is defined
similarly, with no extra interaction with the algebraic structure.
\end{definition}
%
\end{comment}

\begin{comment}
\subsection{Perturbations}

\begin{definition}[Value and Computation Objects and Morphisms]
    A \textbf{value object} consists of a predomain $A$, a monoid $M_A$ and a
    homomorphism of monoids $i_A \colon M_A \to \morbisimid{A} := \{ f : A \to A
    \mid f \bisim \id_A \}$.

    A \textbf{computation object} consists of an error domain $B$, a monoid
    $M_B$ and a homomorphism of monoids $i_B \colon M_B \to \morbisimid{B}$.

    We define a \textbf{value morphism} to be a morphism of the underlying
    predomains, and a \textbf{computation morphism} to be a morphism of the
    underlying error domains.
\end{definition}

\begin{definition}[Push-Pull Structure]
    Let $(A, M_A, i_A)$ and $(A', M_{A'}, i_{A'})$ be value objects, and let $c
    : A \rel A'$ be a predomain relation. A \textbf{push-pull structure} for $c$
    (with respect to $M_A$ and $M_{A'}$) consists of monoid homomorphisms $\push
    : M_A \to M_A'$ and $\pull : M_A' \to M_A$ such that for all $m_A \in M_A$
    there is a square $i_A(m_A) \ltsq{c}{c} i_{A'}(\push\, m_A)$ and for all
    $m_{A'} \in M_{A'}$ there is a square $i_A(\pull\, m_{A'}) \ltsq{c}{c}
    i_{A'}(m_{A'})$. A push-pull structure for error domain relations is defined
    similarly.
\end{definition}

\begin{definition}[Value and Computation Relations]
    A \textbf{value relation} between value objects $(A, M_A, i_A)$ and $(A',
    M_{A'}, i_{A'})$ consists of a predomain relation $c : A \rel A'$ that has a
    push-pull structure, is quasi-left-representable by a morphism $e_c : A \to
    A'$, and is such that $\li c$ is quasi-right-representable by an error
    domain morphism $p_c : \li A' \to \li A$.

    A \textbf{computation relation} between computation objects $(B, M_B, i_B)$
    and $(B', M_{B'}, i_{B'})$ consists of an error domain relation $d$ that has
    a push-pull structure, is quasi-right-representable by a morphism $p_d : B'
    \to B$, and is such that $Ud$ is quasi-left-representable by a morphism $e_d
    : UB \to UB'$
\end{definition}
\end{comment}

\subsection{Adjusting the Term Semantics for Casts}
We provide more details on the change we make to the denotational semantics so
that the denotational and operational semantics are in lock-step up to the
presence of bureaucratic delays in the operational semantics. For an overview,
see Section \ref{sec:perturbations} in the paper.

We begin by recalling that term precision $M \ltdyn N : A$ is modeled in the
semantics of Giovannini et al. by $\sem{M} \ltbisim \sem{N}$ where
%
\[ f \ltbisim g := \exists f', g'. f \bisim f' \ltdyn g' \bisim g. \]
%
That is, we do not require that $f$ and $g$ are in lock-step; we are allowed to
use any $f'$ bisimilar to $f$ and $g'$ bisimilar to $g$ such that $f'$ is
lock-step related to $g'$ in the error ordering.
%
This is useful because the downcast associated to the composition of type
precision $c : A \ltdyn A'$ and $c' : A' \ltdyn A''$ is a morphism from
$F(\sem{A''})$ to $F(\sem{A})$, but in the prior model it is not simply the
composition of the downcast morphisms for $c$ and $c'$. Instead, an additional
perturbation is inserted in order to make the relational squares line up in the
proof of quasi-representability. More details can be found in the Appendix of
their paper \cite{gtt-sgdt-extended-version}.

On the other hand, observe that we could instead define the downcast morphism
for the composition $cc'$ to be the composition of the corresponding downcasts,
without the additional perturbation, and then note that the resulting morphism
is \emph{weakly bisimilar} to the version with the additional perturbation. This
is sufficient for graduality, since the additional perturbation will be absorbed
by the weak bisimilarity in the definition of $f \ltbisim g$.

% Using the above definition for term precision, we see that the additional
% perturbations in the cast corresponding to a composition will be absorbed by the
% weak bisimilarity.

With this observation in mind, we augment the denotation of type precision in
the prior denotational model by associating to each value relation $c$ in the
semantics two morphisms: a ``relational'' embedding morphism $e$ such that $c$
is quasi-left-representable by $e$, and a ``term'' embedding morphism $e' \bisim
e$. We similarly associate two projection morphisms to each computation relation
$d$. We use the term versions of the embedding and projection morphisms as our
denotations of upcasts and downcasts, and we use the relational versions
internally in the proof of graduality.
%
The result is that our term semantics for casts no longer involves unnecessary
perturbations that were inserted to make the graduality proof work.



\section{Omitted Details and Proofs for Section \ref{sec:relating-the-semantics}}

\subsection{Relational Lifting that Ignores Bureaucratic Steps}

  \begin{mathpar}
    \inferrule*[right=RetRet]
    { x \binrel{R} y }
    { \eta^{T_{\mho s}}(x) \binrel{(\bbt(R))} \eta^{T_{\mho sb}}(y) }

    \inferrule*[right=OpOp]
    { \oper \in \calA_{\mho s}.\Op 
      \and \forall s.\, (u\, s) \binrel{(\bbt{R})} (u'\, s) \\
      \and \forall s.\, (g\, s) \binrel{(\later(\bbt{R}))} (g'\, s)  
    }
    { (T_{\mho s}(X).\oper(u, g)) \binrel{(\bbt{R})} (T_{\mho sb}(Y).\oper(u', g')) }

    \inferrule*[right=IdStep]
    { \bar{x} \binrel{(\bbt(R))} \bar{y} }
    { \bar{x} \binrel{(\bbt(R))} \delta_b(\bar{y}) }

  \end{mathpar}

\subsection{Properties of the Logical Relation}

First, we prove two anti-reduction lemmas, one involving bureaucratic steps and
the other involving true steps.

%Proof of the anti-reduction lemma (Lemma \ref{lem:anti-reduction}):
\begin{lemma}[Anti-Reduction]\label{lem:anti-reduction}

  Let $\cdot \vdash M : A$, and suppose $M \stepsin{0} \cpure{M'}$. If
  $\erel{A}{\xbar}{M'}$, then $\erel{A}{\xbar}{M}$.

\end{lemma}
\begin{proof}
    Let $\cdot \vdash M : A$, where $M \stepsin{0} \cpure{M'}$, and suppose
    $\erel{A}{\xbar}{M'}$. We will show that $\erel{A}{\xbar}{M}$.

    % Let $\sty$ be $\bstep$ or $\tstep$ according to whether $k$ is $0$ or $1$,
    % respectively.

    By the definition of $\run$, we have that
    \begin{align*}
        \run{M} &= \delta_{\bstep} (\run^{\dag} {(f(\cpure{M'}))}) \\
                &= \delta_{\bstep} (\run^{\dag} {\eta(M')}) \\
                &= \delta_{\bstep} (\run{M'}).
    \end{align*}

    Since $\erel{A}{\xbar}{M'}$, we have that $\liftvrel{A}{\xbar}{(\run M')}$.
    By definition of the relational lifting, we have that
    $\liftvrel{A}{\xbar}{\delta_{\bstep}(\run{M'})}$. Thus, we have
    $\liftvrel{A}{\xbar}{\run{M}}$, which means that $\erel{A}{\xbar}{M}$.

\end{proof}

\begin{lemma}[Anti-Reduction for True Steps]\label{lem:anti-reduction-true}
    Let $\cdot \vdash M : A$, and suppose $M \stepsin{1} \cpure{M'}$. Then for
    all $\tilde{x} : \later \Ts(\sem{A})$, if $\later_t
    (\erel{A}{\tilde{x}_t}{M'})$, then $\erel{A}{\theta(\tilde{x})}{M}$.
\end{lemma}
\begin{proof}
    Suppose $\later_t (\erel{A}{\tilde{x}_t}{M'})$. We will show that
    $\erel{A}{\theta(\tilde{x})}{M}$. By the definition of $\run$, we have that
    \begin{align*}
        \run{M} &= \delta_{\tstep} (\run^{\dag} {(f(\cpure{M'}))}) \\
                &= \delta_{\tstep} (\run^{\dag} {\eta(M')}) \\
                &= \delta_{\tstep} (\run{M'}).
    \end{align*}

    It suffices to show $\liftvrel{A}{\theta(\tilde{x})}{\run{M}}$,
    i.e., $\liftvrel{A}{\theta(\tilde{x})}{\delta_{\tstep}(\run{M'})}$.
    Now using the definition of the relational lifting, it suffices to show that
    \[ \later_t(\liftvrel{A}{\tilde{x}_t}{(\nxt (\run{M'}))_t}), \]
    i.e.,
    \[ \later_t(\liftvrel{A}{\tilde{x}_t}{(\run{M'})}). \]
    But this is by definition the same as
    \[ \later_t(\erel{A}{\tilde{x}_t}{M'}), \]
    which is our assumption.
\end{proof}


Next, we have a semantic bind lemma for the logical relation, which allows us to
replace arbitrary terms in an evaluation context with values.

%Proof of the Semantic Bind Lemma:
\begin{lemma}[Monadic Bind]\label{lem:monadic-bind}
    % Suppose for all values $V, W$ that $\vrel{A}{V}{W}$ implies
  % $\erel{B}{\phi(V)}{(E[W])}$. Then, for all $\erel{A}{\bar{x}}{N}$, we have
  % $\erel{B}{(\phi(\bar{x}))}{(E[N])}$.
  Let $\phi : T_{\mho s}(\sem{A}) \to T_{\mho s}(\sem{B})$, and let $\bullet : A \vdash E : B$.
  %
  Suppose for all values $x, V$ that $\vrel{A}{x}{V}$ implies $\erel{B}{\phi(T_{\mho s}.\eta(x))}{(E[V])}$.
  %
  Then for all $\bar{x} \in T_{\mho s}(\sem{A})$ and terms $N$ of type $A$ such that
  $\erel{A}{\bar{x}}{N}$, we have $\erel{B}{(\phi(\bar{x}))}{(E[N])}$
\end{lemma}
\begin{proof}
    Follows from Lemma \ref{lem:run-fill} and the universal property of the
    relational lifting $\bbt(R)$.
\end{proof}

\begin{comment}
\begin{proof}
    Suppose for all values $x, V$ that $\vrel{A}{x}{V}$ implies
    $\erel{B}{\phi(T_{\mho s}.\eta(x))}{(E[V])}$. Further, suppose that
    $\erel{A}{\bar{x}}{N}$.
    %
    We will show that $\erel{B}{(\phi(\bar{x}))}{(E[N])}$, which by definition
    is $\liftvrel{B}{\phi(\bar{x})}{(\run (E[N]))}$

    The proof is by guarded recursion and induction over the proof that
    $\erel{A}{\bar{x}}{N}$, since this is defined as
    $\liftvrel{A}{\bar{x}}{(\run N)}$. We have the following cases:

    \begin{itemize}
        \item Case \textsc{RetRet}: This means that $\bar{x} = \eta\, x$ for
        some $x \in \sem{A}$ and $\run N = \eta'\, V$ for some $V \in \val{A}$
        where $\vrel{A}{x}{V}$. (Here, $\eta$ is the unit for $T_{\mho s}$ and
        $\eta'$ is the unit for $T_{\mho sb}$.) Then by our assumption, we have
        $\erel{B}{\phi(x)}{(E[V])}$, as required.

        \item Case \textsc{OpOp}: In this case, we have
        %
        $\bar{x} = T_{\mho s}(\sem{A}).\oper(u, g)$ for some $u$ and $g$.
        %
        We also know that
        %
        $\run N = (T_{\mho sb}(\val A).\oper(u', g'))$
        %
        for some $u'$ and $g'$, where for all $s$, we have $\liftvrel{A}{(u\,
        s)}{(u'\, s)}$ and $(g\, s) \binrel{(\later(\bbt{\vrel{A}}))} (g'\, s)$.

        Now since $\phi$ is a homomorphism, we have
        %
        \[ \phi(T_{\mho s}(\sem{A}).\oper(u, g)) = 
           T_{\mho s}(\sem{B}).\oper(\phi \circ u, (\later \phi) \circ g).
        \]
        %
        Furthermore, by Lemma \ref{lem:run-fill}, we have $\run (E[N]) = (\run
        E)(\run N)$ where we recall that $\run E$ is defined as $(\lambda V.
        \run{(E[V])})^\dag$. In particular, this means that
        \begin{align*}
            \run (E[N]) &= (\lambda V.\run{(E[V])})^\dag (\run N)  \\
                        &= (\lambda V.\run{(E[V])})^\dag (T_{\mho sb}(\val A).\oper(u', g')) \\
                        &= T_{\mho sb}(\val B).\oper((\run E) \circ u', (\later (\run E)) \circ g')
        \end{align*}

        Thus, using rule $\textsc{OpOp}$ and the induction hypotheses, this case is complete.
        % TODO: more detail?

        
        \item Case \textsc{IdStep}: In this case, we have
        %
        $\run N = \delta_b(\bar{V})$
        %
        for some $\bar{V} \in T_{\mho sb}(\val{A})$, and
        $\liftrel{A}{\bar{x}}{\bar{V}}$.
                
        Then as above, we have
        %
        \begin{align*}
            \run (E[N]) &= (\run E)(\run N) \\
                &= T_{\mho sb}(\val B).\delta_b[(\run E)(\bar{V})] \\
                &= 
        \end{align*} 
        %
        so by rule \textsc{IdStep}, it suffices to show
        %
        \[ \liftrel{B}{\phi(\bar{x})}{(\run E)(\bar{V})} \]
    \end{itemize}
\end{proof}
\end{comment}

\subsection{Proof of the Fundamental Lemma for the Logical Relation}

Towards proving the fundamental lemma, we prove a series of compatibility
lemmas, one for each term constructor in the cast calculus.

\begin{lemma}[Logical Relation for Natural Number Values]\label{lem:lr-nat}
    Let $n \in \mathbb{N}$, and let $V \in \val{\nat}$.
    Then if $\vrel{\nat}{n}{V}$, then $\vrel{\nat}{(\suc\, n)}{(\suc\, V)}$
\end{lemma}
\begin{proof}
    If $\vrel{\nat}{n}{V}$, then by definition we have $V = \suc^n(\zro)$. This
    means that $\suc\,V = \suc^{n+1}(\zro)$, i.e., $\vrel{\nat}{(\suc\,
    n)}{(\suc\, V)}$.
\end{proof}

\begin{lemma}[Compatibility for Zero]
    We have $\simrel{\Gamma}{\nat}{\sem{\zro}}{\zro}$.
\end{lemma}
\begin{proof}
    Follows by definition of $\vrel{\nat}{}{}$.
\end{proof}


\begin{lemma}[Compatibility for Successor]
    If $\simrel{\Gamma}{\nat}{\sem{M}}{M}$, then
    $\simrel{\Gamma}{\nat}{\sem{\suc\, M}}{\suc\, M}$.
\end{lemma}
\begin{proof}
    Let $(\grel{\Gamma} {\gambar} {\gamma'})$. We need to show 
    %
    \[ \erel{\nat} {(\sem{\suc\, M}\, \gambar)} {\suc\, M[\gamma']}. \]
    %
    We have $\sem{\suc\, M}\, \gambar = (\lambda x.\eta(\suc\, x))^\dag(\sem{M}\, \gambar)$
    
    By the bind lemma, it suffices to show that (1) $\erel{\nat}{\sem{M}\,
    \gambar}{M[\gamma']}$ and (2) for all $x, V$ such that $\vrel{\nat}{x}{V}$,
    we have $\erel{\nat}{\eta(\suc\,x)}{\suc\,V}$.
    Note that (1) follows from our assumption.

    To show (2), we note that both are values, so it suffices to show that
    $\vrel{\nat}{\suc\, x}{\suc\, V}$. Then by Lemma \ref{lem:lr-nat}, it
    suffices to show $\vrel{\nat}{x}{V}$, which is what we have assumed.    
\end{proof}

\begin{lemma}[Compatibility for Variables]
    For all $\Gamma_1, \Gamma_2$, we have $\simrel{\Gamma_1, x : A, \Gamma_2}{A}{\sem{x}}{x}$.
\end{lemma}
\begin{proof}
    Let $\grel{\Gamma_1, x : A, \Gamma_2}{\gambar}{\gamma'}$. We need to show
    \[ \erel{A}{\sem{x}\, \gambar}{x[\gamma']}.  \]
    We have $\sem{x}\, \gambar = \eta(\gambar(x))$, and $x[\gamma'] = \gamma'(x)$.
    Since both are values, it suffices to show
    \[ \vrel{A}{\gambar(x)}{\gamma'(x)}, \]
    which follows from the assumption that $\gambar$ is related to $\gamma'$.
\end{proof}

\begin{lemma}[Compatibility for Lambda]
    Let $\fbar : \sem{\Gamma} \times \sem{A_i} \to T\sem{A_o}$.
    %
    Suppose $\simrel{\Gamma, x : A_i}{A_o}{\fbar}{M}$.
    %
    Then we have $\simrel{\Gamma}{A_i \ra A_o}{\bar{m}}{\lambda x.M}$, where
    $\bar{m}(\gambar) = \eta(\lambda x. \fbar(\gambar, x))$.
\end{lemma}
\begin{proof}
    % \sem{M} : \sem{\Gamma, x : A_i} \to T(\sem{A_o})
    % = \sem{\Gamma} \times \sem{A_i} \to T(\sem{A_o})
    %
    % \sem{\lambda x.M} : \sem{\Gamma} \to T(\sem{A_i \ra A_o})
    % \sem{\lambda x.M} \gambar = \eta( \lambda x. \sem{M}(\gambar , x) )

    Let $\grel{\Gamma}{\gambar}{\gamma'}$. We need to show
    %
    \[ \erel{A_i \ra A_o} {(\bar{m}\, \gambar)} {(\lambda x.M)[\gamma']}. \]
    %
    Since both are values, it suffices to show that
    %
    \[ \vrel {A_i \ra A_o} {(\lambda x. \fbar(\gambar, x))} {(\lambda x.M[\gamma'])}. \]
    %
    Now, by definition of $\vrel{A_i \ra A_o}{}{}$, let $x : \sem{A_i}$ and $V : \val{A_i}$,
    and suppose $\vrel{A_i}{x}{V}$. We will show that
    \[ \erel{A_o} {((\lambda x. \fbar(\gambar, x))\, x)} {((\lambda x.M[\gamma'])\, V)}, \]
    i.e.,
    \[ \erel{A_o} {(\fbar(\gambar, x))} {((\lambda x.M[\gamma'])\, V)}. \]
    Since $((\lambda x.M[\gamma'])\, V) \stepsin{0} (M[\gamma'][V/x])$,
    by anti-reduction it suffices to show
    \[ \erel{A_o} {(\fbar(\gambar, x))} {(M[\gamma'][V/x])}. \]

    Let $\bar{\gamma_1} := (\gambar, x) \in \sem{\Gamma, x : A_i}$.
    Let $\gamma_1' := (\gamma', x \mapsto V)$.
    Then we have $\grel{\Gamma, x : A_i}{\bar{\gamma_1}}{\gamma'}$,
    so by our assumption that $\simrel{\Gamma, x : A_i}{A_o}{\fbar}{M}$,
    we have $\erel{A_o} {\fbar\, \bar{\gamma_1}} {M[\gamma_1']}$.
    This is what we needed to show.

\end{proof}

\begin{lemma}[Compatibility for Function Application]
    Let $\fbar : \sem{\Gamma} \to T\sem{A_i \ra A_o}$ and 
    $\xbar : \sem{\Gamma} \to T\sem{A_i}$.
    %
    Suppose $\simrel{\Gamma}{A_i \ra A_o}{\fbar}{M}$, and
    $\simrel{\Gamma}{A_i}{\xbar}{N}$. Then we have
    $\simrel{\Gamma}{A_o}{\bar{m}}{M\, N}$ where

    $\bar{m}\, \gambar := (\lambda g. g^\dag (\xbar\, \gambar))^\dag (\fbar\, \gambar)$
\end{lemma}
\begin{proof}
    Let $\grel{\Gamma}{\gambar}{\gamma'}$. We need to show
    %
    \[ \erel{A_o} {\bar{m}\, \gambar} {(M[\gamma']\, N[\gamma'])}. \]

    By the Semantic Bind Lemma with
    %
    $\phi(\bar{x_f}) = (\lambda g. g^\dag (\xbar\, \gambar))^\dag (\bar{x_f})
    %
    : \Ts(\sem{A_i \ra A_o}) \to \Ts(A_o)$, and 
    %
    $E = \bullet\, N[\gamma']$, it suffices to show:
    %
    (1) $\erel{A_i \ra A_o}{(\fbar\, \gambar)}{M[\gamma']}$ and
    %
    (2) for all $x_f : \sem{A_i \ra A_o}$ and $\Gamma \vdash V_f : A_i \ra A_o$,
    if $\vrel{A_i \ra A_o}{x_f}{V_f}$, then $\erel{A_o} {(\phi(\eta\, x_f))} {V_f\, (N[\gamma'])}$.

    (1) follows from our assumption that $\simrel{\Gamma}{A_i \ra
    A_o}{\fbar}{M}$.
    %
    To show (2), we again apply the Semantic Bind Lemma. This time, it suffices
    to show that for all $x : \sem{A_i}$ and $\Gamma \vdash V : A_i$, if
    $\vrel{A_i}{x}{V}$, then $\erel{A_o}{x_f\, x}{V_f\, V}$. This follows from
    our assumption that $\vrel{A_i \ra A_o}{x_f}{V_f}$.

    % Note: \phi(\eta\, x_f) = (\lambda f. f^\dag (\sem{N} \gambar))^\dag (\eta\, x_f)
    %                        = x_f^\dag (\sem{N} \gambar) 


\end{proof}

\begin{lemma}[Compatibility for Product Split]
Suppose $\simrel{\Gamma}{A_1 \times A_2}{\sem{M}}{M}$ and 
$\simrel{(\Gamma, x : A_1, y : A_2)}{A}{\sem{N}}{N}$. Then we have
$\simrel {\Gamma} {A} {\sem {M_{\text{split}}}} {{M_{\text{split}}}}$,
where ${M_{\text{split}}} = \textrm{split } (x,y) = M \textrm{ in } N$.
\end{lemma}
\begin{proof}
    Follows by the semantic bind lemma and the given hypotheses.
\end{proof}

\begin{lemma}[Compatibility for Pairing]
Suppose $\simrel{\Gamma}{A_1}{\sem{M_1}}{M_1}$ and $\simrel{\Gamma}{A_2}{\sem{M_2}}{M_2}$.
Then we have $\simrel{\Gamma}{A_1 \times A_2}{\sem{(M_1, M_2)}}{(M_1, M_2)}$.
\end{lemma}
\begin{proof}
    Follows by the semantic bind lemma and the given hypotheses.
\end{proof}


\begin{lemma}[Compatibility for Effects]
    Let $\oper$ be an effect with arity $n$. Suppose for each $i \in \{1, \dots,
    n\}$ we have $\simrel{\Gamma}{A}{\sem{M_i}}{M_i}$. Then we have
    $\simrel{\Gamma}{A}{\sem{\synoper{M}{n}}}{(\synoper{M}{n})}$.
\end{lemma}
\begin{proof}
    Let $(\grel{\Gamma}{\gambar}{\gamma'})$. We need to show that
    %
    \[ \erel{A} {(\sem{\synoper{M}{n}}\, \gambar)} {(\synoper{M}{n})[\gamma']}. \]
    %
    First note that
    \[ \sem{\synoper{M}{n}}\, \gambar = T_{\mho s}.\oper(\lambda i.\, (\sem{M_i}\, \gambar)). \]
    %
    Furthermore, note that $\synopernoidx{M_i[\gamma']}{n}$ steps to
    $\hat{N} := \ceff{(\oper, (\lambda\, (i : [n]).\, M_i[\gamma']))}$, so by definition of
    $\run$, we have
    %
    \begin{align*}
        \run{\synopernoidx{M_i[\gamma']}{n}}
        &= \delta_b (\run^{\dag} {(f(\hat{N}))}) \\
        &= \delta_b (\run^{\dag} \Tsb(\term{A}).\oper(\eta \circ (\lambda\, i.\, M_i[\gamma']))) \\
        &= \delta_b (\Tsb(\term{A}).\oper(\run^{\dag} \circ \eta \circ (\lambda\, i.\, M_i[\gamma']))) \\
        &= \delta_b (\Tsb(\term{A}).\oper(\run \circ (\lambda\, i.\, M_i[\gamma']))) \\
        &= \delta_b (\Tsb.\oper(\lambda i.\, \run(M_i[\gamma']))). 
    \end{align*}

    Now, using rules \textsc{IdStep} and \textsc{OpOp} of the relational
    lifting, it suffices to show that for each $i \in [n]$, we have
    %
    \[ \liftvrel{A}{(\sem{M_i}\, \gambar)}{(\run (M_i[\gamma']))}, \]
    %
    which follows from our assumption.


\end{proof}

\begin{lemma}[Compatibility for Casts]
    Let $c : A \ltdyn A'$, and let $\fbar : \sem{\Gamma} \to T\sem{A}$.
   
    If $\simrel{\Gamma}{A}{\fbar}{M}$, then we have
    $\simrel{\Gamma}{A'}{(T(\sem{\upc c}) \circ \fbar)}{\upc{c} M}$.

    Likewise, let $\gbar : \sem{\Gamma} \to T\sem{A'}$.
    If $\simrel{\Gamma}{A'}{\gbar}{M}$, then we have
    $\simrel{\Gamma}{A}{\sem{\dnc c} \circ \gbar}{\dnc{c} M}$.
\end{lemma}
\begin{proof}
    The proof is by guarded recursion and induction on the type precision
    derivation $c$. Let $(\grel{\Gamma}{\gambar}{\gamma'})$.

    We consider the upcasts first. We need to show
    %
    \[ \erel {A'} {T\sem{\upc c} (\fbar\, \gambar)} {(\upc c M)[\gamma']}. \]
    %
    Since casts are evaluation contexts, it suffices by the semantic bind lemma
    to assume that the inputs to the casts are values.
    Specifically, let $x : \sem{A}$ and $V : \val{A}$, and suppose that $\vrel{A}{x}{V}$.
    We will show that

    \[ \erel{A'}{(\eta\, e_c(x))} {\upc c V}.  \]

    \begin{itemize}
        \item Case $c = \inat$. Then our assumption says that $\vrel{\nat}{x}{V_n}$, and we need to show that
        \[ \erel {\dyn} {(\eta\, e_{\inat} (n))} {(\upc {\inat}(V_n))}. \]
        Since both are values, it suffices to show that 
        \[ \vrel{\dyn}{e_{\inat}(n)}{\upc {\inat}(V_n)}, \]
        which follows from the definition of $\vrel{\dyn}{}{}$ and our assumption.

        \item Case $c = \itimes$. Then our assumption says that $\vrel{\dyn \times \dyn}{(x, y)}{V_\times}$,
        and we need to show that
        \[ \erel {\dyn} {(\eta\, e_{\itimes}(x, y))} {(\upc {\itimes}(V_\times))}.  \]
        Since both are values, it suffices to show that
        \[ \vrel{\dyn}{e_{\itimes}(x, y)}{\upc {\itimes}(V_\times)}, \] 
        which follows from the definition of $\vrel{\dyn}{}{}$ and our assumption.
        
        % \item Case $c = \itimes(c_\times)$ where $c_\times : A_1 \times A_2 \ltdyn \dyn \times \dyn$.
        % Then our assumption says that $\vrel{A_1 \times A_2}{(x, y)}{V_\times}$, and we need to show that
        % \[ \erel {\dyn} {(\eta, (e_{\itimes(c_\times)}(x, y)))} {(\upc {\itimes(c_\times)} V_\times)}. \]
        % Since $(\upc {\itimes(c_\times)} (V_\times)) \stepsin{0} (\tagtimes (\upc c_\times (V_\times)))$,
        % then by anti-reduction it suffices to show that
        % \[ \erel {\dyn} {(\eta, (e_{\itimes(c_\times)}(x, y)))} {(\tagtimes (\upc c_\times (V_\times)))}. \]
        
        % Note that $e_{\itimes(c_\times)} = e_{\itimes} \circ e_{c_\times}$,
        % so by the induction hypothesis for $c_\times$, we have
        % \[ \erel {\dyn \times \dyn} {e_{c_\times}(x, y)} {\upc c_\times (V_\times)}. \]
        % Then the result follows using the same reasoning as in case $\itimes$.
        % Note: the terms are closed, so when we apply the IH we don't need to worry about the substitutions.
        
        \item Case $c = \iarr$.
        Then our assumption says that $\vrel{\dyn \ra \dyn}{f}{V_f}$, and we need to show that
        \[ \vrel{\dyn}{e_{\iarr}(f)}{\upc \iarr (V_f)}. \]
        By the definition of $\vrel{\dyn}{}{}$, it suffices to show that for all $x, V$ we have
        \[ (\later(\vrel{\dyn}{x}{V}) \to \erel{\dyn}{(f\, x)}{(V_f\, V)}) \]
        But we know that this holds now by our assumption, so it also holds later.

        % \item Case $c = \iarr(c_\arr)$ where $c_\arr : A_i \ra A_o \ltdyn \dyn \ra \dyn$.
        % Then our assumption says that $\vrel{A_i \ra A_o}{f}{V_f}$, and we need to show that
        % \[ \erel {\dyn} {(\eta\, (e_{\iarr(c_\arr)}(f)))} {(\upc {\iarr(c_\arr)} V_f)}. \]
        % Since $(\upc {\iarr(c_\arr)} V_f) \stepsin{0} (\tagarr (\upc c_\arr (V_f)))$, by anti-reduction it suffices to show that
        % \[ \erel {\dyn} {(\eta\, (e_{\iarr(c_\arr)}(f)))} {(\tagarr (\upc c_\arr (V_f)))}. \]
        
        % Note that $e_{\iarr(c_\arr)} = e_\tagarr \circ e_{c_\arr}$,
        % so by the induction hypothesis for $c_\arr$, we have
        % \[ \erel {\dyn \ra \dyn} {(e_{\iarr(c_\arr)}(f))} {(\upc c_\arr (V_f))}. \]
        % Then the result follows using the same reasoning as in case $\tagarr$.


        \item Case $c = c_1 c_2$ where $c_1 : A_1 \ltdyn A_2$ and $c_2 : A_2 \ltdyn A_3$.
        Then our assumption says that $\vrel{A_1}{x}{V}$, and we need to show that
        \[ \erel{A_3} {(\eta(e_{c_1 c_2}(x)))} {(\upc {c_1 c_2} V)}. \]
        Note that $e_{c_1 c_2} = e_{c_2} \circ e_{c_1}$ and that 
        $\upc {(c_1 c_2)} V \stepsin{0} (\upc c_2 (\upc c_1 V))$.
        Thus, it suffices to show that
        \[ \erel{A_3} {(\eta(e_{c_2}(e_{c_1}(x))))} {(\upc c_2 (\upc c_1 V))}. \]
        
        Observe that $(\eta(e_{c_2}(e_{c_1}(x)))) = T(e_{c_2})(\eta(e_{c_1}(x)))$.
        Thus, by the semantic bind lemma, with $\phi(\xbar) = T(e_{c_2})(\xbar)$
        and $E = \upc {c_2} \bullet$ it suffices to show that
        \[ \erel{A_2}{\eta(e_{c_1}(x))}{\upc {c_1} V}, \]
        and that for any $\vrel{A_2}{x'}{V'}$ we have % $\erel{A_3}{\phi(\eta(x'))}{E[V']}$
        \[ \erel{A_3}{\eta(e_{c_2}(x'))}{\upc {c_2} V'}. \]
        These hold by the induction hypotheses for $c_1$ and $c_2$ respectively.


        % TODO: need to justify why the value is a pair of values?
        \item Case $c = c_1 \times c_2$ where $c_1 : A_1 \ltdyn A_1'$ and $c_2 : A_2 \ltdyn A_2'$.
        Then our assumption says that $\vrel{A_1 \times A_2}{(x_1, x_2)}{V_\times}$.
        Note that since $V_\times$ is a closed value of type $A_1 \times A_2$,
        then we have $V_\times = (V_1, V_2)$, where $V_1$ and $V_2$ are closed
        values of type $A_1$ and $A_2$ respectively. We need to show that
        \[ \erel {A_1' \times A_2'} {(\eta(e_{c_1 \times c_2}(x_1, x_2)))} {(\upc {c_1 \times c_2} (V_1, V_2))}. \]
        Now note that $e_{c_1 \times c_2}(x_1, x_2) = (e_{c_1}(x_1), e_{c_2}(x_2))$.
        Furthermore, we have
        $\upc {(c_1 \times c_2)} (V_1, V_2) \stepsin{0} (\upc {c_1} V_1, \upc {c_2} V_2)$,
        so by anti-reduction, it suffices to show
        \[ \erel {A_1' \times A_2'} {(\eta(e_{c_1}(x_1), e_{c_2}(x_2)))} {(\upc {c_1} V_1, \upc {c_2} V_2)}. \]
        Now using the semantic bind lemma, with $\phi(\xbar) = T(\lambda x. (x, e_{c_2}(x_2)))(\xbar)$ and 
        $E = (\bullet, \upc c_2 V_2)$, it suffices to show that (1)
        $\erel{A_1'}{\eta(e_{c_1}(x_1))}{(\upc {c_1} V_1)}$,
        and that (2) for all $\vrel{A_1'}{x_1'}{V_1'}$, we have 
        \[ \erel{A_1' \times A_2'}{\eta(x_1', e_{c_2}(x_2))}{(V_1', \upc c_2 V_2)}. \]
        Note that (1) follows from the inductive hypothesis for $c_1$, and to show (2), we apply the
        semantic bind lemma again, with $\phi'(\xbar) = T(\lambda y. (x_1', y))$ and $E' = (V_1', \bullet)$.
        This time, it suffices to show that $\erel{A_2'}{\eta(e_{c_2}(x_2))}{(\upc c_2 V_2)}$,
        and that for all $\vrel{A_2'}{x_2'}{V_2'}$, we have
        \[ \erel{A_1' \times A_2'}{\eta(x_1', x_2')}{(V_1', V_2')}. \]
        The former follows from the inductive hypothesis for $c_2$, and the latter follows
        from the assumptions that $\vrel{A_1'}{x_1'}{V_1'}$ and $\vrel{A_2'}{x_2'}{V_2'}$. 

        \item Case $c = c_i \ra c_o$, where $c_i : A_i \ltdyn A_i'$ and $c_o :
        A_o \ltdyn A_o'$. Follows by anti-reduction, the semantic bind lemma,
        and the induction hypotheses for $c_i$ and $c_o$.
        
        % We need to show
        % %
        % \[ \erel {A_i' \ra A_o'} 
        %     {\sem{\upc {(c_i \ra c_o)} (f\, \gambar)}} 
        %     {\upc {(c_i \ra c_o)} V_f[\gamma']}.
        % \]
        % %
        % We have 
        % %
        % \begin{align*}
        %     \sem{\upc {(c_i \ra c_o)} (f\, \gambar)} 
        %         &= ((\sem{\dnc {c_i}} \tok T\sem{A_o'}) \circ (\sem{A_i} \tok UT(\sem{\upc {c_o}}))) (f\, \gambar) \\
        %         &= \lambda x'. T\sem{\upc {c_o}}[(f\, \gambar)^\dag (\sem{\dnc {c_o} (x')})],
        % \end{align*}
        % %
        % where $\tok$ is the Kleisli action of $\to$.

        % Let 
        % %
        % $\bar{M} := \lambda x'. T\sem{\upc {c_o}}[(f\, \gambar)^\dag (\sem{\dnc {c_o} (x')})]$,
        % %
        % and let
        % %
        % $N := (\lambda x. (\upc {c_o} (V_f[\gamma']\, (\dnc {c_i} x))))$.
        % %
        % Note that
        % %
        % \[ \upc {(c_i \ra c_o)} V_f[\gamma'] \stepsin{0} N. \] 
        % %
        % Thus, by the anti-reduction lemma, it suffices to show that
        % %
        % \[ \erel {A_i' \ra A_o'} {\bar{M}} {N}. \]
        % %
        
    \end{itemize}

    Now we consider the downcasts. We again apply the semantic bind lemma. Let
    $x' : \sem{A'}$ and $V' : \val{A'}$, and suppose that $\vrel{A'}{x'}{V'}$.
    We will show that
    %
    \[ \erel{A}{p_c(\eta\, x')} {\dnc c V'}.  \]

    \begin{itemize}
        \item Case $c = \inat$. Then $x' \in D$, and $V' \in \val{\dyn}$, and our
        assumption says that $\vrel{\dyn}{x'}{V'}$. We need to show that
        \[ \erel {\nat} {p_{\inat}(\eta\, x')} {\dnc {\inat} V'}.  \]
        We proceed by cases on the value $x'$ of type $D$. If it's $\inat(n)$
        for some $n$, then we know that $V' = \upc {\inat} V_n$ for some value
        $V_n$, where $V_n = \suc^n(\zro)$.
        %
        We have that $p_{\inat}(\eta\, \inat(n)) = \eta\, n$, and that $\dnc
        \inat (\upc \inat V_n) \stepsin{0} V_n$. Thus, by anti-reduction, it
        suffices to show that
        %
        \[ \erel {\nat} {\eta\, n} {V_n}. \]
        %
        It then suffices to show $\vrel{\nat}{n}{V_n}$, which holds by the fact
        that $V_n = \suc^n(\zro)$.

        If $x'$ is a product or a function, then $V'$ will be as well. Then
        $p_{\inat}(\eta\, x') = \mho$, and $\dnc \inat V' \stepsin{0} \mho$.
        Then since $\erel{\nat}{\mho}{\mho}$ by definition of the relation, this
        completes the proof in this case.

        \item Case $c = \itimes$: this is similar to the $\inat$ case.
        
        \item Case $c = \iarr$. Then $x' \in D$ and $V' \in \val{\dyn}$, and our
        assumption says that $\vrel{\dyn}{x'}{V'}$. We need to show that
        %
        \[ \erel {\dyntodyn} {p_{\iarr}(\eta\, x')} {\dnc {\iarr} V'}. \]
        %
        We proceed by cases on the value $x'$ of type $D$. If it's a natural
        number or product, then $V'$ will be as well. Then both sides will
        error, and the proof is complete.
        %
        Otherwise, if $x' = \iarr(h)$ for some $h$, then we know that $V' = \upc
        {\iarr} V_f$ for some $V_f \in \val{\dyntodyn}$, and that $\later (\vrel{\dyntodyn}{h}{V_f})$.

        We have that $p_{\iarr}(\eta\, \iarr(h)) = \theta(\nxt(\eta\, h))$, i.e., a
        step is inserted. We also know that $\dnc {\iarr} (\upc {\iarr} V_f)
        \stepsin{1} V_f$. Thus, by anti-reduction for true steps (Lemma
        \ref{lem:anti-reduction-true}), it suffices to show that
        %
        \[ \later_t(\erel{\dyntodyn}{(\nxt(\eta\, h))_t}{V_f}), \]
        i.e.,
        \[ \later (\erel{\dyntodyn}{(\eta\, h)}{V_f}). \]
        %
        This follows from the fact that $\later (\vrel{\dyntodyn}{h}{V_f})$.

        \item Case $c = c_1 c_2$: this is similar to the corresponding upcast case.
        \item Case $c = c_i \ra c_o$: this is similar to the corresponding upcast case.
        \item Case $c = c_1 \times c_2$: this is similar to the corresponding upcast case.
    \end{itemize}


\end{proof}


Now we can prove the fundamental lemma:
\begin{lemma}[Fundamental Lemma]
  Let $\Gamma \vdash M : A$. Then $\Gamma \vDash \sem{M} \sim M : A$.
\end{lemma}
\begin{proof}
    By induction on the term $M$, using the compatibility lemmas in each case.
\end{proof}



%\section{Omitted Proofs Section \ref{sec:globalization}}

\section{Omitted Details and Proofs for Section \ref{sec:framework}}

\subsection{Graduality Schemata for Free Theories}

We define by induction on execution traces satisfaction relations $M \vDash t$
between an ITree $M$ and a $\textrm{PTrace}$ $t$ and $f \vDash t$ between an
$R$-indexed family of ITrees and an $\textrm{OTrace}R$.
\begin{align*}
  {M \vDash \eta(V)} &= M \leadstomany \leaf(\eta(V))\\
  {M \vDash \mho} &= M \leadstomany \leaf(\mho)\\
  {M \vDash \sigma; o} &= M \leadstomany \sigma(f) \wedge f \vDash o\\
\end{align*}
\begin{align*}
  f \vDash \textrm{end} &= \top \\
  f \vDash r; p &= f(r) \vDash p
\end{align*}
%
Additionally, we inductively define an error ordering $\sqsubseteq$ on player
and opponent traces which intuitively means if $p \sqsubseteq q$ then $p$ is
either equal to $q$ or it is a prefix of $q$ that ends in $\mho$.
\begin{mathpar}
  %\footnotesize
  \mho \sqsubseteq p

  \eta V \sqsubseteq \eta V

  \inferrule{o \sqsubseteq o'}{\sigma; o \sqsubseteq \sigma; o'}

  {\textrm{end} \sqsubseteq \textrm{end}}

  \inferrule{p \sqsubseteq q}{r; p \sqsubseteq r; q}
\end{mathpar}

Next, we state two lemmas about the coinductive error ordering and bisimilarity
on ITrees. The proofs are by a straightforward coinductive argument.

\begin{lemma}[Properties of $\sqsubseteq$]
  \label{lem:itree-strong-err}
  Assume $M \sqsubseteq N$.
  \begin{enumerate}
  \item If $M \leadsto \eta x$ then $N \leadsto \eta x$
  \item If $M \leadsto \sigma(f)$ then $N \leadsto \sigma(g)$ with $\forall r\in \ar(\sigma). f(r) \sqsubseteq g(r)$.
  \item If $N \leadsto \err$ then $M \leadsto \err$
  \item If $N \leadsto \eta x$ then $M \leadsto \err$ or $M \leadsto \eta x$
  \item If $N \leadsto \sigma(f)$ then $M \leadsto \err$ or $M \leadsto \sigma(g)$ with $\forall r \in \ar(\sigma). f(r) \sqsubseteq g(r)$
  \end{enumerate}
\end{lemma}

\begin{lemma}[Properties of $\bisim$]
  \label{lem:itree-bisim}
  Assume $M \sqsubseteq N$.
  \begin{enumerate}
  \item If $M \leadsto \err$ then $N \leadsto \err$
  \item If $M \leadsto \eta x$ then $N \leadsto \eta x$
  \item If $M \leadsto \sigma(f)$ then $N \leadsto \sigma(g)$ with $\forall r\in \ar(\sigma). f(r) \bisim g(r)$.
  \item If $N \leadsto \err$ then $M \leadsto \err$
  \item If $N \leadsto \eta x$ then $M \leadsto \eta x$
  \item If $N \leadsto \sigma(f)$ then $M \leadsto \sigma(g)$ with $\forall r \in \ar(\sigma). f(r) \bisim g(r)$
  \end{enumerate}
\end{lemma}

Now we prove Theorem \ref{thm:itrees-ext-ord} about the extensional ordering $\ltbisim$ on ITrees:
\begin{theorem}[Properties of $\ltbisim$]
  $\ltbisim$ is transitive.
  If $M \sqsubseteq N$ then $M \ltbisim N$.
  If $M \bisim N$ then $M \ltbisim N$.
\end{theorem}
\begin{proof}
  Transitivity is clear. The other two implications follow by induction on the
  structure of the input trace and Lemmas \ref{lem:itree-strong-err} and
  \ref{lem:itree-bisim}.
\end{proof} 

\subsection{Exception Transformer for Graduality Schemata}

In this section, we show that the exception monad transformer acts on
$\calA$-graduality schemata. More concretely, the goal of this section is to
prove the following theorem:
%
\begin{theorem}[Exception Transformer for Graduality Schemata]
  Let $E$ be a finite type. Let $\calA$ be an ordinary algebraic theory, and
  suppose we are given an $\calA$-graduality schema. Then there is a graduality
  schema for the transformed theory $(\calA + \exn{E})$.
\end{theorem}

We first show that if $T$ gives the free models of $\calA$ on Set, then the
error monad transformer applied to $T$ gives the free models of the theory of
$\calA + \exn{E}$ on Set:
%
\begin{lemma}\label{lem:error-monad-transformer-free-models}
    Let $\calA$ be a guarded algebraic theory whose free models are given by
    $T$. Then the free models of the theory $\calA + \exn{E}$ are given by
    $T'(A) = T(A + E)$.
\end{lemma}
\begin{proof}
    Let $\calA = (\Sigma_\calA, \text{Eqn}_\calA)$ be a guarded algebraic theory with
    free-model monad given by $T_\calA$. Let $T'(A) = T(A + E)$.

    We first show that $T'(A)$ implements each operation in $\calA +
    \exn{E}$. For each $e : E$, we define 
    %
    \[ \sem{\ex{e}} = T.\eta(\inr\, e). \]
    %
    For each operation $\oper : (\tau, \gamma)$ in $A$, we define $\sem{\oper} :
    ((\tau_\oper \to T'(A)) \times (\gamma_\oper \to \later (T'(A)))) \to T'(A)$
    using the fact that $T(A+E)$ implements the operations, i.e, we define
    %
    \[ \sem{\oper}(u, g) = T(A+E).\sem{\oper}(u, g) \]
    
    We define the inclusion of generators $\eta' : A \to T'(A)$ by
    %
    \[ \eta'(x) = T(A+E).\eta(\inl\, x). \]

    Next, we show that the equations are satisfied. Since there are no equations
    in $\exn{E}$, we need only show that the equations of $\calA$ hold.
    Specifically, let $e = (\Gamma, l, r)$ be an equation in $\text{Eqn}_\calA$, where
    $l, r$ are terms in the signature $\Sigma_\calA$ with variables in $\Gamma$.
    Consider an arbitrary $f : \Gamma \to T'(A)$. We need to show that
    $\sem{l}_f = \sem{r}_f$. Since we know that $l$ and $r$ do not mention the operations
    $\ex{e}$, we know that this equality holds by the fact that $T(A + E)$ is a
    model for $\calA$.

    Finally, we show the universal property. Let $B$ be a model for $\calA +
    \exn{E}$ and let $f : A \to |B|$. We define $\fbar : T'(A) \to |B|$ using
    the universal property of $T(A + E)$. Namely, since $B$ is also a model for
    $\calA$, it suffices to define a function $g : A + E \to |B|$. We define
    %
    \[ g(\inl\, x) = f(x), \] 
    and
    \[ g(\inr\, e) = B.\ex{e}. \]
    %
    Then by the universal property this extends uniquely to a function $\ghat :
    T(A + E) \to |B|$ that preserves the operations in $\calA$.
    %
    It is clear that this is a homomorphism of models of $\calA +
    \exn{E}$ since for each $e : E$, we have $g(\sem{\ex{e}}) = \ghat(T.\eta(\inr\, e)) = g(\inr\, e) =
    B.\ex{e}$.

    For uniqueness, suppose $h, h' : T(A + E) \to B$ are homomorphisms of
    $(\calA + \exn{E})$-models such that $h \circ \eta' = h' \circ \eta'$. We
    show that $h = h'$. Since $h$ and $h'$ are in particular homomorphisms of
    $\calA$-models, by the universal property for $T(A + E)$, it suffices to
    show that they agree when composed with $T(A+E).\eta$. Thus, let $x \in A +
    E$. We proceed by cases on $x$. If $x = \inl\, a$, then we have
    %
    \[ 
      h (\eta (\inl\, a)) 
      = h (\eta'(a)) 
      = h' (\eta'(a)) 
      = h' (\eta (\inl\, a)). 
    \]
    %
    Now suppose $x = \inr\, e$. Then we have
    \[ 
      h (\eta (\inr\, e)) 
      = h(\sem{\ex{e}}) 
      = B.\ex{e} 
      = h'(\sem{\ex{e}}) 
      = h' (\eta (\inr\, e)).  
    \]

\end{proof}

Now we show that the monad transformer lifts to relations:
%
\begin{lemma}\label{lem:error-monad-transformer-free-domains}
    Let $\calA$ be an ordinary algebraic theory. Let $\calA^G \in \calA[\mho,
    \theta]$, i.e., $\calA^G$ is a guarded extension of $\calA$ with an error
    and a $\theta$ operation. Suppose that the adjunction between sets and
    models of $\calA^G$ lifts to a strong adjunction between predomains and
    $\calA^G$-domains.
    
    Then for the theory $\calA^{G+} := (\calA^{G} + \exn{E})$, the adjunction
    between sets and models of $\calA^{G+}$ lifts to a strong adjunction between
    predomains and $\calA^{G+}$-domains.
\end{lemma}
\begin{proof}
    % We define $\calA^{G+}$ to be $\calA^{G} + \mho$. % i.e., exceptions
    Let the free models of $\calA^G$ be given by $T$, and let the free domains
    of $\calA^G$ be given by $\That$. By assumption, for any predomain $A$, we
    have $|\That(A)| = T(|A|)$.
    
    We know by Lemma \ref{lem:error-monad-transformer-free-models} that at the
    level of sets, the free models of $\calA^{G+}$ are given by $T'(A) = T(A +
    E)$.
    
    We show that $\That'(A) := \That(A + E)$ gives the free domains of
    $\calA^{G+}$. In $A + E$, the $+$ symbol denotes the coproduct of
    predomains, where in the ordering and bisimilarity relations, $\inl$ is
    related to $\inl$ and $\inr$ to $\inr$, and since $E$ is merely a type, the
    ordering and bisimilarity are equality.

    The proof proceeds in several steps:
    \begin{itemize}
        \item We define the strong error ordering and weak bisimilarity on
        $\That'(A)$ to be those of $\That(A + E)$. These relations satisfy the
        properties required of the strong error ordering and weak bisimilarity
        relations, since the relations on $\That(A)$ satisfy the properties for
        all $A$.

        \item For a predomain relation $c : A \rel A'$, we define the lifting
        $\That'(c) = \That(c + =_E)$.

        \item We check the universal property of the lifting, as follows.
        Consider a square $f \ltsq{c}{Ud} g$ where $c : A \rel A'$ is a
        predomain relation, and $d : B \rel B'$ is a relation of
        $\calA^{G+}$-domains.
        %
        Then we need to show that $f^* \ltsq{\That'(c)}{d} g^*$, where $-^*$
        denotes the extension operation turning a function $A \to UB$ into a
        homomorphism $T'(A) \to B$. By the definition of $-^*$ and the universal
        property of the lifting for $\That$, it suffices to construct a
        predomain square
        %
        $(f + B.\ex{\cdot}) \ltsq{c + =_E}{UB} (g + B.\ex{\cdot})$.
        %
        We proceed by cases on the input. In case both are $\inl$, then we use
        the square $f \ltsq{c}{Ud} g$. If both are $\inr e$, then we need to
        show that $(B.\ex{e}) \binrel{d} (B'.\ex{e})$ which is true as $d$ is a
        relation of $\calA^{G+}$ domains and hence is a congruence with respect
        to $\ex{e}$.


        \item It is easily checked that the operations are monotone and
        bisimilarity-preserving, and that the strength
        $t : A \times T'(A') \to T'(A \times A')$ is a morphism of predomains.
        \item Lastly, for any predomain $A$, we have   
        \begin{align*}
            |\That'(A)| 
            &= |\That(A + E)| \\
            &= T(|A + E|)\\ 
            &= T(|A| + E)\\ 
            &= T'(|A|).
        \end{align*}
    \end{itemize}
\end{proof}

Now we are ready to define the action of the exception transformer on graduality
schemata. Let $\calA$ be an ordinary algebraic theory, and let $\calG$ be a
graduality schema for $\calA$.
%
We write $\calAs, \calAsb$ for the guarded extensions of $\calA$ associated with
$\calG$, and we write $T$ for the free-model monad for $\calAs$ on Set, and
$\That$ for its lifting to predomains.
%
We define a graduality schema for $(\calA + \exn{E})$, as follows.
\begin{itemize}
    \item We define the guarded extensions 
    $\calAgplus := \calAs + \exn{E}$ and 
    $\calAgplus_b := \calAsb + \exn{E}$.

    \item Then by Lemma \ref{lem:error-monad-transformer-free-models}, the free models of
    $\calAgplus$ on Set are given by $T'(A) := T(A + E)$.

    \item By Lemma \ref{lem:error-monad-transformer-free-domains}, we have that 
    the free domains for $\calAgplus$ are given by $\That(A + E)$.

    \item We define the extensional types $Z'(A) := Z(A + E)$ where $Z$ is the
    family of extensional types for $\calG$ (we use $Z$ here to avoid confusion
    with the type $E$ of exceptions). Next, we define the function
    %
    $\col'_A : T'^{gl}(A) \to Z'(A)$, i.e., $T^{gl}(A + E) \to Z(A + E)$ as 
    %
    \[ \col'_A := \col_{A + E}, \]
    %
    where $\col$ is the family of functions associated with the schema $\calG$.
    
    \item Lastly, the relation $\ltbisim_{Z'(A)}$ is defined to be
    $\ltbisim_{Z(A + E)}$. Then since each $\ltbisim_{Z(A)}$ is a partial order,
    in particular $\ltbisim_{Z(A + E)}$ is a partial order. Furthermore, it is
    clear that the error ordering and bisimilarity relations on the globalized
    monad imply $\ltbisim_{Z'(A)}$ on the output of $\col'_A$, since this holds for
    $\ltbisim_{Z(A)}$ for all $A$.
\end{itemize}



\subsection{State Transformer for Graduality Schemata}

In this section, we show that the state monad transformer acts on
$\calA$-graduality schemata. We prove the following theorem:
\begin{theorem}[Global State Transformer for Graduality Schemata]
  Let $S$ be a finite type. Let $\calA$ be an ordinary algebraic theory, and
  suppose we are given an $\calA$-graduality schema. Then there is a graduality
  schema for the transformed theory $(\calA \otimes \state{S})$.
\end{theorem}

% In particular:
% \begin{itemize}
%     \item For any guarded algebraic theory $\calA$, if the free models of
%     $\calA$ are given by $T$, then the free models of $\calA \otimes \state{S}$
%     are given by $S \to T(S \times A)$.
% \end{itemize}

First, we show that for any guarded algebraic theory $\calAg$, if the free
models of $\calAg$ are given by $T$, then the free models of $\calAg \otimes
\state{S}$ are given by $S \to T(S \times A)$. (This is Theorem
\ref{thm:state-monad-transformer} in the paper.)

Let $\calAg = (\Sigma, E)$ be a guarded algebraic theory with free-model monad
$T$. We will show that the free-model monad for the theory $\calAg \otimes
\state{S}$ is given by
%
\[ T'(A) := S \to T(S \times A). \]
%
We begin by defining the operations. We let

\[ \sem{\get} f\, s = f\, s\, s, \]
%
\[ \sem{\putop_s} m\, s' = m\, s, \]  
%
and we define $\sem{\oper} : [(\tau_\oper \to T'(A)) \times (\gamma_\oper \to \later(T'(A)))] \to T'(A)$ by
%
\[ 
  \sem{\oper}(u, g)\, s = 
    T(A \times S).\sem{\oper}
      [(\lambda w. u\, w\, s), (\lambda z. \lambda t. g\, z\, t\, s)].
\]

Now, we show that the required equations hold when interpreted in $T'(A)$.
%
The proofs of the get-put laws are analogous to those for the ordinary state
monad and hence omitted.
%
We prove the tensor equations stating that $\get$ and $\putop$ commute with each
operation in $\calA$:
\begin{enumerate}
    \item \emph{Get commutes with each operation in $\calAg$}: Let $\oper :
    (\tau, \gamma) \in \calAg$. Let $u : \tau \to S \to T'(A)$ and $g : \gamma \to
    \later(S \to T'(A))$. Let
    %
    \[ 
        M = \sem{\oper}[
              (\lambda w. \sem{\get}(\lambda s. u\, w\, s)),
              (\lambda z. \lambda t.\sem{\get}(\lambda s. g\, z\, t\, s))]
    \] 
    and
    \[
        N = \sem{\get}[\lambda s. 
              \sem{\oper}((\lambda w. u\, w\, s), (\lambda z. \lambda t. g\, z\, t\, s))]
    \]
    %
    We need to show that $M = N$. For brevity, let
    %
    $u_1 = (\lambda w. \sem{\get}(\lambda s. u\, w\, s))$ and let 
    %
    $g_1 = (\lambda z. \lambda t.\sem{\get}(\lambda s. g\, z\, t\, s))$.
    %
    We first observe that for any $w : \tau$ and $z : \gamma$, we have
    %
    $u_1(w) = (\lambda s. u\, w\, s\, s)$ of type $S \to T(S \times A)$ and 
    %
    $g_1(z) = (\lambda t. \lambda s. g\, z\, t\, s\, s)$ of type $\later (S \to T(S \times A))$.
    
    Now let $s \in S$. We have
    \begin{align*}
        M(s) &= \sem{\oper}(u_1, g_1)(s) \\
             &= T(A \times S).\sem{\oper}[(\lambda w. u_1\, w\, s), (\lambda z. \lambda t. g_1\, z\, t\, s)] \\
             &= T(A \times S).\sem{\oper}[(\lambda w. u\, w\, s\, s), (\lambda z. \lambda t. g\, z\, t\, s\, s) ]. \\
    \end{align*}
    
    Now let $p(s) = (\lambda w. u\, w\, s)$ and $q(s) = (\lambda z. \lambda t. g\, z\, t\, s)$.
    We have
    %
    \begin{align*}
        N(s) &= (\sem{\get}[\lambda s'. \sem{\oper}(p(s'), q(s'))])(s) \\
             &= (\lambda s'. \sem{\oper}(p(s'), q(s')))(s)(s) \\
             &= [\sem{\oper}(p(s), q(s))](s) \\
             &= T(A \times S).\sem{\oper}[(\lambda w. p(s)\, w\, s), (\lambda z. \lambda t. q(s)\, z\, t\, s)] \\
             &= T(A \times S).\sem{\oper}[(\lambda w. u\, w\, s\, s), (\lambda z. \lambda t. g\, z\, t\, s\, s)].
    \end{align*}
    %
    Thus, we conclude that $M = N$.
    
    \item \emph{Put commutes with each operation in $\calAg$}: This is similar to
    the above and hence omitted.
\end{enumerate}

Finally, we prove that the equations in $\calAg$ hold. For this, we first prove
the following lemma:
%
\begin{lemma}
    % Let $t$ be a term in the signature $F_\Sigma(A)$ (recall $\Sigma$ is the
    % signature of $\calA$), and let $f : \Gamma \to T'(A)$ be a function. 

    Let $f : \Gamma \to T'(A)$ be a function. Then for all $s \in S$, there exists a
    function $\fhat_s : \Gamma \to T(S \times A)$ such that, for all terms $t$ in
    $F_\Sigma(\Gamma)$, we have
    %
    \[ \sem{t}_{\fhat_s}^{T(S \times A)} = (\sem{t}_f^{T'(A)})(s). \]
\end{lemma}

\begin{proof}
    For $s \in S$, we define $\fhat_s(x) = f(x)(s)$.
    %
    The equality follows by induction on the term $t$:
    \begin{itemize}
        \item If $t$ is a variable $x \in \Gamma$, then we have
        %
        \[ \sem{x}_{\fhat_s}^{T(S \times A)} 
          = \fhat_s(x) 
          = f(x)(s) 
          = (\sem{x}_f^{T'(A)})(s).
        \]

        \item Now suppose $t = \oper(u, g)$, where (1) $\oper \in \Sigma$, 
        (2) $u : \tau_\oper \to F_\Sigma(\Gamma)$ and 
        (3) $g : \gamma_\oper \to \later (F_\Sigma(\Gamma))$. Then the left-hand side is
        %
        \begin{align*}
            \sem{\oper(u, g)}_{\fhat_s}^{T(S \times A)}
              &= T(S \times A).\sem{\oper}(u_1, g_1),
        \end{align*}
        %
        where $u_1(w) := \sem{u(w)}_{\fhat_s}^{T(S \times A)}$ and 
        $g_1(z)(t) := \sem{g(z)(t)}_{\fhat_s}^{T(S \times A)}$.

        The right-hand side is
        \begin{align*}
            (\sem{\oper(u, g)}_f^{T'(A)})(s)
              &= [T'(A).\sem{\oper}(u_2, g_2)](s) \\
              &= T(S \times A).\sem{\oper}[
                  (\lambda w. u_2\, w\, s), 
                  (\lambda z. \lambda t. g_2\, z\, t\, s)],
        \end{align*}
        where $u_2(w) := \sem{u(w)}_f^{T'(A)}$ and $g_2(z)(t) := \sem{g(z)(t)}_f^{T'(A)}$.
        
        Then by the inductive hypothesis, we have $u_1(w) = u_2(w)(s)$ and
        $g_1(z)(t) = g_2(z)(t)(s)$, so the above becomes
        %
        \[ T(S \times A).\sem{\oper}[
                  (\lambda w. u_1\, w), 
                  (\lambda z. \lambda t. g_1\, z\, t)].
        \]
        %
        Thus, both sides are equal.

    \end{itemize}
\end{proof}

Now we can show that the equations in $\calAg$ hold:
\begin{proof}
    Let $e = (\Gamma, l, r) \in
    E$, and let $f : \Gamma \to T'(A)$ be a function. We need to show

    \[ \sem{l}_f^{T'(A)} = \sem{r}_f^{T'(A)}, \]

    where both sides are functions $S \to T(S \times A)$. To this end, let $s \in
    S$. Then by the above Lemma, we have that there is $\fhat_s : \Gamma \to T(S \times
    A)$ with the property described in the Lemma. In particular, we have
    %
    \begin{align*}
        (\sem{l}_f^{T'(A)})(s) 
        &= \sem{l}_{\fhat_s}^{T(S \times A)} \\
        &= \sem{r}_{\fhat_s}^{T(S \times A)} \\
        &= (\sem{t}_f^{T'(A)})(s),
    \end{align*}
    %
    where the second equality holds because $T(S \times A)$ is a model of $\calAg$
    and hence the equation holds when interpreted in $T(S \times A)$.
\end{proof}

Next, we define $\eta' : A \to T'(A)$ by 
%
$\eta'(x) = \lambda s.\eta(s, x)$, 
%
where $\eta : (S \times A) \to T(S \times A)$ is the unit for $T$.

Finally, we show the universal property.
\begin{proof}
    Let $B$ be a model of the theory $\calAg \otimes \state{S}$, and let $f : A \to UB$.
    %
    First, define $g : A \times S \to B$ as the following composition:
    %
    \[ g = A \times S \xrightarrow{f \times \id} UB \times S \xrightarrow{B.\sem{\putop} B} \]
    %
    We then define $\fbar : T'(A) \to B$ as the following composition:
    %
    \[ \fbar = 
        S \to T(S \times A) \xrightarrow{S \to \gbar} S \to B \xrightarrow{B.\sem{\get}} B.  \]
    %
    where $\gbar$ is the unique extension of $g$ to a homomorphism of $\calAg$-models.


    Then it is easily checked that $\fbar$ is a homomorphism of models of $\calAg
    \otimes {\state{S}}$, and that $\fbar \circ \eta' = f$.
\end{proof}

Now we show uniqueness. Let $\phi, \phi' : T'(A) \to B$ be homomorphisms of
$(\calAg \otimes {\state{S}})$-models, and suppose $\phi(\eta'\, x) =
\phi'(\eta'\, x)$ for all $x$. That is, $\phi(\lambda s. \eta(s,x)) =
\phi'(\eta(s,x))$. Then we claim that $\phi = \phi'$.

First, we need a small lemma:
\begin{lemma}
    For all $m : T(S \times A)$, we have $\phi(\lambda s.m) = \phi'(\lambda s.m)$.
\end{lemma}
\begin{proof}
    Define $\psi, \psi' : T(S \times A) \to B$ by $\psi(m) = \phi(\lambda s.m)$
    and $\psi'(m) = \phi'(\lambda s.m)$. First, note that $\psi$ and $\psi'$ are
    homomorphisms of $\calAg$-models.
    %
    Thus, by the universal property of $T(S \times A)$ as the free-model monad
    for $\calAg$, it suffices to show that $\psi \circ \eta = \psi' \circ \eta$.
    We have
    \begin{align*}
        \psi(\eta(s,a)) 
        &= \phi(\lambda s'.\eta(s, a)) \\
        &= \phi(\putop_s (\lambda s''.\eta(s'', a))) \\
        &= B.\putop_s(\phi(\lambda s''.\eta(s'', a))) \\
        &= B.\putop_s(\phi'(\lambda s''.\eta(s'', a))) \\
        &= \phi'(\putop_s(\lambda s''.\eta(s'', a))) \\
        &= \phi'(\lambda s'.\eta(s, a)) \\
        &= \psi'(\eta(s,a)).
    \end{align*}
\end{proof}

Now let $x : T'(A)$. We have
\begin{align*}
    \phi(x)
    &= \phi(\get(\lambda s. \putop_s(x))) \\
    &= B.\get(\lambda s. \phi(\putop_s(x))) \\
    &= B.\get(\lambda s. \phi(\lambda s'. x\, s)) \\
    &= B.\get(\lambda s. \phi'(\lambda s'. x\, s)) \\
    &= B.\get(\lambda s. \phi'(\putop_s(x))) \\
    &= \phi'(\get(\lambda s.\putop_s(x))) \\
    &= \phi'(x).
\end{align*}

Next, we show that the monad transformer lifts to relations:
\begin{lemma}\label{lem:state-monad-transformer-free-domains} Let $\calA$ be an
    ordinary algebraic theory. Let $\calA^G \in \calA[\mho, \theta]$, be a
    guarded extension of $\calA$. Suppose that the adjunction between sets and
    models of $\calA^G$ lifts to a strong adjunction between predomains and
    $\calA^G$-domains, with the free domains being given by $\That$.
    
    Then for the theory $\calA^{G\otimes} := (\calA^{G} \otimes \state{S})$, the
    adjunction between sets and models of $\calA^{G\otimes}$ lifts to a strong
    adjunction between predomains and $\calA^{G\otimes}$-domains.
    The free domains for $\calA^{G\otimes}$ are given by $S \to \That(S \times A)$.
\end{lemma}
\begin{proof}
    % We define $\calA^{G+}$ to be $\calA^{G} + \mho$. % i.e., exceptions
    Let the free models of $\calA^G$ be given by $T$, and let the free domains
    of $\calA^G$ be given by $\That$. By assumption, for any predomain $A$, we
    have $|\That(A)| = T(|A|)$.
    
    We know by the above that at the level of sets, the free models of
    $\calA^{G+}$ are given by $T'(A) = S \to T(S \times A)$.
    
    We show that $\That'(A) := S \to \That(S \times A)$ gives the free domains
    of $\calA^{G+}$. Here, $S$ is a discrete predomain, so the relations are
    given by equality. In addition, $\to$ denotes the predomain of morphisms of
    predomains between $S$ and $\That(S \times A)$ (this is the same as ordinary
    functions from $S$ to $\That(S \times A)$), and $\times$ denotes the product
    of predomains.

    \begin{itemize}
        \item The strong error ordering and weak bisimilarity on $\That'(A)$ are
        those of the predomain $S \to \That(S \times A)$.

        \item For a predomain relation $c : A \rel A'$, the lifting is given by
        $\That'(c) = S \to \That(S \times c)$, where $\That(S \times c)$ is the
        lifting of the relation $=_S \times c$.
        %
        More concretely, we have $f \binrel{\That'(c)} g$ iff for all $s \in S$,
        we have $(f\, s) \binrel{\That(S \times c)} (g\, s)$.

        \item To verify the universal property of the lifting of a relation, we
        use the definition of the unique extension of a function $A \to UB$ to a
        homomorphism out of $T'(A)$, as defined in the proof of Theorem
        \ref{thm:state-monad-transformer} above. The universal property for the
        lifted relation then follows using the actions of $\to$ and $\times$ on
        squares as well as the universal property of $\That$ on relations.

        \item It is easily checked that the operations are monotone and
        bisimilarity-preserving, and that the strength $t : A \times T'(A') \to
        T'(A \times A')$ is a morphism of predomains.
        
        \item Lastly, for any predomain $A$, we have
        %
        \begin{align*}
            |\That'(A)| &= |S \to \That(S \times A)| \\
                        &= S \to |\That(S \times A)| \\
                        &= S \to T(|S \times A|)  \\
                        &= T'(|A|).
        \end{align*}
    \end{itemize}
\end{proof}

Now we are ready to define the action of the global state transformer on
graduality schemata. Let $\calA$ be an ordinary algebraic theory, and let
$\calG$ be a graduality schema for $\calA$.
%
We write $\calAs, \calAsb$ for the guarded extensions of $\calA$ associated with
$\calG$, and we write $T$ for the free-model monad for $\calAs$ on Set, and
$\That$ for its lifting to predomains.
%
We define a graduality schema for $(\calA \otimes \state{S})$, as follows.
\begin{itemize}
    \item We define the guarded extensions 
    $\calAgtimes := \calAs \otimes \state{S}$ and 
    $\calAgtimes_b := \calAsb \otimes \state{S}$.

    \item Then by Theorem \ref{thm:state-monad-transformer}, the free models of
    $\calAgtimes$ on Set are given by $T'(A) := S \to T(S \times A)$.

    \item By Lemma \ref{lem:state-monad-transformer-free-domains}, we have that 
    the free domains for $\calAgtimes$ are given by $S \to \That(S \times A)$.

    % \item The functor $T'$ commutes with clock quantification. First note that
    % $S$ is finite and is hence clock-irrelevant. Then it follows from Theorem
    % 4.2 of kristensen et al. \cite{kristensen-mogelberg-vezzosi2022}, that $T'$
    % commutes with clock quantification.
    % %
    % Then if $A$ is clock-irrelevant, we have $\forall k. T'_k(A) \cong S \to T'$ provided $A$ is clock-irrelevant.

    \item We define the extensional types $E'(A) := S \to E(S \times A)$, where
    $E$ is the family of extensional types for $\calG$. Next, we define the
    family of functions
    %
    $\col'_A : T'^{gl}(A) \to E'(A)$, i.e., $(S \to T^{gl}(S \times A)) \to (S \to E(S \times A))$, via 
    %
    \[ \col'_A(f) := \col_{S \times A} \circ f, \]
    %
    where $\col$ is the family of functions associated with the schema $\calG$.

    \item Lastly, the relation $\ltbisim_{E'(A)}$ is defined to be $S \to
    \ltbisim_{E(S \times A)}$, i.e., we have $f \binrel{\ltbisim_{E'(A)}} g$ iff
    for all $s \in S$, we have $(f\, s) \binrel{\ltbisim_{E(S \times A)}} (g\, s)$.
    
    Then we observe that $\ltbisim_{E'(A)}$ is a partial order for each $A$. We
    now show that the error ordering and bisimilarity relations on the
    globalized monad imply $\ltbisim_{E'(A)}$ on the output of $\col'_A$.
    %
    More concretely, let $f, g : S \to T^{gl}(S \times A)$, and suppose that 
    $f \ltdyn_{S \to T^{gl}(S \times A)} g$. Then we want to show that
    $\col'_A(f) \ltbisim_{S \to E(S \times A)} \col'_A(g)$. Let $s \in S$.
    It suffices to show
    %
    \[ \col_{S \times A}(f(s)) \ltbisim_{E(S \times A)} \col_{S \times A}(g(s)). \] 
    %
    This follows from the fact that $f(s) \ltdyn_{T^{gl}(S \times A)} g(s)$.

    The same reasoning holds for bisimilarity.
\end{itemize}

\subsection{Writer Transformer for Graduality Schemata}

In this section, we show that the writer monad transformer acts on
$\calA$-graduality schemata. We prove the following theorem:
\begin{theorem}[Writer State Transformer for Graduality Schemata]
  Let $W$ be a monoid. Let $\calA$ be an ordinary algebraic theory, and
  suppose we are given an $\calA$-graduality schema. Then there is a graduality
  schema for the transformed theory $(\calA \otimes \writer{W})$.
\end{theorem}

We first show that if $T$ gives the free models of $\calA$, then the writer monad
transformer applied to $T$ gives the free models of the theory of $\calA \otimes
\writer{W}$:
\begin{lemma}\label{lem:writer-monad-transformer-free-models}
    Let $\calA$ be a guarded algebraic theory whose free models are given by
    $T$. Then the free models of the theory $\calA \otimes \writer{W}$ are given by
    $T'(A) = T(W \times A)$.
\end{lemma}
\begin{proof}
    Let $\calA = (\Sigma_\calA, \text{Eqn}_\calA)$ be a guarded algebraic theory with
    free-model monad given by $T_\calA$. Let $T'(A) = T(W \times A)$.

    We first show that $T'(A)$ implements each operation in $\calA \otimes
    \writer{W}$. For each $w : W$, we define 
    %
    \[ \sem{\wrt{w}} = T(\lambda(w',a).\, (w'w, a)). \]
    %
    For each operation $\oper : (\tau, \gamma)$ in $A$, we define $\sem{\oper} :
    ((\tau_\oper \to T'(A)) \times (\gamma_\oper \to \later (T'(A)))) \to T'(A)$
    using the fact that $T(W \times A)$ implements the operations, i.e, we define
    %
    \[ \sem{\oper}(u, g) = T(W \times A).\sem{\oper}(u, g), \]
    %
    where $T$ is the functorial action of the monad.
    
    We define the inclusion of generators $\eta' : A \to T'(A)$ by
    %
    \[ \eta'(x) = T(W \times A).\eta(\epsilon, x), \]
    %
    where $\epsilon$ is the unit of the monoid $W$.

    Next, we show that the equations are satisfied. First, the writer equations
    $\wrt{\epsilon}(x) = x$ and $\wrt{w_1}(\wrt{w_2}(x)) = \wrt{w_1 \cdot
    w_2}(x)$ follow by functoriality and associativity of the monoid.

    Now we show that the equations of $\calA$ hold. Specifically, let $e =
    (\Gamma, l, r)$ be an equation in $\text{Eqn}_\calA$, where $l, r$ are terms
    in the signature $\Sigma_\calA$ with variables in $\Gamma$. Consider an
    arbitrary $f : \Gamma \to T'(A)$. We need to show that $\sem{l}_f =
    \sem{r}_f$. Since we know that $l$ and $r$ do not mention the operations
    $\wrt{w}$, we know that this equality holds by the fact that $T(W \times A)$
    is a model for $\calA$.

    Lastly, we claim that each operation $\wrt{w}$ commutes with each operation
    $\oper$ in $\calA$. This follows from the fact that the morphism $T(W \times
    A).\sem{\oper}(u, g)$ is a homomorphism, i.e., preserves the operations in
    $\calA$.

    Finally, we show the universal property. Let $B$ be a model for $\calA \otimes
    \writer{W}$ and let $f : A \to |B|$. We define $\fbar : T'(A) \to |B|$ using
    the universal property of $T(W \times A)$. Namely, since $B$ is also a model for
    $\calA$, it suffices to define a function $g : W \times A \to |B|$. We define
    %
    \[ g(w,x) = B.\wrt{w}(f(x)). \] 
    %
    Then by the universal property this extends uniquely to a function $\ghat :
    T(W \times A) \to |B|$ that preserves the operations in $\calA$.
    %
    Then using the universal property of $T(W \times A)$ and the fact that $B$
    satisfies the associativity equation of $\writer{W}$, we can show that
    $\ghat$ is a homomorphism of models of $\calA \otimes \writer{W}$.
    %
    It is also straightforward to show that $\ghat \circ \eta' = f$, using the
    fact that $B$ satisfies the identity equation of $\writer{W}$.

    Uniqueness of $\ghat$ follows from the corresponding property for $T$.
\end{proof}


Next, we show that the monad transformer lifts to relations:
\begin{lemma}\label{lem:writer-monad-transformer-free-domains}
    Let $\calA$ be an ordinary algebraic theory. Let $\calA^G \in \calA[\mho,
    \theta]$, be a guarded extension of $\calA$. Suppose that the adjunction
    between sets and models of $\calA^G$ lifts to a strong adjunction between
    predomains and $\calA^G$-domains, with the free domains being given by
    $\That$.
    
    Then for the theory $\calA^{G\otimes} := (\calA^{G} \otimes \writer{W})$, the
    adjunction between sets and models of $\calA^{G\otimes}$ lifts to a strong
    adjunction between predomains and $\calA^{G\otimes}$-domains.
    The free domains for $\calA^{G\otimes}$ are given by $\That(W \times A)$.
\end{lemma}
\begin{proof}
    Let the free models of $\calA^G$ be given by $T$, and let the free domains
    of $\calA^G$ be given by $\That$. By assumption, for any predomain $A$, we
    have $|\That(A)| = T(|A|)$.
    
    We know by the above that at the level of sets, the free models of
    $\calA^{G+}$ are given by $T'(A) = T(W \times A)$.
    
    We show that $\That'(A) := \That(W \times A)$ gives the free domains of
    $\calA^{G+}$. Here, $W$ is a discrete predomain, so the relations are given
    by equality. In addition, $\times$ denotes the product of predomains.

    \begin{itemize}
        \item The strong error ordering and weak bisimilarity on $\That'(A)$ are
        those of the predomain $\That(W \times A)$.

        \item For a predomain relation $c : A \rel A'$, the lifting is given by
        $\That'(c) = \That(W \times c)$, i.e., the lifting of the relation $=_W \times c$.

        \item To verify the universal property of the lifting of a relation, we
        use the definition of the unique extension of a function $A \to UB$ to a
        homomorphism out of $T'(A)$, as defined in the proof of Lemma
        \ref{lem:writer-monad-transformer-free-models} above. The universal
        property for the lifted relation then follows using the action of
        $\times$ on squares as well as the universal property of $\That$ on
        relations.

        \item It is easily checked that the operations are monotone and
        bisimilarity-preserving, and that the strength $t : A \times T'(A') \to
        T'(A \times A')$ is a morphism of predomains.
        
        \item Lastly, for any predomain $A$, we have
        %
        \begin{align*}
            |\That'(A)| &= |\That(W \times A)| \\
                        &= T(|W \times A|) \\
                        &= T(W \times |A|)  \\
                        &= T'(|A|).
        \end{align*}
    \end{itemize}
\end{proof}

Now we are ready to define the action of the writer transformer on
graduality schemata. Let $\calA$ be an ordinary algebraic theory, and let
$\calG$ be a graduality schema for $\calA$.
%
We write $\calAs, \calAsb$ for the guarded extensions of $\calA$ associated with
$\calG$, and we write $T$ for the free-model monad for $\calAs$ on Set, and
$\That$ for its lifting to predomains.
%
We define a graduality schema for $(\calA \otimes \writer{W})$, as follows.
\begin{itemize}
    \item We define the guarded extensions 
    $\calAgtimes := \calAs \otimes \writer{W}$ and 
    $\calAgtimes_b := \calAsb \otimes \writer{W}$.

    \item Then by Lemma \ref{lem:writer-monad-transformer-free-models}, the free models of
    $\calAgtimes$ on Set are given by $T'(A) := S \to T(S \times A)$.

    \item By Lemma \ref{lem:writer-monad-transformer-free-domains}, we have that 
    the free domains for $\calAgtimes$ are given by $\That(W \times A)$.

    \item We define the extensional types $E'(A) := E(W \times A)$, where
    $E$ is the family of extensional types for $\calG$. Next, we define the
    family of functions
    %
    $\col'_A : T'^{gl}(A) \to E'(A)$, i.e., $(T^{gl}(W \times A)) \to (E(W \times A))$, via 
    %
    \[ \col'_A := \col_{W \times A} \]
    %
    where $\col$ is the family of functions associated with the schema $\calG$.

    \item Lastly, the relation $\ltbisim_{E'(A)}$ is defined to be $\ltbisim_{E(W \times A)}$.
    %
    Then we observe that $\ltbisim_{E'(A)}$ is a partial order for each $A$. We
    now show that the error ordering and bisimilarity relations on the
    globalized monad imply $\ltbisim_{E'(A)}$ on the output of $\col'_A$.
    %
    More concretely, let $x, y : T^{gl}(W \times A)$, and suppose that 
    $x \ltdyn_{T^{gl}(W \times A)} y$. Then we want to show that
    $\col'_A(x) \ltbisim_{E(W \times A)} \col'_A(y)$.
    But by definition of $\col'$, this is the same as showing
    %
    \[ \col_{W \times A}(x) \ltbisim_{E(W \times A)} \col_{W \times A}(y). \] 
    %
    This follows from the fact that $x \ltdyn_{T^{gl}(W \times A)} y$.

    The same reasoning holds for bisimilarity.
\end{itemize}


\subsection{Reader Transformer for Graduality Schemata}

In this section, we show that the reader monad transformer acts on
$\calA$-graduality schemata. We prove the following theorem:
\begin{theorem}[Reader Transformer for Graduality Schemata]
  Let $R$ be a finite type. Let $\calA$ be an ordinary algebraic theory, and
  suppose we are given an $\calA$-graduality schema. Then there is a graduality
  schema for the transformed theory $(\calA \otimes \reader{R})$.
\end{theorem}

First, we relate the free models of a guarded algebraic theory $\calAg$ to those
of $\calAg \times \reader{R}$.
\begin{lemma}\label{lem:reader-monad-transformer-free-models} Let $\calA$ be a
    guarded algebraic theory whose free models are given by $T$. Then the free
    models of the theory $\calA \otimes \reader{R}$ are given by $T'(A) = R \to
    T(A)$.
\end{lemma}
\begin{proof}

Let $\calAg = (\Sigma, E)$ be a guarded algebraic theory with free-model monad
$T$. We will show that the free-model monad for the theory $\calAg \otimes
\state{S}$ is given by
%
\[ T'(A) := R \to T(A). \]
%
We begin by defining the operations. We let
%
$\sem{\lookup} f\, r = f\, r\, r,$
%
and we define $\sem{\oper} : [(\tau_\oper \to T'(A)) \times (\gamma_\oper \to \later(T'(A)))] \to T'(A)$ by
%
\[ 
  \sem{\oper}(u, g)\, r = 
    T(A).\sem{\oper}
      [(\lambda w. u\, w\, r), (\lambda z. \lambda t. g\, z\, t\, r)].
\]

Now, we show that the required equations hold when interpreted in $T'(A)$.
%
The proofs of the two laws from the reader theory are analogous to those for the
ordinary reader monad and hence omitted.
%
Next, we consider the tensor equations stating that $\get$ commutes with each
operation in $\calA$. The proof is analogous to the corresponding proof for the
global state monad, and hence omitted.

To prove that the equations in $\calAg$ hold, we first prove a following lemma,
analogous to the one used in the proof for the global state monad transformer.
%
\begin{lemma}
    Let $f : \Gamma \to T'(A)$ be a function. Then for all $r \in R$, there
    exists a function $\fhat_r : \Gamma \to T(R)$ such that, for all terms $t$
    in $F_\Sigma(\Gamma)$, we have
    %
    \[ \sem{t}_{\fhat_r}^{T(A)} = (\sem{t}_f^{T'(A)})(r). \]
\end{lemma}
The proof is similar to that of the analogous lemma for global state, and hence
omitted.

Now we can show that the equations in $\calAg$ hold:
Let $e = (\Gamma, t_1, t_2) \in
E$, and let $f : \Gamma \to T'(A)$ be a function. We need to show

\[ \sem{t_1}_f^{T'(A)} = \sem{t_2}_f^{T'(A)}, \]

where both sides are functions $R \to T(A)$. To this end, let $r \in R$.
Then by the above Lemma, we have that there is $\fhat_r : \Gamma \to T(A)$
with the property described in the Lemma. In particular, we have
%
\begin{align*}
    (\sem{t_1}_f^{T'(A)})(s) 
    &= \sem{t_1}_{\fhat_s}^{T(S \times A)} \\
    &= \sem{t_2}_{\fhat_s}^{T(S \times A)} \\
    &= (\sem{t_2}_f^{T'(A)})(s),
\end{align*}
%
where the second equality holds because $T(A)$ is a model of $\calAg$
and hence the equation holds when interpreted in $T(A)$.

Next, we define $\eta' : A \to T'(A)$ by 
%
$\eta'(x) = \lambda r.\eta(x)$, 
%
where $\eta : A \to T(A)$ is the unit for $T$.

Finally, we show the universal property.
%
Let $B$ be a model of the theory $\calAg \otimes \reader{R}$, and let $f : A \to UB$.
%
We define $\fbar : T'(A) \to B$ as the following composition:
%
\[ \fbar = 
    R \to T(A) \xrightarrow{R \to f^*} R \to B \xrightarrow{B.\sem{\lookup}} B.  \]
%
where $f^*$ is the unique extension of $f$ to a homomorphism of $\calAg$-models.

Then it is easily checked that $\fbar$ is a homomorphism of models of $\calAg
\otimes {\reader{R}}$, and that $\fbar \circ \eta' = f$.

The proof of uniqueness is similar to the corresponding proof for global state,
and is hence omitted.
\end{proof}

Next, we show that the monad transformer lifts to relations:
\begin{lemma}\label{lem:reader-monad-transformer-free-domains}
    Let $\calA$ be an ordinary algebraic theory. Let $\calA^G \in \calA[\mho,
    \theta]$, be a guarded extension of $\calA$. Suppose that the adjunction
    between sets and models of $\calA^G$ lifts to a strong adjunction between
    predomains and $\calA^G$-domains, with the free domains being given by
    $\That$.
    
    Then for the theory $\calA^{G\otimes} := (\calA^{G} \otimes \reader{R})$,
    the adjunction between sets and models of $\calA^{G\otimes}$ lifts to a
    strong adjunction between predomains and $\calA^{G\otimes}$-domains. The
    free domains for $\calA^{G\otimes}$ are given by $R \to \That(A)$.
\end{lemma}
\begin{proof}
    % We define $\calA^{G+}$ to be $\calA^{G} + \mho$. % i.e., exceptions
    Let the free models of $\calA^G$ be given by $T$, and let the free domains
    of $\calA^G$ be given by $\That$. By assumption, for any predomain $A$, we
    have $|\That(A)| = T(|A|)$.
    
    We know by the above that at the level of sets, the free models of
    $\calA^{G+}$ are given by $T'(A) = R \to T(A)$.
    
    We show that $\That'(A) := R \to \That(A)$ gives the free domains of
    $\calA^{G+}$. Here, $R$ is a discrete predomain, so the relations are given
    by equality. In addition, $\to$ denotes the predomain of morphisms of
    predomains between $R$ and $\That(A)$ (this is the same as ordinary
    functions from $R$ to $\That(A)$), and $\times$ denotes the product of
    predomains.

    \begin{itemize}
        \item The strong error ordering and weak bisimilarity on $\That'(A)$ are
        those of the predomain $R \to \That(A)$.

        \item For a predomain relation $c : A \rel A'$, the lifting is given by
        $\That'(c) = R \to \That(c)$.
        %
        More concretely, we have $f \binrel{\That'(c)} g$ iff for all $r \in R$,
        we have $(f\, r) \binrel{\That(c)} (g\, r)$.

        \item To verify the universal property of the lifting of a relation, we
        use the definition of the unique extension of a function $A \to UB$ to a
        homomorphism out of $T'(A)$, as defined in the proof of Lemma
        \ref{lem:reader-monad-transformer-free-models} above. The universal
        property for the lifted relation then follows using the actions of $\to$
        and on squares as well as the universal property of $\That$ on
        relations.

        \item It is easily checked that the operations are monotone and
        bisimilarity-preserving, and that the strength $t : A \times T'(A') \to
        T'(A \times A')$ is a morphism of predomains.
        
        \item Lastly, for any predomain $A$, we have
        %
        \begin{align*}
            |\That'(A)| &= |R \to \That(A)| \\
                        &= R \to |\That(A)| \\
                        &= R \to T(|A|)  \\
                        &= T'(|A|).
        \end{align*}
    \end{itemize}
\end{proof}

Now we are ready to define the action of the reader transformer on graduality
schemata. Let $\calA$ be an ordinary algebraic theory, and let $\calG$ be a
graduality schema for $\calA$.
%
We write $\calAs, \calAsb$ for the guarded extensions of $\calA$ associated with
$\calG$, and we write $T$ for the free-model monad for $\calAs$ on Set, and
$\That$ for its lifting to predomains.
%
We define a graduality schema for $(\calA \otimes \reader{R})$, as follows.
\begin{itemize}
    \item We define the guarded extensions 
    $\calAgtimes := \calAs \otimes \reader{R}$ and 
    $\calAgtimes_b := \calAsb \otimes \reader{R}$.

    \item Then by Lemma \ref{lem:reader-monad-transformer-free-models}, the free models of
    $\calAgtimes$ on Set are given by $T'(A) := R \to T(A)$.

    \item By Lemma \ref{lem:reader-monad-transformer-free-domains}, we have that 
    the free domains for $\calAgtimes$ are given by $R \to \That(A)$.

    \item We define the extensional types $E'(A) := R \to E(A)$, where
    $E$ is the family of extensional types for $\calG$. Next, we define the
    family of functions
    %
    $\col'_A : T'^{gl}(A) \to E'(A)$, i.e., $(R \to T^{gl}(A)) \to (R \to E(A))$, via 
    %
    \[ \col'_A(f) := \col_{A} \circ f, \]
    %
    where $\col$ is the family of functions associated with the schema $\calG$.

    \item Lastly, the relation $\ltbisim_{E'(A)}$ is defined to be $R \to
    \ltbisim_{E(A)}$, i.e., we have $f \binrel{\ltbisim_{E'(A)}} g$ iff
    for all $r \in R$, we have $(f\, r) \binrel{\ltbisim_{E(A)}} (g\, r)$.
    
    Then we observe that $\ltbisim_{E'(A)}$ is a partial order for each $A$. We
    now show that the error ordering and bisimilarity relations on the
    globalized monad imply $\ltbisim_{E'(A)}$ on the output of $\col'_A$.
    %
    More concretely, let $f, g : R \to T^{gl}(A)$, and suppose that 
    $f \ltdyn_{R \to T^{gl}(A)} g$. Then we want to show that
    $\col'_A(f) \ltbisim_{R \to E(A)} \col'_A(g)$. Let $r \in R$.
    It suffices to show
    %
    \[ \col_{A}(f(r)) \ltbisim_{E(A)} \col_{A}(g(r)). \] 
    %
    This follows from the fact that $f(r) \ltdyn_{T^{gl}(A)} g(r)$.

    The same reasoning holds for bisimilarity.
\end{itemize}














\end{document}
